% Options for packages loaded elsewhere
\PassOptionsToPackage{unicode}{hyperref}
\PassOptionsToPackage{hyphens}{url}
\PassOptionsToPackage{dvipsnames,svgnames,x11names}{xcolor}
\documentclass[
  10pt,
]{article}
\usepackage{xcolor}
\usepackage[margin=2.2cm]{geometry}
\usepackage{amsmath,amssymb}
\setcounter{secnumdepth}{-\maxdimen} % remove section numbering
\usepackage{iftex}
\ifPDFTeX
  \usepackage[T1]{fontenc}
  \usepackage[utf8]{inputenc}
  \usepackage{textcomp} % provide euro and other symbols
\else % if luatex or xetex
  \usepackage{unicode-math} % this also loads fontspec
  \defaultfontfeatures{Scale=MatchLowercase}
  \defaultfontfeatures[\rmfamily]{Ligatures=TeX,Scale=1}
\fi
\usepackage{lmodern}
\ifPDFTeX\else
  % xetex/luatex font selection
\fi
% Use upquote if available, for straight quotes in verbatim environments
\IfFileExists{upquote.sty}{\usepackage{upquote}}{}
\IfFileExists{microtype.sty}{% use microtype if available
  \usepackage[]{microtype}
  \UseMicrotypeSet[protrusion]{basicmath} % disable protrusion for tt fonts
}{}
\makeatletter
\@ifundefined{KOMAClassName}{% if non-KOMA class
  \IfFileExists{parskip.sty}{%
    \usepackage{parskip}
  }{% else
    \setlength{\parindent}{0pt}
    \setlength{\parskip}{6pt plus 2pt minus 1pt}}
}{% if KOMA class
  \KOMAoptions{parskip=half}}
\makeatother
\usepackage{longtable,booktabs,array}
\usepackage{caption}
\captionsetup[table]{skip=6pt}
\newcounter{none} % for unnumbered tables
\usepackage{calc} % for calculating minipage widths
% Correct order of tables after \paragraph or \subparagraph
\usepackage{etoolbox}
\makeatletter
\patchcmd\longtable{\par}{\if@noskipsec\mbox{}\fi\par}{}{}
\makeatother
\usepackage{graphicx}
\makeatletter
\newsavebox\pandoc@box
\newcommand*\pandocbounded[1]{% scales image to fit in text height/width
  \sbox\pandoc@box{#1}%
  \Gscale@div\@tempa{\textheight}{\dimexpr\ht\pandoc@box+\dp\pandoc@box\relax}%
  \Gscale@div\@tempb{\linewidth}{\wd\pandoc@box}%
  \ifdim\@tempb\p@<\@tempa\p@\let\@tempa\@tempb\fi% select the smaller of both
  \ifdim\@tempa\p@<\p@\scalebox{\@tempa}{\usebox\pandoc@box}%
  \else\usebox{\pandoc@box}%
  \fi%
}
% Set default figure placement to htbp
\def\fps@figure{htbp}
\makeatother
\setlength{\emergencystretch}{3em} % prevent overfull lines
\providecommand{\tightlist}{%
  \setlength{\itemsep}{0pt}\setlength{\parskip}{0pt}}
\usepackage{fontspec}
\setmainfont{STIX Two Text}
\setmonofont{Menlo}[Scale=0.85]
\usepackage{xeCJK}
\setCJKmainfont{Songti SC}
\usepackage{microtype}
\usepackage{etoolbox}
\AtBeginEnvironment{longtable}{\footnotesize}
\usepackage{xurl}
\setlength{\emergencystretch}{3em}
\sloppy
\setCJKmonofont{Songti SC}
\usepackage{unicode-math}
\setmathfont{STIX Two Math}
\usepackage{newunicodechar}
\newunicodechar{→}{\ensuremath{\rightarrow}}
\newunicodechar{∎}{\ensuremath{\blacksquare}}
\newunicodechar{≠}{\ensuremath{\neq}}
\newunicodechar{≈}{\ensuremath{\approx}}
\newunicodechar{↔}{\ensuremath{\leftrightarrow}}
\newunicodechar{≪}{\ensuremath{\ll}}
\newunicodechar{↛}{\ensuremath{\nrightarrow}}
\newunicodechar{・}{\textperiodcentered}
\newunicodechar{⇒}{\ensuremath{\Rightarrow}}
\newunicodechar{∪}{\ensuremath{\cup}}
\newunicodechar{≥}{\ensuremath{\geq}}
\newunicodechar{≤}{\ensuremath{\leq}}
\newunicodechar{⟨}{\ensuremath{\langle}}
\newunicodechar{⟩}{\ensuremath{\rangle}}
\usepackage{bookmark}
\IfFileExists{xurl.sty}{\usepackage{xurl}}{} % add URL line breaks if available
\urlstyle{same}
% fallback for those not using the hyperref driver hyperxmp:
\makeatletter
\@ifundefined{xmpquote}{\newcommand{\xmpquote}[1]{#1}}{}
\makeatother
\hypersetup{
  pdftitle={Thermodynamic Bounds and an Optimal Tempo for Goal-Holding Agents},
  pdfauthor={David Charlot, Open Interface Engineering},
  colorlinks=true,
  linkcolor={blue},
  filecolor={Maroon},
  citecolor={Blue},
  urlcolor={blue},
  pdfcreator={LaTeX via pandoc}}

\title{Thermodynamic Bounds and an Optimal Tempo for Goal-Holding
Agents}
\author{David Charlot, Open Interface Engineering}
\date{Version 1.1, 11 October 2026. HTML:
https://research.openie.dev/papers/agency/}

\begin{document}
\maketitle

{
\setcounter{tocdepth}{2}
\tableofcontents
}
\emph{Version 1.1, 11 October 2026. Data and code:
\href{https://research.openie.dev/data/\#agency}{research.openie.dev/data/}.
LaTeX PDF:
\href{https://research.openie.dev/pdfs/latex/agency.pdf}{/pdfs/latex/agency.pdf}
(\href{https://research.openie.dev/pdfs/latex/agency.tex}{source}).}

\subsection{Abstract}\label{abstract}

This paper defines an agent as a system that holds a goal under an
energy budget, derives physical limits on how fast it can act per second
and per joule, and measures intelligence as progress per joule after a
run is complete.

\textbf{In plain words}

\begin{enumerate}
\def\labelenumi{\arabic{enumi}.}
\tightlist
\item
  To count as an agent, a system must aim at a goal, act on the world,
  and stay within an energy budget.
\item
  Physics sets a minimum energy cost for every decision, so no agent can
  act faster or cheaper than that limit.
\item
  Among systems that get the job done correctly, the smarter one is the
  one that used less energy.
\end{enumerate}

\subsection{1. Thesis}\label{thesis}

The proposal has three layers, and the order is fixed.

\begin{enumerate}
\def\labelenumi{\arabic{enumi}.}
\tightlist
\item
  \textbf{Agency comes first, and its definition is required.} This
  track proposes it as the Universal Law of Agency. In form it is a
  definition plus a gate: four clauses say what counts as an agent, and
  the fourth clause refuses any act that breaks the joule budget.
  Nothing is an agent until the definition is met. Nothing downstream is
  defined without it.
\item
  \textbf{The kinematic bounds of agency follow from the definition.}
  They say how fast agency can move, per second and per joule, and what
  it must pay to move.
\item
  \textbf{Intelligence is the calculus of agency.} Intelligence is a
  hindsight characteristic. It is read off a completed run, the way
  velocity is read off a recorded path. It is therefore a derivative of
  agency: \(\iota = dX/dJ\), where \(X\) is confirmed agency and \(J\)
  is joules.
\end{enumerate}

The thesis: among systems that pass a task's correctness test, the
ranking is by energy to completion, in joules.

The thesis is read through its gate. Capability is not discarded: it is
the gate. Burnell and colleagues ask that evaluation report results by
instance rather than one aggregate {[}333{]}, and Zhou and colleagues
and Laber and colleagues find that capability and efficiency vary
independently across models and inference budgets {[}334, 335{]}. Those
findings are about systems that do not all pass the same predicate. Once
a predicate \(C(z)\) and a failure rate are fixed, every factor other
than joules has already been spent in passing the gate, and that is the
sense of the thesis' second sentence. The gated ranking is also
practice: SPECpower ranks server work per watt at stated load {[}338{]},
cost-of-pass ranks models by the cost of a correct answer {[}336{]}, and
intelligence per watt ranks local models by accuracy per unit power
{[}337{]}. This paper's contribution is the joule form of that ranking
with a refusal clause and a hindsight derivative, not the idea of
ranking by cost under correctness.

The gate and the ranking separate two things that are often merged:
capability and intelligence. Capability is passing \(C(z)\) at all.
Intelligence is how few joules the pass takes. Krakauer, Mitchell and
Krakauer draw the same line from complexity science: impressive
capabilities are not mistaken for intelligence, which is reserved for
solving problems ``quickly and parsimoniously, with a minimal
expenditure of energy'' {[}339{]}. Two systems that hand over the same
correct answer, one directly and one after a long detour, are equally
capable. Only the first is the more intelligent, and \(J^*\) is the
quantity that tells them apart. A system can be made more capable by
spending more. It is made more intelligent only by spending less for the
same pass.

The thesis has an exact form. Correctness is a gate. Joules are the
ranking. For a task \(q\) with completeness predicate \(C(z)\), the
joules to correct completion of system \(S\) are

\[J^*_S(q) = \min_{\pi} \mathbb{E}[J_\pi] \quad \text{subject to} \quad \Pr[C(z) = 1] \ge 1 - \delta\]

The minimum runs over the policies \(\pi\) that \(S\) can execute, and
\(\delta\) is the failure rate the task allows. A deadline, when the
task has one, is written into \(C(z)\): an answer after its deadline
fails the predicate. So correctness, deadlines and task size appear in
the formulas of this track, and each appears as a constraint. None
appears as a score. Among systems that pass the gate, joules alone order
them. Parameters, tokens, benchmark scores and prices carry zero weight
in that order. That is what ``collapse to zero'' means. Correctness
becomes the constraint, and energy becomes the only metric. That is the
thesis. Energy-to-solution in high-performance computing and the
accuracy-gated measurements of MLPerf Power have the same form
{[}260{]}. This track puts an eligibility definition in front of the
gate and reads the joules as a metered density. Section 2.3 states the
ranking. Proposition 8 proves that it is a total order on systems that
survives bounded meter error, and that integrated \(\iota\) is not.

The order matters. A derivative needs a path. A path needs something
that moves. The definition says what moves. The kinematics say how it
can move. Intelligence is read off the path afterward.

This track tests that order against the literature. It sorts every
source into three bins.

\begin{itemize}
\tightlist
\item
  \textbf{Proven:} theorems and derivations. The result follows from
  stated assumptions.
\item
  \textbf{Demonstrated:} experiments, measurements and simulations. The
  result was observed.
\item
  \textbf{Opinion:} position papers, frameworks, definitions and
  philosophy. The result is argued.
\end{itemize}

No source is excluded for its method. Each source is placed by what it
gives the definition, the kinematics or the calculus.

\subsubsection{Terms}\label{terms}

\begin{itemize}
\tightlist
\item
  \textbf{SI:} the International System of Units.
\item
  \textbf{Joule (J):} the SI unit of energy. One watt for one second.
\item
  \textbf{\(k_B\):} the Boltzmann constant, \(1.380649 \times 10^{-23}\)
  J/K, exact by SI definition.
\item
  \textbf{Entropy production (\(\Sigma\)):} the total entropy a process
  creates, in units of \(k_B\). It is zero only for a reversible
  process.
\item
  \textbf{Mutual information (\(I\)):} the number of bits one variable
  carries about another.
\item
  \textbf{Dynamical activity (\(A\)):} the rate of jumps between states
  in a stochastic process. It measures how busy the dynamics are.
\item
  \textbf{TUR (thermodynamic uncertainty relation):} a theorem that ties
  the precision of a current to its entropy production.
\item
  \textbf{NMR (nuclear magnetic resonance):} a method that controls and
  reads molecular spins.
\item
  \textbf{Instability rate (\(\Lambda\)):} how fast an uncontrolled
  world drifts from the goal, in nats per second. For a linear plant run
  at \(f\) steps per second,
  \(\Lambda = f\sum_{|\lambda_i|>1}\ln|\lambda_i|\).
\item
  \textbf{Reset error (\(\varepsilon\)):} the probability that a reset
  record is not in its most likely state.
\item
  \textbf{Counting observable:} a quantity of a trajectory that counts
  selected jumps. It never decreases.
\item
  \textbf{Radon-Nikodym derivative:} the density of one measure with
  respect to another. It exists when the first measure gives zero mass
  to every set that the second gives zero mass to.
\item
  \textbf{Mixture of Limits:} the rule of using the cheapest mechanism
  (gear) that is sufficient. There exist floors past which additional
  tokens, parameters, or joules do not purchase verifiable progress.
  Spend past the floor is waste. That floor is the plateau (threshold
  convergence). Mixture of Limits is not one fixed cascade. A fixed
  cascade order is one instance of it. The σ-law selector is another.
\item
  \textbf{Receipt:} a record attached to every act, holding its
  prediction, its outcome and its joules.
\end{itemize}

\subsection{2. Formal framework}\label{formal-framework}

Definitions fix the terms. Propositions are proven from cited theorems,
with their assumptions stated.

Figure 1 shows how the three layers stack.

\includegraphics[width=0.95\linewidth,height=\textheight,keepaspectratio]{fig/three-layers.png}

\emph{Figure 1. Agency in three layers. Diagram of Sections 2.1 to 2.3:
the definition and gate, the kinematic bounds, and intelligence as the
calculus of agency. A
\href{https://research.openie.dev/figures/agency/three-layers.png}{PNG
copy} is provided.}

\subsubsection{2.1 Layer 1. The definition and its gate:
eligibility}\label{layer-1.-the-definition-and-its-gate-eligibility}

For a system \(S\), an act \(a\) and an episode budget \(B\) in joules,
eligibility \(E(S,a) \in \{0,1\}\) equals 1 only when all four clauses
hold.

\begin{enumerate}
\def\labelenumi{\arabic{enumi}.}
\tightlist
\item
  \textbf{Prediction first.} A result template \(R^*(a)\) is fixed at
  time \(t_0\), before the act, and it discriminates:
  \(D(R^*(a)) \ge d_0\). Here \(D\) is the probability that the template
  rejects an outcome drawn without the act, and \(d_0 > 0\) is fixed
  with the template. A template that accepts every outcome predicts
  nothing.
\item
  \textbf{Coupling to the world.} \(I(\mathrm{do}(a); W') > 0\). Here
  \(W'\) is the world state downstream of the actuator, and
  \(\mathrm{do}(a)\) is Pearl's intervention operator {[}1{]}. The
  coupling must survive changes to the system's own inputs. Editing your
  own sensors does not count {[}2{]}.
\item
  \textbf{Consequence borne.} The comparison
  \(m = M(R, R^*(a)) \in \{1, \mu, 0\}\) (match, partial, mismatch)
  changes the system: \(s_{t+1} = f(s_t, m)\).
\item
  \textbf{Physical floor: the joule budget.} \(J(a) \le B\). If no act
  meets clauses 1 to 4, the act is refusal, with a receipt.
\end{enumerate}

Physics already prices clauses 1 to 3. The second law with feedback
{[}3, 4{]} and the measurement and erasure bound {[}5{]} give, for one
bath at temperature \(T\):

\[W_{\mathrm{ext}} \le -\Delta F + k_B T\, I, \qquad W_{\mathrm{meas}} + W_{\mathrm{erase}} \ge k_B T\, I\]

Here \(I\) is the mutual information the act's comparison acquires, in
nats. One bit is \(\ln 2\) nats, so each bit costs \(k_B T\ln 2\). An
acceptor that compares and then resets pays at least \(k_B T \ln 2\) per
bit over a closed cycle {[}6, 7{]}. The definition is physical from its
first clause.

\paragraph{Definition 1. Coupling}\label{definition-1.-coupling}

Clause 2 is measured by causal information flow. Let act \(a\) be drawn
from the policy \(p(a)\) and set by intervention, and let \(W'\) be the
world downstream of the actuator. Then

\[I(\mathrm{do}(a); W') = \sum_{a} p(a) \sum_{w} p(w \mid \mathrm{do}(a)) \log_2 \frac{p(w \mid \mathrm{do}(a))}{\sum_{a'} p(a')\, p(w \mid \mathrm{do}(a'))}\]

This is the information flow of Ay and Polani from the act node to
\(W'\) {[}8{]}. It equals the mutual information of the interventional
joint distribution \(p(a)\,p(w \mid \mathrm{do}(a))\). Its maximum over
policies is the capacity of the channel from acts to the world:

\[0 \;\le\; I(\mathrm{do}(a); W') \;\le\; \max_{p(a)} I(\mathrm{do}(a); W') \;\le\; \log_2 |\mathcal{A}|\]

Here \(\mathcal{A}\) is the set of available acts. Klyubin, Polani and
Nehaniv defined empowerment as this capacity, with the agent's own
future sensors in place of \(W'\) {[}9{]}. Clause 2 takes the world
side, so editing a sensor adds nothing {[}2{]}.

\textbf{Realized, not capacity.} The quantity credited to an act is the
realized flow under the policy the agent ran: \(I(\mathrm{do}(a); W')\)
at that \(p(a)\). The capacity is an upper bound and is never credited.
A system is not credited for coupling it could have made and did not.

\textbf{Estimation.} The flow is estimated from receipts. The agent's
randomized policy assigns acts, so the receipt's joint counts of
\((a, w)\) are interventional samples. The plug-in mutual information of
those counts is biased upward by about
\((|\mathcal{A}|-1)(|\mathcal{W}|-1)/(2n)\) nats for \(n\) samples, with
\(\mathcal{W}\) the set of distinguishable world outcomes. That
Miller-Madow term is subtracted, and a bootstrap interval is reported
{[}315{]}. At receipt-scale counts that correction is not enough. The
Bayesian estimator of Nemenman, Shafee and Bialek {[}325{]} is the
stated replacement for the next run, and the receipt schema is to record
which estimator produced each figure. Where the policy is deterministic
in a context, acts are not randomized there, and the flow is not
identified from that run. The receipt then records the capacity bound in
its own column, labeled as a bound, and \(X\) for that context is
reported as an upper bound. The run of Section 8.11 credited the
capacity bound \(b = n_{\mathrm{empty}}\log_2 9\). Its \(X\) and its
\(\iota\) are upper bounds on realized coupling and realized
intelligence. Its \(J^*\) ranking does not use \(X\), so it stands.

Why this measure:

\begin{itemize}
\tightlist
\item
  \textbf{It is interventional.} Pearl's do-calculus gives the operator
  {[}1{]}. Kenton, Kumar, Farquhar, Richens, MacDermott and Everitt
  define agents as systems that would adapt their policy if their acts
  influenced the world in a different way, and they discover agents from
  interventional data {[}10{]}. A receipt is interventional by
  construction. The agent sets its acts, and the outcome is read from
  the world.
\item
  \textbf{Observational measures fail.} James, Barnett and Crutchfield
  showed by example that transfer entropy does not measure information
  flow {[}11{]}. Lizier and Prokopenko separated information transfer
  from causal effect and assigned causal effect to causal information
  flow {[}12{]}.
\item
  \textbf{It is priced.} Under intervention the flow is a mutual
  information. Mutual information is the quantity that the feedback
  second law and the measurement and erasure bound price in joules {[}3,
  5{]}.
\item
  \textbf{It needs no causal graph.} Janzing, Balduzzi, Grosse-Wentrup
  and Schölkopf derived a measure of causal influence from postulates,
  given the causal graph and the joint distribution {[}13{]}. A receipt
  carries acts and outcomes, not the agent's full graph. The flow is
  computed from them directly.
\item
  \textbf{Credit goes to the act that caused it.} When several acts
  touch one outcome, the bits credited to act \(k\) follow the account
  of actual causation for single transitions given by Albantakis,
  Marshall, Hoel and Tononi {[}14{]}.
\end{itemize}

\paragraph{Definition 2. Agent
boundary}\label{definition-2.-agent-boundary}

Let the physical system be a finite set \(V\) of components, each behind
a metered wall. A candidate agent is a subset \(S \subseteq V\), or a
coarse-graining of one, that holds the record of every comparison
\(m_k\) (clause 3), with \(W'\) outside \(S\). For a completed run, let
\(X_S\) be the confirmed coupled bits credited with \(S\) as the agent,
and let \(J_S\) be the joules metered at the wall of \(S\).

\textbf{The agent is a candidate that maximizes \(X_S / J_S\).}

\(V\) is finite, so there are finitely many candidates. A maximizer
exists whenever some candidate has \(J_S > 0\). The receipt names the
maximizer and records any tie.

\textbf{Act-closed metering.} Clause 4 charges an act's joules to the
agent that acts. So a candidate credited with the coupling of an act is
charged every joule that realizes it: the joules of the actuator that
executes the act and of the link that carries the command. \(J_S\) in
Definition 2 is read under this rule. Takahashi and Hayashi show why the
rule is needed. With an open boundary, information gain and in-boundary
dissipation decouple, and bits per joule can be made as large as one
likes {[}239{]}.

\paragraph{Lemma 1. The maximizing boundary contains the
actuator}\label{lemma-1.-the-maximizing-boundary-contains-the-actuator}

\textbf{Setting.} A component \(\alpha \in V\), the actuator, is a cut
between the rest of \(V\) and the world: every causal path from a
component of \(V \setminus \{\alpha\}\) to \(W'\) passes through the
output \(a\) of \(\alpha\). A candidate \(S\) with \(\alpha \notin S\)
acts on \(W'\) only through a command \(c\) that it sends to \(\alpha\).
Walls are disjoint, as in Proposition 7, so
\(J_{S \cup \{\alpha\}} = J_S + J_\alpha + J_0\), with link joules
\(J_0 \ge 0\).

\textbf{Assumptions.}

\begin{itemize}
\tightlist
\item
  (B1) Cut. For a candidate without \(\alpha\), \(c \to a \to W'\) is a
  Markov chain under intervention on \(c\).
\item
  (B2) Act-closed metering. A candidate without \(\alpha\) that is
  credited coupling through \(\alpha\) is charged
  \(J_S + J_\alpha + J_0\).
\item
  (B3) A tie between candidates goes to the larger by inclusion.
\end{itemize}

\textbf{Statement.} For every candidate \(S\) with \(\alpha \notin S\)
and positive charged joules,
\(X_S / (J_S + J_\alpha + J_0) \le X_{S\cup\{\alpha\}} / J_{S\cup\{\alpha\}}\).
So some maximizer of Definition 2 contains \(\alpha\), and under (B3)
the named maximizer does. The inequality is strict whenever the actuator
adds coupling, \(X_S < X_{S\cup\{\alpha\}}\).

\textbf{Proof.}

\begin{enumerate}
\def\labelenumi{\arabic{enumi}.}
\tightlist
\item
  Fix an act. Both candidates hold the same comparison records, because
  \(S\) holds every record (Definition 2) and \(S \cup \{\alpha\}\)
  contains \(S\). So \(E_k\) and \(m_k\) are the same for both. By (B1)
  and the data-processing inequality, applied to the interventional
  joint distribution,
  \(I(\mathrm{do}(c); W') \le I(\mathrm{do}(a); W')\). Summing over
  acts, \(X_S \le X_{S\cup\{\alpha\}}\).
\item
  By (B2) and the composition of walls (Proposition 7, part 1), the
  joules charged to \(S\) equal \(J_{S\cup\{\alpha\}}\).
\item
  Step 1 over the common denominator of step 2 gives the inequality. It
  is strict when \(X_S < X_{S\cup\{\alpha\}}\).
\item
  \(V\) is finite, so a maximizer \(S^*\) exists. If
  \(\alpha \notin S^*\), step 3 makes \(S^* \cup \{\alpha\}\) a
  maximizer too, and (B3) names it. \(\blacksquare\)
\end{enumerate}

\textbf{Reading.} Without (B2), Definition 2 is degenerate. A small
subset that holds the records and leaves the actuator's draw outside its
wall reports the same coupled bits on fewer joules. Act-closed metering
removes that, at no cost in practice, because the actuator's joules are
already on the receipt under clause 4. Components off the causal path
are still removed by the ratio: dropping one leaves \(X\) unchanged and
lowers \(J\). The search does not run over all \(2^{|V|}\) subsets. It
runs over candidates that hold the records and contain the actuator. On
a host those are few: processor, accelerator, memory, actuator. One
meter suffices when its domains attribute power to those components
(Section 8.1).

Why this boundary:

\begin{itemize}
\tightlist
\item
  \textbf{The ratio, not \(X_S\) alone.} Maximizing \(X_S\) alone puts
  no price on components that add joules and no confirmed bits. The
  ratio removes them.
\item
  \textbf{Not the Markov blanket.} Aguilera, Millidge, Tschantz and
  Buckley derived that the Markov blanket condition holds only in a
  narrow range of parameters, even for weakly coupled linear stochastic
  systems {[}15{]}. Biehl, Pollock and Kanai proved by counterexample
  that the definitions of the Markov blanket in use are not equivalent
  {[}16{]}. Bruineberg, Dołęga, Dewhurst and Baltieri separate the
  blanket as a tool of inference from the blanket as a physical boundary
  {[}17{]}.
\item
  \textbf{No observer.} Orseau, McGregor McGill and Legg define agent
  and device by a Bayesian comparison of two descriptions of the same
  behavior, so the label depends on the observer's priors {[}18{]}.
  \(X_S\) and \(J_S\) are read off receipts and meters.
\item
  \textbf{A boundary found by a measure.} IIT 4.0 (integrated
  information theory) takes the system to be the set of elements with
  maximal integrated cause-effect power {[}19{]}. Krakauer,
  Bertschinger, Olbrich, Flack and Ay identify individuals as aggregates
  that propagate information from their past into their future {[}20{]}.
  Definition 2 keeps the form and uses the definition's own currency:
  confirmed coupled bits per joule.
\item
  \textbf{Coarse-grainings count.} Hoel, Albantakis and Tononi showed
  that a macro description can carry more causal power than the micro
  one {[}21{]}.
\item
  \textbf{Tools move the boundary.} When a person uses a tool, the
  candidates include the person alone, the tool alone and the pair.
  Definition 2 credits the coupling to whichever candidate maximizes
  \(X_S/J_S\), with every joule of the act charged to it. Krakauer
  models the person, the tool and the task as one dynamical system in
  which retained competence and reliance co-evolve. Above a critical
  availability of the tool, competence collapses to a dependent floor,
  and under an uncertain goal agency itself can pass to the tool
  {[}340{]}. In that model, a tool whose working the user can
  reconstruct, such as an abacus, leaves competence with the person
  after the tool is removed, while satellite navigation lies in the
  region where competence diminishes with use {[}340{]}. The ratio of
  Definition 2 makes this transfer measurable: test the person without
  the tool, and the confirmed bits that remain are the ones the person
  holds.
\item
  \textbf{Every candidate has an interface.} Rosas, Mediano, Biehl,
  Chandaria and Polani proved that every bipartite stochastic process
  has a causal blanket, and gave an algorithm that finds it from data
  with no steady-state or Markov assumption {[}22{]}.
\end{itemize}

\subsubsection{2.2 Layer 2. Kinematics}\label{layer-2.-kinematics}

Cumulative joules are \(J(t) = \int_0^t P(t')\,dt'\), with \(P\) the
metered power.

Agency displacement counts bits of confirmed, coupled change:

\[X(t) = \sum_{k:\,t_k \le t} E_k\, m_k\, I(\mathrm{do}(a_k); W'_k)\]

Per act, \(I(\mathrm{do}(a_k); W'_k)\) is the realized flow of the
policy that chose \(a_k\), estimated over the run's acts in the same
context (Definition 1).

Velocity per second is \(v_t = dX/dt\). Velocity per joule is
\(v_J = dX/dJ\). Acceleration per joule is \(\alpha_J = d^2X/dJ^2\). The
budget is \(J(T) \le B\).

The kinematic bounds are the limits on these quantities.

\textbf{K1. Landauer floor, closed cycle.} Each confirmed bit is
recorded and later erased to reset the acceptor. So

\[v_J \le \frac{1}{k_B T \ln 2}.\]

At \(T = 300\) K, \(k_B T \ln 2 \approx 2.87 \times 10^{-21}\) J,
computed from the SI constants. That is about \(3.5 \times 10^{20}\)
bits per joule. Proposition 1 proves K1 act by act, for the mean joules
of each act, under stated assumptions on the record.

\textbf{K2. Classical speed limit.} Shiraishi, Funo and Saito proved
this bound for a Markov jump process with local detailed balance
{[}23{]}:

\[\tau \;\ge\; \frac{L\big(p(0),p(\tau)\big)^2}{2\,\Sigma_{\mathrm{tot}}\,\langle A\rangle_\tau}\]

\begin{itemize}
\tightlist
\item
  \(L(p,p') = \sum_i |p_i - p'_i|\) is the full \(L_1\) distance between
  two probability distributions. There is no factor of 1/2. With total
  variation distance \(d_{TV} = L/2\), the same bound reads
  \(\tau \ge 2 d_{TV}^2/(\Sigma_{\mathrm{tot}}\langle A\rangle_\tau)\).
\item
  \(\Sigma_{\mathrm{tot}} = \int_0^\tau \dot\Sigma\,dt\) is total
  entropy production, system plus baths, in units of \(k_B\).
\item
  \(A(t) = \sum_{i\ne j} W_{ij}(t)\,p_j(t)\) is the dynamical activity,
  with \(W_{ij}\) the jump rate from state \(j\) to state \(i\).
  \(\langle A\rangle_\tau = \frac{1}{\tau}\int_0^\tau A(t)\,dt\) is its
  time average.
\end{itemize}

Rearranged,
\(\Sigma_{\mathrm{tot}} \ge L^2/(2\tau\langle A\rangle_\tau)\). Moving a
distribution far, fast, with few jumps, costs entropy.

\textbf{Joule form.} For one bath,
\(k_B T\,\Sigma_{\mathrm{tot}} = W - \Delta F_{\mathrm{neq}}\), with
\(F_{\mathrm{neq}}\) the nonequilibrium free energy {[}24{]}. So the
dissipated work obeys

\[W_{\mathrm{diss}} \;\ge\; \frac{k_B T\, L^2}{2\,\tau\,\langle A\rangle_\tau}.\]

\textbf{Stationary currents.} The same paper gives
\(\tau \ge c^* L^2 / (2\,\Sigma_{HS}\,\langle A\rangle_\tau)\), with
\(c^* = 0.896\ldots\) and \(\Sigma_{HS}\) the Hatano-Sasa, or excess,
entropy production {[}23{]}.

\textbf{K3. Precision costs dissipation.} Acts are driven and not
stationary, so the form that applies is Koyuk and Seifert's uncertainty
relation for time-dependent driving {[}323{]}; the steady-state form is
stated here for reference. For a time-antisymmetric current \(Y\) in a
steady state,
\(\mathrm{Var}(Y)/\langle Y\rangle^2 \ge 2/\Sigma_{\mathrm{tot}}\)
{[}25, 26{]}. Agents use feedback, and feedback beats this bound in
experiment {[}27{]}. The bound for agents is the feedback form, in which
mutual information enters the bound {[}28{]}. Proposition 3 states it
for confirmed agency.

\textbf{K4. Finite-time cost.} Erasing a bit in finite time \(\tau\)
costs more than \(k_B T \ln 2\). The excess grows as \(\tau\) shrinks
{[}29{]}. Near equilibrium, the excess work along a protocol is bounded
below by the squared thermodynamic length over the duration {[}30{]}.
Proposition 1 gives the floor for a whole act.

\textbf{K5. Quantum ceiling on rate.} A system with mean energy
\(\langle E\rangle\) above its ground state makes at most
\(2\langle E\rangle/(\pi\hbar)\) orthogonal transitions per second
{[}31{]}. That is about \(6.04 \times 10^{33}\) transitions per second
per joule, computed from the constants. The figure of about
\(3 \times 10^{33}\) sometimes quoted is \(2E/h\), half the
Margolus-Levitin rate, since \(2E/(\pi\hbar) = 4E/h\).

\paragraph{Proposition 1. Joule floor of one act in finite
time}\label{proposition-1.-joule-floor-of-one-act-in-finite-time}

\textbf{Setting.} Act \(k\) writes its comparison to a record \(M_k\)
with \(N\) states, then resets the record. The record is a Markov jump
process with local detailed balance, in contact with one bath at
temperature \(T\). During the write, the compared outcome \(R_k\) is
held fixed, and the record's jump rates may depend on it. During the
reset, the record's rates depend on nothing else, so the record evolves
alone. The reset lasts \(\tau_k\). \(\langle A\rangle_k\) is the
record's time-averaged dynamical activity during the reset. Joules and
works are means over the ensemble of act \(k\).

\textbf{Assumptions.}

\begin{itemize}
\tightlist
\item
  (R1) The record's states have equal energy. The record shares no
  interaction energy with \(R_k\) at the start or the end of the write.
  At the start of the act the record is in its ready distribution
  \(q^0\), independent of \(R_k\).
\item
  (R2) The credited bits do not exceed the information the write leaves
  in the record: \(b_k \le I(M_k; R_k)\), in bits, at the end of the
  write. The record is unchanged between the write and the reset.
\item
  (R3) The reset returns the record to \(q^0\). The reset error is
  \(\varepsilon = 1 - \max_s q^0_s\).
\item
  (R4) The work of the write and of the reset passes through the meter:
  \(J_k \ge W_k^{\mathrm{write}} + W_k^{\mathrm{reset}}\).
\end{itemize}

\textbf{Statement.}

\[J_k \;\ge\; b_k\, k_B T \ln 2 \;+\; \frac{2\, k_B T\, \big[\big(\phi_N^{-1}(b_k) - \varepsilon\big)_+\big]^2}{\tau_k\, \langle A\rangle_k}\]

Here \((x)_+ = \max(x, 0)\), \(h_2(x) = -x\log_2 x - (1-x)\log_2(1-x)\)
is the binary entropy, and \(\phi_N(x) = h_2(x) + x\log_2(N-1)\). On
\([0, 1 - 1/N]\), \(\phi_N\) increases from 0 to \(\log_2 N\), so its
inverse is defined on \([0, \log_2 N]\). For a one-bit record that
confirms one bit, \(\phi_2^{-1}(1) = 1/2\) and

\[J_k \;\ge\; k_B T \ln 2 \;+\; \frac{2\, k_B T\, \big[(1/2 - \varepsilon)_+\big]^2}{\tau_k\, \langle A\rangle_k}\]

As \(\varepsilon \to 0\) this tends to
\(k_B T\ln 2 + k_B T/(2\tau_k\langle A\rangle_k)\).

\textbf{Proof.}

\begin{enumerate}
\def\labelenumi{\arabic{enumi}.}
\tightlist
\item
  For an isothermal Markov jump process with local detailed balance,
  \(W = \Delta F_{\mathrm{neq}} + k_B T\, \Sigma_{\mathrm{tot}}\) with
  \(\Sigma_{\mathrm{tot}} \ge 0\). Here
  \(F_{\mathrm{neq}} = \langle E\rangle - k_B T\, S\), with \(S\) the
  Shannon entropy in nats {[}24{]}.
\item
  The write. Apply step 1 to the pair \((M_k, R_k)\) with \(R_k\) fixed.
  By (R1) the mean energy does not change and the pair starts
  independent. The entropy of \(R_k\) does not change. So the joint
  entropy changes by \(H(q') - I - H(q^0)\), where \(q'\) is the
  record's distribution after the write and \(I = I(M_k; R_k)\), all in
  nats. So \(W_k^{\mathrm{write}} \ge k_B T\,[\,I - H(q') + H(q^0)\,]\).
\item
  The reset. The record evolves alone, so step 1 applies to the record.
  Its energy does not change. By (R2) and (R3) its entropy goes from
  \(H(q')\) to \(H(q^0)\). So
  \(W_k^{\mathrm{reset}} = k_B T\,[\,H(q') - H(q^0)\,] + k_B T\, \Sigma^{\mathrm{reset}}\).
\item
  Steps 2 and 3 and (R4) give
  \(J_k \ge k_B T\, I + k_B T\, \Sigma^{\mathrm{reset}}\). The record's
  entropies cancel over the cycle. In bits,
  \(k_B T\, I = k_B T\ln 2\; I(M_k;R_k) \ge b_k\, k_B T\ln 2\) by (R2).
  This is the measurement and erasure bound of Sagawa and Ueda {[}5{]},
  here for any ready distribution.
\item
  The classical speed limit gives
  \(\Sigma^{\mathrm{reset}} \ge L^2/(2\tau_k\langle A\rangle_k)\), with
  \(L = \sum_s \lvert q'_s - q^0_s\rvert\) {[}23{]}. Take \(s\) with
  \(q^0_s = \max q^0 = 1 - \varepsilon\). The other states carry the
  opposite change in total, so
  \(L \ge 2(q^0_s - q'_s) \ge 2\,[(1 - q'_{\max}) - \varepsilon]\).
\item
  Fano's inequality, with no side information and the most likely state
  as the guess, gives
  \(H(q') \le h_2(1-q'_{\max}) + (1-q'_{\max})\log_2(N-1) = \phi_N(1 - q'_{\max})\),
  in bits. Since \(1 - q'_{\max} \le 1 - 1/N\) and \(\phi_N\) increases
  there, \(1 - q'_{\max} \ge \phi_N^{-1}(H(q'))\). By (R2),
  \(H(q') \ge I(M_k;R_k) \ge b_k\), and \(\phi_N^{-1}\) increases. So
  \(L \ge 2\big(\phi_N^{-1}(b_k) - \varepsilon\big)_+\).
\item
  Steps 4 to 6 give the statement. \(\blacksquare\)
\end{enumerate}

\textbf{Reading.} The floor is a sum of two prices. The information
price is \(k_B T\ln 2\) per confirmed bit, paid over the closed cycle
whatever the ready state. The speed price belongs to the reset. A reset
with finite rates never reaches a single state in finite time, so
\(\varepsilon > 0\) at every finite \(\tau_k\), and the speed price
carries that error. K1 holds act by act at every \(\tau_k\). Tighter
erasure bounds for particular dynamics are also lower bounds on the same
reset (Section 4.1). Per act, the velocity per joule is

\[v_{J,k} = \frac{b_k}{J_k} \;\le\; \frac{b_k}{b_k\, k_B T \ln 2 + 2 k_B T\,\big[\big(\phi_N^{-1}(b_k) - \varepsilon\big)_+\big]^2/(\tau_k\langle A\rangle_k)} \;\le\; \frac{1}{k_B T \ln 2}\]

The last inequality is strict whenever
\(\phi_N^{-1}(b_k) > \varepsilon\) at finite
\(\tau_k\langle A\rangle_k\).

\paragraph{Proposition 2. Standing joules to hold a
goal}\label{proposition-2.-standing-joules-to-hold-a-goal}

\textbf{Setting.} The world is a linear plant
\(x_{t+1} = A x_t + B u_t + w_t\), with \(x_t \in \mathbb{R}^n\). At
step \(t\) the agent writes a finite-valued record \(c_t\) of the plant.
Its memory \(s_t\) is a finite-valued function of \((s_{t-1}, c_t)\),
with \(s_0\) fixed, and its act \(u_t\) is a function of \(s_t\).
Keeping every record, \(s_t = (s_{t-1}, c_t)\), is one such memory. The
goal is to hold the plant: \(\sup_t \mathbb{E}\|x_t\|^2 < \infty\).

\textbf{Assumptions.}

\begin{itemize}
\tightlist
\item
  (C1) The noise \(w_t\) is independent of \((x_t, s_t)\).
\item
  (C2) \(x_0\) has a density with finite differential entropy and a
  finite second moment.
\item
  (C3) Closed-cycle accounting per step. \(J_t\) is the mean joules
  attributed to the record \(c_t\): its write and its later erasure,
  whenever they are paid. \(J_t \ge k_B T\, I(x_t; c_t \mid s_{t-1})\),
  with the information in nats. This is Proposition 1, steps 1 to 4, at
  each value of the memory \(s_{t-1}\), with a fresh record in its ready
  state and \(x_t\) held fixed during the write {[}5{]}. A memory of
  bounded size erases each record within a bounded number of steps, so
  the attributed joules fall inside the run up to a boundary term.
\end{itemize}

\textbf{Statement.} This is a joule-priced corollary of the data-rate
theorem of Tatikonda and Mitter {[}169{]} and of Nair and Evans
{[}170{]}. The closed-loop inequality of Touchette and Lloyd {[}32,
33{]} carries the rate to the record, and Landauer's erasure price
{[}6{]} converts each bit to \(k_B T\ln 2\) joules. It is new only as
that composition. Let \(\lambda_1, \dots, \lambda_n\) be the eigenvalues
of \(A\). Then

\[\liminf_{N\to\infty} \frac{1}{N}\sum_{t=1}^{N} J_t \;\ge\; k_B T \ln 2 \sum_{i:\,|\lambda_i| > 1} \log_2 |\lambda_i|\]

\textbf{Proof.}

\begin{enumerate}
\def\labelenumi{\arabic{enumi}.}
\tightlist
\item
  Choose real coordinates in which \(A = \mathrm{diag}(A_u, A_s)\), with
  \(A_u\) carrying exactly the eigenvalues of modulus above 1. Let
  \(y_t\) be the unstable block of \(x_t\). Then
  \(y_{t+1} = A_u y_t + v_t + e_t\), where \(v_t\) is a function of
  \(s_t\) and \(e_t\) is independent of \((y_t, s_t)\) by (C1). A linear
  change of coordinates shifts differential entropies by a constant, so
  finiteness is unchanged.
\item
  Let \(D_t = h(y_t \mid s_t)\), a conditional differential entropy in
  nats. Adding independent noise does not lower a conditional entropy.
  Shifting by a function of \(s_t\) does not change it. The linear map
  \(A_u\) adds \(\ln|\det A_u|\). So
  \(h(y_{t+1} \mid s_t) \ge D_t + \ln|\det A_u|\).
\item
  \(s_{t+1}\) is a function of \((s_t, c_{t+1})\), and conditioning on
  less never lowers entropy. So
  \(D_{t+1} \ge h(y_{t+1}\mid s_t, c_{t+1}) = h(y_{t+1}\mid s_t) - I(y_{t+1}; c_{t+1} \mid s_t)\).
  Since \(y\) is a function of \(x\),
  \(I(y_{t+1}; c_{t+1}\mid s_t) \le I(x_{t+1}; c_{t+1}\mid s_t)\). This
  is the closed-loop inequality of Touchette and Lloyd {[}32, 33{]}.
\item
  Steps 2 and 3 give
  \(I(x_{t+1}; c_{t+1}\mid s_t) \ge \ln|\det A_u| - (D_{t+1} - D_t)\).
  Summing over \(N\) steps,
  \(\sum_{t=1}^{N} I(x_t; c_t\mid s_{t-1}) \ge N \ln|\det A_u| - (D_N - D_0)\).
\item
  Each record is finite-valued, so each information term is finite.
  \(D_0 = h(y_0)\) is finite: the second moment bounds it above, and
  \(h(x_0) \le h(y_0) + h(x^{\mathrm{s}}_0)\), with \(x^{\mathrm{s}}_0\)
  the stable block, bounds it below, by (C2). \(D_N \le h(y_N)\), and a
  Gaussian has the largest entropy for a given covariance, so the
  bounded second moment bounds \(D_N\) above. Dividing by \(N\), the
  average information per step is at least
  \(\ln|\det A_u| = \sum_{|\lambda_i|>1}\ln|\lambda_i|\) in the limit.
\item
  (C3) prices each nat at \(k_B T\), which is \(k_B T\ln 2\) per bit.
  \(\blacksquare\)
\end{enumerate}

\textbf{Reading.} An agent that holds a goal against an unstable world
pays a standing power of at least
\(f\, k_B T\ln 2\sum_{|\lambda_i|>1}\log_2|\lambda_i|\) at \(f\) control
steps per second. One unstable mode with \(|\lambda| = 2\) needs one bit
per step, which is \(k_B T\ln 2 \approx 2.87\times10^{-21}\) J per step
at 300 K. The record's bits per step bound its information from above,
so the proof also gives a non-strict, averaged form of the data-rate
theorem (Section 4.1). The composition is short. It is the data-rate
theorem of Tatikonda and Mitter and of Nair and Evans {[}169, 170{]},
priced at \(k_B T\ln 2\) per bit through Touchette and Lloyd {[}32,
33{]}. What it adds is the price in joules, which is the floor
Prediction 1 needs.

\paragraph{Proposition 3. Precision of confirmed
agency}\label{proposition-3.-precision-of-confirmed-agency}

\textbf{Setting.} The agent, its world and its feedback form one
time-homogeneous, irreducible Markov jump process on finitely many
states, in its stationary state. Each confirmed act is one jump of a
marked kind. \(N(t)\) counts confirmed acts in \([0,t]\).
\(\langle K\rangle\) is the mean number of jumps of all kinds in
\([0,t]\), and \(\langle k\rangle = \langle K\rangle/t\) is the mean
jump rate.

\textbf{Statement.}

\begin{enumerate}
\def\labelenumi{\arabic{enumi}.}
\tightlist
\item
  In the long-time limit,
  \(\mathrm{Var}(N)/\langle N\rangle^2 \ge 1/\langle K\rangle\). If
  every confirmed act credits the same \(b\) bits, then \(X = bN\) and
  the same bound holds for \(X\).
\item
  Let \(T_n\) be the time at which the \(n\)-th confirmed act closes.
  For large \(n\),
  \(\mathrm{Var}(T_n)/\langle T_n\rangle^2 \ge 1/(\langle T_n\rangle\langle k\rangle)\).
\item
  Let \(X\) instead be odd under the map to the backward experiment of a
  fluctuation relation that includes an information term,
  \(P_B(-X,-\sigma_I)/P(X,\sigma_I) = e^{-\sigma_I}\). Then
\end{enumerate}

\[\frac{\mathrm{Var}(X) + \mathrm{Var}_B(X)}{\big(\langle X\rangle + \langle X\rangle_B\big)^2} \;\ge\; \frac{1}{\exp\!\big[(\langle\sigma_I\rangle + \langle\sigma_I\rangle_B)/2\big] - 1}\]

When the backward experiment equals the forward one, this reads
\(\mathrm{Var}(X)/\langle X\rangle^2 \ge 2/\big(e^{\langle\sigma_I\rangle} - 1\big)\).

\textbf{Proof.} \(N\) is a counting observable: it counts a subset of
jumps and never decreases. Garrahan proved parts 1 and 2 for counting
observables in stationary Markov jump processes {[}34{]}. Part 1 is his
bound on the variance of a counting observable by the mean activity.
Part 2 is his first-passage bound, stated for counts whose jump weights
are 0 or 1, which holds here because \(N\) counts one marked kind of
jump. Scaling by \(b\) leaves a relative variance unchanged. Part 3 is
the uncertainty relation that Potts and Samuelsson derived from any
fluctuation relation, measurement and feedback included {[}28{]},
applied with \(X\) as the observable. \(\blacksquare\)

\textbf{Reading.} Part 1 needs no time-reversal property and no model of
the feedback. The precision of confirmed agency is bought with activity.
Part 3 is the dissipation price, and it holds under feedback because it
rests on the feedback fluctuation relation. The standard TUR does not
bind feedback engines in experiment {[}27{]}. That is why the agent form
carries the information term.

\paragraph{Proposition 4. Quantum agents keep the
floor}\label{proposition-4.-quantum-agents-keep-the-floor}

\textbf{Setting.} The comparison \(m_k\) comes from a quantum
measurement and is written to a classical register, the record of clause
3. The register's reset meets the setting of Proposition 1 and (R2) and
(R3), with \(I(M_k;R_k)\) the quantum-classical mutual information of
the measurement. The reset's work passes through the meter:
\(J_k \ge W_k^{\mathrm{reset}}\).

\textbf{Statement.}

\[J_k \;\ge\; k_B T\ln 2\,\big[\,b_k - \phi_N(\varepsilon)\,\big] \;+\; \frac{2\, k_B T\, \big[\big(\phi_N^{-1}(b_k) - \varepsilon\big)_+\big]^2}{\tau_k\, \langle A\rangle_k}\]

As \(\varepsilon \to 0\) this is the floor of Proposition 1. When the
write is metered as well, the measurement and erasure bound of Sagawa
and Ueda gives \(J_k \ge b_k\, k_B T\ln 2\) at every \(\varepsilon\),
under the additional assumption they state for quantum measurement
{[}5{]}.

\textbf{Proof.} Step 3 of Proposition 1 uses only the register and its
reset dynamics. It gives
\(W_k^{\mathrm{reset}} = k_B T\ln 2\,[H(q') - H(q^0)] + k_B T\,\Sigma^{\mathrm{reset}}\),
with entropies in bits. The quantum-classical mutual information is at
most the entropy of the outcomes {[}5{]}, so \(H(q') \ge b_k\). Fano's
inequality gives \(H(q^0) \le \phi_N(\varepsilon)\), since
\(\varepsilon \le 1 - 1/N\). Steps 5 and 6 of Proposition 1 bound
\(\Sigma^{\mathrm{reset}}\). \(\blacksquare\)

\textbf{Reading.} Quantum resources change what an agent can do per act.
They do not lower the price of its record. The record costs
\(k_B T\ln 2\) per bit, less a residue \(\phi_N(\varepsilon)\) that
vanishes as the reset becomes exact. The measurement that writes the
record carries its own price, given in Section 4.1.

\subsubsection{2.3 Layer 3. Calculus: intelligence as a derivative, read
in
hindsight}\label{layer-3.-calculus-intelligence-as-a-derivative-read-in-hindsight}

A run produces a receipt: the ordered record
\(\{(J_k, X_k)\}_{k=0}^{n}\) of metered joules and confirmed
displacement. Each \(m_k\) is known only after act \(k\) closes. So the
path \(X(J)\) exists only once the run is complete.

\begin{itemize}
\tightlist
\item
  \textbf{Intelligence, marginal:} \(\iota(J) = \dfrac{dX}{dJ}\),
  estimated from the receipt as \(\iota_k = \Delta X_k/\Delta J_k\).
\item
  \textbf{Run integral:} \(X(B) = \int_0^B \iota(J)\,dJ\).
\item
  \textbf{Joules to close a requirement:}
  \(J^*(X_{\mathrm{req}}) = \int_0^{X_{\mathrm{req}}} \frac{dX}{\iota}\).
\item
  \textbf{Optimum:} \(\pi^* = \arg\min_\pi J_\pi\) subject to
  \(X_\pi \ge X_{\mathrm{req}}\), \(E_k = 1\) for every \(k\), and
  \(J_\pi \le B\). If no policy is feasible, refuse with a receipt.
\item
  \textbf{Ranking:}
  \(S_1 \succ_q S_2 \iff J^*_{S_1}(q) < J^*_{S_2}(q)\).
\end{itemize}

\(X_{\mathrm{req}}\) is set by the task's completeness predicate
\(C(z)\). It is a gate, not a score. Parameters, tokens and benchmarks
enter only through \(J^*\). All other factors collapse to zero.

Hindsight is literal here. Velocity is read off a recorded path.
Intelligence is read off a recorded run. The receipt is the path.

\paragraph{Proposition 5. Intelligence exists as a bounded density on
the joule
axis}\label{proposition-5.-intelligence-exists-as-a-bounded-density-on-the-joule-axis}

\textbf{Setting.} A run on \([0, T]\) has metered power \(P(t) \ge 0\),
integrable, and joule measure \(dJ = P\,dt\). Act \(k\) occupies
\([s_k, e_k]\) and credits
\(b_k = E_k\, m_k\, I(\mathrm{do}(a_k); W'_k) \ge 0\) bits. Acts may
overlap. Each act has an attributed power \(P_k \ge 0\), zero outside
\([s_k, e_k]\), with \(\sum_k P_k \le P\). Its joules are
\(J_k = \int P_k\,dt > 0\). For acts that do not overlap, \(P_k = P\) on
the act's interval.

\textbf{Crediting rule.} Each act's bits are spread over its own joules:

\[dX = \sum_k \frac{b_k}{J_k}\, P_k\, dt\]

\textbf{Assumption (K1 per act).} \(b_k\, k_B T \ln 2 \le J_k\) for
every \(k\). Proposition 1 gives it under (R1) to (R4) for the mean
joules of each act, so the statement holds for the mean run. On a single
run it holds whenever each act draws more than \(b_k\, k_B T\ln 2\).
Every confirmed act in the run of Section 8.11 drew more than
\(3 \times 10^{14}\) times that.

\textbf{Statement.}

\begin{enumerate}
\def\labelenumi{\arabic{enumi}.}
\tightlist
\item
  For every measurable set \(A \subseteq [0,T]\) of times,
  \(0 \le X(A) \le J(A)/(k_B T\ln 2)\). So \(dX\) is absolutely
  continuous with respect to \(dJ\).
\item
  The Radon-Nikodym derivative \(\iota = dX/dJ\) exists, is unique up to
  sets of zero joules, and satisfies
\end{enumerate}

\[0 \;\le\; \iota \;\le\; \frac{1}{k_B T \ln 2}\]

Where \(P > 0\), \(\iota = \sum_k (b_k/J_k)(P_k/P)\). For acts that do
not overlap, \(\iota = b_k/J_k\) on act \(k\), which is the receipt's
\(\iota_k\).

\begin{enumerate}
\def\labelenumi{\arabic{enumi}.}
\setcounter{enumi}{2}
\tightlist
\item
  Confirmed agency is a function of cumulative joules:
  \(X(t) = g(J(t))\), with \(g\) nondecreasing and Lipschitz with
  constant \(1/(k_B T\ln 2)\). So \(g\) is differentiable at almost
  every joule, and \(g'(J(t)) = \iota(t)\) for almost every joule.
\item
  If each act's bits are instead credited at the instant it closes, then
  \(dX = \sum_k b_k\,\delta_{e_k}\) is singular with respect to \(dJ\)
  whenever some \(b_k > 0\). Then \(g\) is a step function, \(g' = 0\)
  at almost every joule, and all agency sits on a set of zero joules.
\end{enumerate}

\textbf{Proof.}

\begin{enumerate}
\def\labelenumi{\arabic{enumi}.}
\tightlist
\item
  \(X(A) = \sum_k (b_k/J_k)\int_A P_k\,dt \le \frac{1}{k_BT\ln 2}\int_A \sum_k P_k\,dt \le \frac{J(A)}{k_BT\ln 2}\),
  by K1 per act and then \(\sum_k P_k \le P\). If \(J(A) = 0\), then
  \(X(A) = 0\).
\item
  Both measures are finite, so the Radon-Nikodym theorem gives
  \(\iota\). If \(\iota > 1/(k_BT\ln 2)\) on a set \(B\) with
  \(J(B) > 0\), then \(X(B) > J(B)/(k_BT\ln 2)\), against step 1. On
  \(\{P > 0\}\), \(dX = \sum_k (b_k/J_k)(P_k/P)\,dJ\), and \(P_k = 0\)
  where \(P = 0\). That gives the explicit form.
\item
  If \(J(t_1) = J(t_2)\) with \(t_1 < t_2\), step 1 gives
  \(X((t_1,t_2]) = 0\), so \(g\) is well defined. For \(t_1 < t_2\),
  \(g(J(t_2)) - g(J(t_1)) = X((t_1,t_2]) \le (J(t_2) - J(t_1))/(k_BT\ln 2)\).
  \(J\) is continuous, so \(g\) is Lipschitz on \([0, J(T)]\). A
  Lipschitz function is absolutely continuous and differentiable almost
  everywhere. A change of variables gives
  \(\int_0^t g'(J(s))\,dJ(s) = g(J(t)) = X([0,t])\). So \(g'\circ J\) is
  also a density of \(X\) with respect to \(J\), and it equals \(\iota\)
  by uniqueness.
\item
  \(X(\{e_k\}) = b_k > 0\) while \(J(\{e_k\}) = 0\), because \(J\) is
  continuous. So \(dX\) is not absolutely continuous. It sits on
  finitely many points of zero joules, so it is singular. A function
  constant between finitely many jumps has zero derivative almost
  everywhere. \(\blacksquare\)
\end{enumerate}

\textbf{Reading.} Hindsight intelligence is a well-defined density on
the joule axis once each act's bits are spread over the joules that act
used. The crediting rule is what makes \(dX/dJ\) exist. Point crediting
puts all agency on a set of zero joules. Overlapping acts are handled by
any attribution that never assigns more joules than the meter read. Each
\(b_k\) is known only after act \(k\) closes, so \(\iota\) is read off
the completed run.

Proposition 5 holds by construction. The crediting rule makes \(X\)
absolutely continuous with respect to \(J\), so the density exists by
definition. Its content is the choice of rule: part 4 shows that the
other natural rule, point crediting, leaves no density at all.

\paragraph{Proposition 6. Sampled intelligence
converges}\label{proposition-6.-sampled-intelligence-converges}

\textbf{Setting.} The meter is read at cumulative joules
\(0 = j_0 < j_1 < \dots < j_M = J(T)\), with mesh
\(\Delta = \max_i (j_{i+1} - j_i)\). On \([j_i, j_{i+1})\) the sampled
intelligence is

\[\hat\iota = \frac{g(j_{i+1}) - g(j_i)}{j_{i+1} - j_i}\]

Time samples refine joule samples, because \(J\) is continuous.

\textbf{Statement.} Under the setting of Proposition 5,
\(\hat\iota(j) \to g'(j)\) at almost every joule \(j\) as
\(\Delta \to 0\), and \(0 \le \hat\iota \le 1/(k_BT\ln 2)\) at every
mesh. If the meter's reading of each sampled interval has relative error
at most \(\epsilon < 1\), the reported ratio lies in
\([\hat\iota/(1+\epsilon),\, \hat\iota/(1-\epsilon)]\).

\textbf{Proof.} \(\hat\iota(j)\) is the average of \(g'\) over a sampled
interval \(I \ni j\) of length \(h \le \Delta\). Since
\(I \subseteq [j - h, j + h]\),

\[\big|\hat\iota(j) - g'(j)\big| \;\le\; \frac{1}{h}\int_{j-h}^{j+h}\big|g'(u) - g'(j)\big|\,du\]

The right side tends to 0 at every Lebesgue point of \(g'\), and almost
every \(j\) is a Lebesgue point by Lebesgue's differentiation theorem.
The bounds on \(\hat\iota\) are the monotonicity and the Lipschitz
constant of \(g\). A reading
\(\Delta\hat J \in [(1-\epsilon)\Delta J, (1+\epsilon)\Delta J]\)
divides the same \(\Delta X\). \(\blacksquare\)

\paragraph{Proposition 7. Composition of joules and coupled
bits}\label{proposition-7.-composition-of-joules-and-coupled-bits}

\textbf{Setting.} A group \(G\) of members \(i = 1, \dots, n\) acts on
one world \(W'\). Each member sits behind its own metered wall. The
walls are disjoint and inside the group's wall, and every member draws
its energy from outside the group. Members choose acts
\(a_1, \dots, a_n\) by independent policies. Member \(i\)'s coupling is
computed with the other members acting under their own policies.

\textbf{Statement.}

\begin{enumerate}
\def\labelenumi{\arabic{enumi}.}
\tightlist
\item
  \(J_G = \sum_i J_i + J_0\), where \(J_0 \ge 0\) is the energy that
  enters the group's wall and no member's wall.
\item
  \(I(\mathrm{do}(a_1,\dots,a_n); W') \ge \sum_i I(\mathrm{do}(a_i); W')\),
  with equality if and only if \(a_i\) is independent of
  \((a_1, \dots, a_{i-1})\) given \(W'\) for every \(i\).
\item
  The inequality can be strict when every member alone is uncoupled.
  With \(a_1, a_2\) independent uniform bits and
  \(W' = a_1 \oplus a_2\), each member's coupling is 0 bits and the
  group's is 1 bit.
\item
  Under K1 per act at the group's wall, Proposition 5 applies to \(G\):
  \(\iota_G\) exists and \(0 \le \iota_G \le 1/(k_BT\ln 2)\). By
  Definition 2, \(G\) is the agent when \(X_G/J_G\) exceeds \(X_S/J_S\)
  for every other candidate, each member included.
\end{enumerate}

\textbf{Proof.}

\begin{enumerate}
\def\labelenumi{\arabic{enumi}.}
\tightlist
\item
  Every joule that enters a member's wall first enters the group's wall,
  and metered energy adds over disjoint walls.
\item
  Under intervention the acts have the product distribution
  \(\prod_i p(a_i)\), and each coupling is a mutual information of the
  interventional joint distribution (Definition 1). Member \(i\)'s
  coupling is \(I(a_i; W')\) in that joint, by marginalizing the other
  acts. By the chain rule,
  \(I(a_{1:n}; W') = \sum_i I(a_i; W' \mid a_{<i})\). Independence gives
  \(I(a_i; a_{<i}) = 0\), so
  \(I(a_i; W' \mid a_{<i}) = I(a_i; W', a_{<i}) = I(a_i; W') + I(a_i; a_{<i} \mid W') \ge I(a_i; W')\).
  Equality holds for every term if and only if
  \(I(a_i; a_{<i}\mid W') = 0\) for every \(i\).
\item
  Given \(a_1\) alone, \(W'\) is still a uniform bit, so
  \(I(a_1; W') = 0\), and likewise for \(a_2\). Given both acts, \(W'\)
  is fixed, so \(I(a_1, a_2; W') = H(W') = 1\) bit.
\item
  Proposition 5 needs only the run, the meter and K1 per act. Apply it
  at the group's wall. \(\blacksquare\)
\end{enumerate}

\textbf{Reading.} Joules add, with the group's links on top. Coupled
bits from independent members add at least, and some exist only at the
group level. A group raises \(\iota_G\) when it gains more coupled bits
than its members hold alone, or spends fewer joules acting together than
its members spend acting apart. Section 4.1 gives the theorems that
price both terms.

\paragraph{Prediction 1. The optimal act time of an agent that holds a
goal}\label{prediction-1.-the-optimal-act-time-of-an-agent-that-holds-a-goal}

\textbf{Setting.} A closed-cycle acceptor runs one act after another.
Each act writes a record with \(N\) states and credits \(b\) bits when
confirmed. The comparison, made before the reset, confirms with
probability \(p\). Each reset meets Proposition 1, with duration
\(\tau\), mean activity \(\langle A\rangle\) and reset error
\(\varepsilon\). Throughout, the agent holds a goal against a world with
instability rate \(\Lambda\), through a separate register. Every joule
is metered at the agent's wall (Definition 2).

\textbf{Derivation.} Each act lasts at least \(\tau\). Proposition 2
prices the holding at \(k_B T\,\Lambda\) per second, in the long run.
Proposition 1 prices the act. The two registers are distinct, so their
metered joules add. Per act, in the mean,

\[J(\tau) \;\ge\; b\,k_B T\ln 2 \;+\; \frac{\kappa}{\tau} \;+\; k_B T\,\Lambda\,\tau, \qquad \kappa = \frac{2k_B T\,\big[\big(\phi_N^{-1}(b) - \varepsilon\big)_+\big]^2}{\langle A\rangle}\]

By Proposition 5, the run's intelligence is confirmed bits per joule,
\(\iota(\tau) = p\,b/J(\tau)\). The right side of the bound on \(J\) is
least at \(\tau^* = \sqrt{\kappa/(k_B T\Lambda)}\), where it equals
\(b\,k_B T\ln 2 + 2\sqrt{\kappa\,k_B T\Lambda}\). For a one-bit record
and an exact reset, this gives two numbers:

\[\tau^* = \frac{1}{\sqrt{2\,\langle A\rangle\,\Lambda}}, \qquad \iota(\tau) \;\le\; \frac{p}{k_B T\,\big[\ln 2 + \sqrt{2\Lambda/\langle A\rangle}\big]}\]

\textbf{Prediction 1.}

\begin{enumerate}
\def\labelenumi{\arabic{enumi}.}
\tightlist
\item
  \textbf{Ceiling.} No acceptor that holds a goal at instability rate
  \(\Lambda\) while it acts reaches \(\iota\) above
  \(p/\big(k_B T[\ln 2 + \sqrt{2\Lambda/\langle A\rangle}]\big)\) for
  one-bit acts, at any act time. The general ceiling is
  \(p\,b/\big(b\,k_B T\ln 2 + 2\sqrt{\kappa\,k_B T\Lambda}\big)\). It
  depends only on \(T\), \(p\), \(b\), \(N\), \(\varepsilon\) and the
  ratio \(\Lambda/\langle A\rangle\).
\item
  \textbf{Optimal act time.} Let an acceptor's measured joules per act
  follow \(J(\tau) = J_0 + a/\tau + s\,\tau\). Measure \(a\) and \(J_0\)
  from a sweep of \(\tau\) with no goal held, and the standing power
  \(s\) from the goal held with no acts. Then the measured
  \(\iota(\tau)\) peaks at \(\tau^* = \sqrt{a/s}\), at the value
  \(p\,b/(J_0 + 2\sqrt{a s})\), with \(a \ge \kappa\) and
  \(s \ge k_B T\Lambda\). The error band is
\end{enumerate}

\[\frac{\delta\tau^*}{\tau^*} = \frac{1}{2}\sqrt{\Big(\frac{\delta a}{a}\Big)^2 + \Big(\frac{\delta s}{s}\Big)^2}\]

\begin{enumerate}
\def\labelenumi{\arabic{enumi}.}
\setcounter{enumi}{2}
\tightlist
\item
  \textbf{Scaling.} Across a sweep of the instability rate,
  \(\ln\tau^*\) falls against \(\ln s\) with slope \(-1/2\).
\end{enumerate}

\textbf{Falsifier.} The measured maximizer of \(\iota(\tau)\) lies
outside \(\tau^* \pm 2\,\delta\tau^*\), or the fitted slope of
\(\ln\tau^*\) against \(\ln s\) over at least one decade of \(s\) lies
outside \(-1/2\) by more than two standard errors. Either defeats the
claim that act joules and standing joules add on one meter, with
\(\iota\) read as their density. A measured \(\iota\) above the ceiling
defeats the accounting of that run (F12, Section 8.7).

Figure 2 draws the curve of Prediction 1.

\includegraphics[width=0.95\linewidth,height=\textheight,keepaspectratio]{fig/act-time.png}

\emph{Figure 2. An optimal act time. Model illustration, not a
measurement: \(J(\tau) = J_0 + a/\tau + s\,\tau\) with \(J_0 = a = 1\)
and \(s \in \{1, 4\}\) in normalized units. Quadrupling the standing
cost halves \(\tau^*\), the slope of \(-1/2\) that the falsifier tests.
A \href{https://research.openie.dev/figures/agency/act-time.png}{PNG
copy} is provided.}

\textbf{Against the closest work.} A finite optimal cycle time is not
new by itself. Paneru and colleagues measured a power-optimal cycle time
in a Brownian information engine {[}330{]}. Panizon derives optimal
measurement schedules as a function of measurement cost {[}331{]}, and
Rolandi and colleagues bound an energy, time and accuracy trade-off in
thermodynamic computing {[}332{]}. What Prediction 1 adds is the source
of the interior optimum, a standing cost of holding a goal against
drift, and its scaling: \(\tau^*\) falls as \(\Lambda^{-1/2}\). The
slope test is the distinctive falsifier.

\textbf{What is new.} The finite-time Landauer bound prices the memory
alone and falls toward \(k_B T\ln 2\) as \(\tau\) grows {[}29{]}. There,
the slowest reset is the cheapest. The data-rate theorem prices holding
alone, per second (Section 4.1). Neither gives an interior optimum. The
definition meters the whole agent, counts confirmed bits only, and reads
\(\iota\) as a density on the joule axis. Then holding a goal makes
slowness expensive, and the best act time is the geometric mean of the
jump time \(1/\langle A\rangle\) and the drift time \(1/\Lambda\), up to
the factor \(1/\sqrt{2}\). The controller's own costs enter the joules
in the same way in the work of Ehrich, Still and Sivak (Section 5.1).

\textbf{Where to test it.} A colloidal feedback trap imposes virtual
potentials in software and computes work from trajectories at the
\(k_B T\) scale (Section 3.2). One coordinate runs a one-bit act cycle
in a virtual double well with reset time \(\tau\). The other is held at
the top of a virtual inverted well of rate \(\Lambda\). A
single-electron box with a charge detector runs the same act cycle with
every jump counted, so \(\langle A\rangle\) is measured (Section 3.2).
Section 8.8 gives the protocol.

\paragraph{Proposition 8. The gated
ranking}\label{proposition-8.-the-gated-ranking}

\textbf{Setting.} Fix a task \(q\), its predicate \(C(z)\) and its
allowed failure rate \(\delta\). For each system \(S\),
\(J^*_S(q) \in (0, \infty]\) is the joules to correct completion of
Section 1, with \(J^*_S(q) = \infty\) when no policy of \(S\) passes the
gate. Define \(S_1 \preceq_q S_2\) when
\(J^*_{S_1}(q) \le J^*_{S_2}(q)\).

\textbf{Statement.}

\begin{enumerate}
\def\labelenumi{\arabic{enumi}.}
\tightlist
\item
  \(\preceq_q\) is a total preorder on systems: reflexive, transitive
  and total. The ranking \(\succ_q\) of Section 2.3 is its strict part.
\item
  Every system that passes the gate ranks above every system that fails
  it.
\item
  Let every reading of a policy's mean joules carry relative error at
  most \(\epsilon < 1\). If
  \((1+\epsilon)\,J^*_{S_1}(q) < (1-\epsilon)\,J^*_{S_2}(q)\), then
  \(S_1\) ranks above \(S_2\) under every admissible reading.
\item
  Integrated \(\iota\) does not induce this order. There are systems
  with \(\iota_{S_1} > \iota_{S_2}\) and \(S_2 \succ_q S_1\).
\end{enumerate}

\textbf{Proof.}

\begin{enumerate}
\def\labelenumi{\arabic{enumi}.}
\tightlist
\item
  \(\preceq_q\) is the pullback of the total order \(\le\) on
  \((0, \infty]\) by the map \(S \mapsto J^*_S(q)\). A pullback of a
  total order is a total preorder. Its strict part is
  \(J^*_{S_1} < J^*_{S_2}\), which is the ranking of Section 2.3.
\item
  Every finite value is below \(\infty\).
\item
  Each policy's reading lies in
  \([(1-\epsilon)J_\pi, (1+\epsilon)J_\pi]\), so the minimum of the
  readings over the feasible policies lies in
  \([(1-\epsilon)J^*, (1+\epsilon)J^*]\). Then
  \(\hat J^*_{S_1} \le (1+\epsilon)J^*_{S_1} < (1-\epsilon)J^*_{S_2} \le \hat J^*_{S_2}\).
\item
  Let \(S_1\) make one attempt per task at 0.1 J that passes with
  probability 1/2 and never retries. Let \(S_2\) make one attempt at 1 J
  that always passes. Each pass credits \(b\) bits. Then
  \(\iota_{S_1} = (b/2)/0.1 = 5b\) bits per joule and
  \(\iota_{S_2} = b\) bits per joule. For any \(\delta < 1/2\), \(S_1\)
  never meets the gate, so \(J^*_{S_1} = \infty\), while
  \(J^*_{S_2} = 1\) J. \(\blacksquare\)
\end{enumerate}

\textbf{Reading.} Part 4 is the mechanism of R-F6 in Section 8.11: cheap
failures raise integrated \(\iota\) and leave \(J^*\) unchanged or
infinite. That is why \(\iota\) is read as a density inside a run and
\(J^*\) is the ranking across systems. Part 3 settles P6 up to a
validated \(\epsilon\) for each meter.

\paragraph{Proposition 9. When cheapest first is
optimal}\label{proposition-9.-when-cheapest-first-is-optimal}

\textbf{Setting.} A task has gears \(g_1, \dots, g_n\). An attempt with
gear \(g_i\) costs \(c_i > 0\) joules, paid in full, and passes \(C(z)\)
with probability \(p_i\). A policy tries gears in some order, stops at
the first pass, and refuses after the last. Two structures are treated.
Independent: passes of different gears are independent events. Nested:
if \(g_i\) passes then every later gear passes, with pass probabilities
\(q_1 \le \dots \le q_n\).

\textbf{Statement.}

\begin{enumerate}
\def\labelenumi{\arabic{enumi}.}
\tightlist
\item
  Independent gears. The expected joules of an order are least when
  gears are tried in decreasing order of \(p_i/c_i\). Cheapest first is
  optimal exactly when the order of rising cost is also an order of
  falling \(p_i/c_i\).
\item
  Nested gears, two of them, with the heavy gear sufficient
  (\(q_H = 1\)). Trying the light gear first lowers the expected joules
  if and only if \(q_L\,c_H > c_L\).
\item
  Budget. Under a hard budget \(B\) with \(c_H \le B < c_L + c_H\),
  light first closes the task with probability \(q_L\). Heavy alone
  closes it with probability 1.
\end{enumerate}

\textbf{Proof.}

\begin{enumerate}
\def\labelenumi{\arabic{enumi}.}
\tightlist
\item
  The expected joules of an order \(\sigma\) are
  \(\sum_k c_{\sigma(k)} \prod_{j<k}(1 - p_{\sigma(j)})\). Swap two
  adjacent gears \(i\) then \(j\) that follow a prefix with survival
  probability \(P\). Their contribution changes from
  \(P[c_i + (1-p_i)c_j]\) to \(P[c_j + (1-p_j)c_i]\), and nothing else
  changes. So \(i\) first is no worse if and only if
  \(p_i c_j \ge p_j c_i\), that is \(p_i/c_i \ge p_j/c_j\). Any order
  not sorted by this ratio has an adjacent pair out of order, and
  swapping it does not raise the expected joules. Finitely many swaps
  sort the order. This is the ratio rule of Simon and Kadane for search
  with all-or-none solutions {[}318{]}.
\item
  Light first costs \(c_L + (1-q_L)\,c_H\), because the heavy gear runs
  only when the light one fails and then passes. Heavy alone costs
  \(c_H\). The first is smaller if and only if \(c_L < q_L\,c_H\).
\item
  A failed light attempt leaves \(B - c_L < c_H\), so the heavy gear
  cannot run. \(\blacksquare\)
\end{enumerate}

\textbf{Reading.} The fixed cascade, cheapest sufficient first in a
fixed order, is optimal when cheap gears pass often enough to pay for
their failures: \(q_L c_H > c_L\) at each step of a nested cascade, or
rising cost matched by falling \(p/c\) among independent gears. When the
condition fails, going straight to the heavier gear wins, and under a
binding budget it can be the only way to close. That is the mechanism of
the selector result in Section 8.11. The σ-law selector estimates the
ratio with calibrated pass probabilities. It is a value-of-computation
rule with joules as the cost, in the sense of Russell and Wefald and of
Horvitz (Section 5.1). It stops when the next act no longer pays for its
joules. Both rules stop at the floor past which more joules buy no
verifiable progress, so both are Mixture of Limits instances. P7
(Section 7.1) lists what remains open: a regret bound when pass
probabilities are estimated or sufficiency is not monotone, and the
nested case with more than two gears, which is a dynamic program over
increasing subsequences of gears.

\subsubsection{2.4 Where the studies sit}\label{where-the-studies-sit}

\begin{itemize}
\tightlist
\item
  \textbf{Mixture of Limits} is the optimum operator. The cheapest
  sufficient gear is the argmin over joules. It is a philosophy, not one
  fixed order: spend stops at the plateau, where more is waste.
\item
  \textbf{The receipt} carries \(R^*\), \(R\), \(m\) and \(J\) for every
  act.
\end{itemize}

\subsection{3. Layer 1. The proposed Universal Law of Agency: a
definition plus a
gate}\label{layer-1.-the-proposed-universal-law-of-agency-a-definition-plus-a-gate}

The definition has four clauses: prediction fixed before the act,
coupling to the world through intervention, consequence borne by the
system, and a physical floor in joules with refusal on a receipt.

\subsubsection{3.0 Objections to the
floor}\label{objections-to-the-floor}

The floor used here is contested, and the assumptions (R1) to (R4) of
Proposition 1 are where the objections are met. Norton argues that
Landauer's principle holds only for protocols with a thermodynamically
irreversible step {[}326{]}. Maroney and Sagawa separate logical from
thermodynamic irreversibility: \(k_B T\ln 2\) prices the entropy change
of a memory, not a bit as such {[}327, 328{]}. López-Suárez, Neri and
Gammaitoni operated an irreversible logic gate below \(k_B T\)
{[}329{]}. Wolpert and colleagues argue that the costs of real computers
are dominated by mismatch, modularity and clocking rather than by
erasure {[}320{]}. Each objection narrows the claim rather than reverses
it. Proposition 1 is stated only for a record that meets (R1) to (R4), a
reset to a fixed state with equal-energy states and no shared
interaction energy, which is the case in which the bound holds; the
floor is a lower bound, and every confirmed act of the runs in Section
8.11 that the meter resolved sits more than 14 orders of magnitude above
it. K5 has a similar scope: Jordan shows that the Margolus-Levitin bound
does not limit the rate of computation by energy without an assumption
on information density {[}324{]}, so K5 is read as a bound on orthogonal
state changes, not on clock speed.

\subsubsection{3.1 Proven}\label{proven}

\textbf{Information has a price.} Landauer showed that a logically
irreversible operation dissipates at least \(k_B T \ln 2\) per bit
{[}6{]}. This is the floor under clause 4. Bennett built a reversible
Turing machine and showed that only erasure must cost {[}35{]}. The
unavoidable cost of agency sits at the reset of the acceptor, not at the
thinking. Bennett then resolved Maxwell's demon by the cost of erasing
the demon's memory {[}7{]}. The demon is a minimal agent. Its accounts
close only when its memory is reset. That is clause 3 plus clause 4.

\textbf{Feedback is priced by information.} Sagawa and Ueda proved that
with feedback, extractable work is bounded by \(-\Delta F + k_B T I\)
{[}3{]}. The value of acting on a measurement is capped by the
information the measurement carries. They proved that measurement work
plus erasure work is at least \(k_B T I\) {[}5{]}. The full cycle of
predict, compare and reset is priced in joules. They derived an exact
Jarzynski equality under feedback {[}4{]} and a fluctuation theorem with
information exchange between subsystems {[}36{]}. Coupling is a
correlation that enters the second law. Hartle, Wolpert, Stier, Kempes
and Manzano priced feedback with noisy measurements, limited memory and
a limited repertoire of protocols. Against random action as the
baseline, the benefit of feedback can exceed the information gained by
orders of magnitude {[}37{]}. The bits bound the work against free
energy. The value against a baseline is read off the meter.

\textbf{Entropy converts to joules.} Esposito and Van den Broeck derived
\(k_B T \Sigma = W - \Delta F_{\mathrm{neq}}\) for an isothermal driven
system {[}24{]}. Every kinematic bound in Layer 2 becomes a joule bound
through this identity. Horodecki and Oppenheim proved work costs for
single shots at small scale {[}38{]}. Agency acts are single shots. The
standard reviews state the settled account: information is a
thermodynamic resource with a price {[}39, 40{]}.

\textbf{Control needs coupling.} Conant and Ashby proved that a good
regulator must be a model of the system it regulates {[}41{]}. A system
that regulates well already carries a prediction. Touchette and Lloyd
proved that the entropy a controller removes is bounded by its mutual
information with the system {[}32, 33{]}. Coupling is the currency of
control. Closed loop control, acting on a comparison, is what the
definition requires. Pearl's do-calculus gives the intervention operator
and the conditions under which its effects are identifiable {[}1{]}.
Clause 2 is written in this operator. Everitt and colleagues classified
how an agent can tamper with its own reward or inputs {[}2{]}. Coupling
must be to the world, not to the agent's own sensor. Richens and Everitt
proved that any agent that meets a regret bound under a large set of
distributional shifts has learned an approximate causal model of the
process that generates its data {[}42{]}. Richens, Abel, Bellot and
Everitt proved that any agent that generalizes to multi-step
goal-directed tasks has learned a predictive model of its environment,
and that the model can be extracted from the policy {[}43{]}. A capable
agent carries a prediction. That is clause 1, proven for modern agents.

\textbf{Information flows inside the agent.} Horowitz and Esposito
showed that information flow between subsystems enters each subsystem's
second law {[}44{]}. Coupling becomes a rate. Ito and Sagawa proved a
second law for arbitrary causal networks with transfer entropy {[}45{]}.
Prediction, act and outcome are nodes on such a network. Hartich, Barato
and Seifert linked transfer entropy, learning rate and dissipation
{[}46{]}. Allahverdyan, Janzing and Mahler bounded the efficiency of
information flow between coupled systems {[}47{]}. Esposito and Schaller
showed that feedback that changes only barriers still enters the entropy
balance {[}48{]}. A purely informational act is in the joule ledger.

\textbf{Autonomous demons.} Mandal and Jarzynski solved a demon that
lifts a mass by writing bits to a tape {[}49{]}. It has no outside
operator, and memory pays for its work. Boyd, Mandal and Crutchfield
bounded the work of such ratchets by the entropy rates of their input
and output tapes {[}50{]}. They showed that modular designs pay
dissipation above Landauer {[}51{]}. How an agent is built changes its
floor. Wolpert and Kolchinsky computed the extra entropy production a
gate pays inside a circuit, set by the circuit's topology {[}52{]}.
Wiring is part of the floor.

\textbf{Prediction quality has a joule price.} Kolchinsky and Wolpert
proved that a process tuned for one input distribution dissipates extra
on any other, by a relative entropy term {[}53{]}. Still, Sivak, Bell
and Crooks showed that memory that does not predict costs dissipation
{[}54{]}. Keeping only what predicts is the cheap design. Kolchinsky and
Wolpert defined semantic information as information that is causally
necessary for a system to keep itself out of equilibrium, with a
thermodynamic value {[}55{]}. It is the closest existing account with
joules. It stops before kinematics and before an intelligence operator.
Ortega and Braun derived bounded-rational decisions from a free-energy
principle with an information cost {[}56{]}. Kolchinsky and Wolpert
bounded the heat of universal computation by program complexity
{[}57{]}. Kolchinsky and Wolpert showed that constraints on the
available protocols limit extractable work and the value of information
{[}58{]}. Manzano and colleagues derived fluctuation relations and
second-law-like inequalities for computations that halt at random times
{[}59{]}. A run that stops when its predicate passes is such a
computation.

\textbf{Coupling and boundary rest on theorems.} Ay and Polani defined
causal information flow through interventions {[}8{]}. Transfer entropy
does not measure that flow, shown by example {[}11{]}. The Markov
blanket condition holds only in a narrow range of parameters, by
derivation {[}15{]}, and its definitions in use differ, by
counterexample {[}16{]}. Every bipartite stochastic process has a causal
blanket, by theorem {[}22{]}. Definitions 1 and 2 stand on these
results.

\subsubsection{3.2 Demonstrated}\label{demonstrated}

\textbf{Landauer's floor, measured.} A colloidal bead in an optical
double well approached \(k_B T \ln 2\) of mean heat as erasure slowed
{[}60{]}. A feedback trap measured the same approach with high precision
and extra cost at short cycle times {[}61{]}. The floor rises when the
act is fast. That is a Layer 2 effect seen in a Layer 1 experiment. An
adiabatic electronic circuit showed the Landauer cost for irreversible
operation and dissipation below \(k_B T\) for reversible operation
{[}62{]}. The floor binds the reset, not the logic. Nanomagnetic memory
bits showed the floor in memory hardware {[}63{]}. Erasure with
near-zero work appears when the free-energy accounting is complete
{[}64{]}. The floor is on total entropy, so the receipt must carry the
full account. The entropy that enters the bound was measured to have the
Gibbs-Shannon form {[}65{]}. A single trapped ion showed the quantum
floor {[}66{]}. The floor holds on every substrate tested. An
underdamped micromechanical oscillator reached the floor at high speed
{[}67{]}. A quantum dot erased bits at minimal dissipation for a given
duration, with protocols designed from thermodynamic length {[}68{]}.

\textbf{Feedback agents turn information into work.} A colloid on a
spiral staircase potential gained free energy above the work done, paid
by information, and confirmed the generalized Jarzynski equality
{[}69{]}. This is the first direct observation of an act whose value is
set by a measurement made before it. A single-electron box extracted
work near \(k_B T \ln 2\) per bit {[}70{]} and measured the fluctuation
theorem with mutual information {[}71{]}. An autonomous on-chip demon
cooled a system with information while its own heating was measured
{[}72{]}. The agent's cost was measured beside its effect.

Quantum agents obey the same bounds. NMR on molecular spins reduced
entropy production by feedback, within the information bound {[}73{]},
and converted information to energy at the floor {[}74{]}.
Superconducting circuits tracked a demon's memory during feedback
{[}75{]}, verified the generalized integral fluctuation theorem
{[}76{]}, and measured information gained and lost during feedback
{[}77{]}. Coupling decays, so the receipt must record it at the time of
the act. Cold atoms in a three-dimensional optical lattice were sorted
by measurement {[}78{]}. Linpeng and colleagues measured the energetic
cost of quantum measurement with quantum, coherent and thermal light in
a circuit quantum electrodynamics setup. Single-photon light cost the
least energy per unit of information gained {[}79{]}. The measurement
that writes a record has a measured price.

Classical engines measured the price at full scale. An optical tweezer
engine converted information to work at the theoretical limit {[}80{]}.
A colloidal engine extracted work that depends on temporal correlations
in its input {[}81{]}. A DNA hairpin under continuous monitoring yielded
more than \(k_B T \ln 2\) per cycle when many measurements happened per
cycle, consistent with the information content {[}82{]}. A levitated
microparticle with delayed feedback obeyed a generalized second law
across two decades of delay, and feedback failed at large delay
{[}83{]}. A late act is a weaker act. Feedback trap engines maximized
power and velocity {[}84{]}, exploited noisy measurements optimally
{[}85{]}, and harvested a nonequilibrium bath {[}86{]}. A quantum dot
linked work fluctuations to dissipation during information-to-work
conversion {[}87{]}. A trapped ion verified fluctuation theorems that
include the demon's own dissipation {[}88{]}, following theory by Zeng
and Wang {[}89{]}.

\textbf{Functional-system agents built and run.} A simulated nematode
with an acceptor of action results learned locomotion and chemotaxis
{[}90{]}. A physical mobile robot ran the same functional-system
controller {[}91{]}. Clauses 1 and 3 run in software and on a body.
Neither meters joules.

\subsubsection{3.3 Positions in the
literature}\label{positions-in-the-literature}

\textbf{Purposive systems and anticipation.} Rosenblueth, Wiener and
Bigelow defined purposeful behavior through negative feedback toward a
goal {[}92{]}. It has no energy term. Anokhin's functional system has
afferent synthesis, decision, an acceptor of action results fixed before
the act, action, and comparison {[}93, 94{]}. This is clauses 1 and 3 in
their first systematic form, built on animal experiments. Sudakov stated
the result as the system-forming factor {[}95{]}. The prediction clause
is the organizing principle. Shvyrkov showed from single-neuron
recordings that the goal organizes neural activity {[}96{]}. Rosen
defined an anticipatory system as one that contains a predictive model
of itself and its world and acts on the prediction {[}97{]}. None of
these sets a joule floor. Sterling defines allostasis as regulation by
prediction {[}98{]}. Physiology puts the prediction before the act.
Virgo, Biehl and McGregor gave conditions under which a system's states
can be read as beliefs updated by Bayes' theorem {[}99{]}. Biehl and
Virgo call a system an agent when it can be read as solving a POMDP
(partially observable Markov decision process) {[}100{]}. Both are
formal definitions. Clause 1 asks for less and fixes more: one template,
fixed at \(t_0\), hashed on the receipt.

\textbf{Autonomy and organization.} Maturana and Varela defined the
living by self-production {[}101{]}. Kauffman put a thermodynamic work
cycle inside the definition of an autonomous agent {[}102{]}. It has no
prediction clause. Di Paolo made adaptivity, the regulation of one's own
viability, a condition for sense-making {[}103{]}. That is clause 3.
Barandiaran and Moreno proposed a minimal criterion for cognitive
organization and grounded adaptive regulation in metabolism {[}104,
105{]}. Barandiaran, Di Paolo and Rohde defined agency by individuality,
interactional asymmetry and normativity {[}106{]}. Ruiz-Mirazo, Moreno
and Mossio framed autonomy as a self-maintaining metabolic organization
{[}107, 108{]}. Montévil and Mossio described organisms as closures of
constraints on thermodynamic flows {[}109{]}. It is the closest of this
school to joules. Aguilera and Barandiaran built a minimal model of
autonomous agency in stochastic thermodynamics {[}110{]}. The autonomy
tradition is moving toward the joule. Barandiaran and Almendros applied
individuality, normativity and interactional asymmetry to large language
models. A model fails the first two and partly fails the third
{[}111{]}. The definition applies its gate per act: a model call is an
act when its receipt carries \(R^*\), \(m\) and \(J\). Levin's TAME
(Technological Approach to Mind Everywhere) treats agency as graded and
testable by experiment at every scale {[}112{]}. The definition keeps
the gate binary per act and puts the grading in \(X\).

\textbf{Free energy and the physics of agency.} Friston proposes that
agents minimize variational free energy {[}113, 114{]}. That free energy
is informational, not joules, and every bounded system qualifies, so it
sets no floor. Rovelli explains agency through the entropy gradient
{[}115{]}. Agency is thermodynamic at root. Jaeger argues that agency
requires self-manufacture and that current artificial intelligence lacks
it {[}116{]}. That is a demand for a gate before any intelligence claim.
Azadi ties agency to computational irreducibility {[}117{]}.
Wissner-Gross and Freer propose a force toward future path diversity
{[}118{]}. Its causal temperature is not a bath temperature, so its
units are not joules. Kappen's comment critiques it {[}119{]}. Fields,
Goldstein and Sandved-Smith separate thermodynamic free energy, in
joules, from variational free energy in active inference agents, and
relate the two {[}120{]}. Only the joule quantity enters K1. Bruineberg
and colleagues separate the Markov blanket as a tool of inference from
the blanket as a physical boundary {[}17{]}. Definition 2 takes the
boundary from the definition. Bartlett and colleagues make goal-directed
information processing, measured by semantic information, the
distinctive feature of living systems {[}121{]}. That is clauses 2 and 3
as a criterion for life.

\subsection{4. Layer 2. The kinematic bounds of
agency}\label{layer-2.-the-kinematic-bounds-of-agency}

Once an act is eligible, physics bounds how fast agency moves per second
and per joule, and what precision and speed cost. The bounds come in
four families: speed limits, the TUR, finite-time costs and rate
ceilings.

\subsubsection{4.1 Proven}\label{proven-1}

\textbf{Speed limits.} Shiraishi, Funo and Saito proved the classical
speed limit in Section 2.2 {[}23{]}. Moving the agent's state
distribution a distance \(L\) in time \(\tau\) costs entropy, and so
joules. The bound extends to open quantum systems {[}122{]} and to
strong coupling with general environments {[}123{]}. Vo, Van Vu and
Hasegawa derived the speed limit and the TUR from one inequality
{[}124{]}. Falasco and Esposito bounded the time to complete a
transition by the dissipation spent {[}125{]}. Ito bounded the
information-geometric speed of a distribution by entropy production
{[}126{]}. Ito and Dechant bounded the rate of change of any observable
by Fisher information and dissipation {[}127{]}. Nakazato and Ito
bounded entropy production by the Wasserstein distance traveled
{[}128{]}. Dechant, Sasa and Ito split entropy production into the part
that moves the system and the part that keeps it running {[}129{]}. Van
Vu and Saito unified the TUR, minimum dissipation and speed limits
through optimal transport {[}130{]}. Nagayama, Yoshimura and Ito derived
a family of speed limits indexed by generalized means of activity
{[}131{]}. Yoshimura and Ito proved the TUR and the speed limit in
deterministic chemical reaction networks {[}132{]}. Aurell,
Mejía-Monasterio and Muratore-Ginanneschi solved minimal finite-time
dissipation as an optimal transport problem {[}133{]}. The kinematic
bounds form one family. The cheapest way to move an agent's state is a
transport map. Hasegawa derived the speed limit, the TUR and the
Heisenberg principle from one bulk-boundary correspondence {[}134{]}.
Van Vu and Hasegawa bounded irreversible entropy production by a
modified Wasserstein distance in classical and quantum Markov systems
{[}135{]}. Delvenne and Falasco bounded entropy production by the
statistics of kinetic observables {[}136{]}.

\textbf{Precision costs dissipation.} Barato and Seifert stated the TUR,
\(\mathrm{Var}(Y)/\langle Y\rangle^2 \ge 2/\Sigma\) {[}25{]}. A reliable
agent pays for its reliability. From \(2/\epsilon^2\) at
\(\epsilon = 0.01\), precision of 1\% costs at least 20,000 \(k_B T\)
{[}137{]}. Gingrich, Horowitz, Perunov and England proved the TUR for
all steady-state currents {[}26{]}. Horowitz and Gingrich proved it for
finite observation times {[}138{]}, which is what agents have. Further
results bound the full distribution of fluctuations {[}139{]}, cover
discrete time steps {[}140{]}, extend the TUR wherever a fluctuation
theorem holds {[}141{]}, and cover Langevin dynamics, underdamped motion
and velocity feedback {[}142, 143, 144{]}. Correlations between currents
tighten the bounds {[}145{]}. Potts and Samuelsson showed that any
fluctuation relation implies a TUR, including with measurement and
feedback {[}28{]}. This is the precision bound for agents. A review
states the field {[}146{]}. Barato and Seifert priced the precision of a
clock {[}147{]}. Salazar bounded entropy production by information from
the detailed fluctuation theorem {[}148{]}. Landi and Paternostro
reviewed entropy production from classical to quantum {[}149{]}.
Feedback forms now cover the cases agents meet. Van Vu and Hasegawa
bounded fluctuations under arbitrary control protocols by entropy
production and a kinetic term {[}150{]}. Tanogami, Van Vu and Saito
bounded the current of a subsystem by its own entropy production and the
information flow between it and its partner {[}151{]}. Kumasaki, Tojo,
Sagawa and Funo derived a TUR for feedback cooling {[}152{]}. Honma and
Van Vu derived a finite-time TUR for open quantum systems under
continuous monitoring and Markovian feedback, with quantum mutual
information beside entropy production {[}153{]}. For counting
observables, Garrahan bounded fluctuations by dynamical activity and
gave the matching bounds for first-passage times {[}34{]}. Gingrich and
Horowitz derived the conjugate relation for the first passage time to
accumulate a net current {[}154{]}. Proposition 3 applies these results
to confirmed agency.

\textbf{Finite-time costs.} Schmiedl and Seifert found optimal
finite-time protocols, with jumps at the ends {[}155{]}. The optimal act
at finite speed is not the slow act. Sivak and Crooks bounded excess
work by the squared thermodynamic length over the duration {[}30{]}.
Proesmans, Ehrich and Bechhoefer found the minimal cost of erasing a bit
in time \(\tau\) {[}29, 156{]}. The Layer 1 floor becomes a Layer 2
curve. Zhen and colleagues stated a universal finite-time bound on reset
energy {[}157{]}. Speed and efficiency trade against each other by
theorem: efficiency at maximum power {[}158, 159{]}, finite power
forbidding Carnot efficiency {[}160{]}, a three-way trade among power,
efficiency and constancy {[}161{]}, and a universal constraint for
low-dissipation engines {[}162{]}. Tsirlin and colleagues found minimal
dissipation at a fixed rate {[}163{]}. Lee, Lee, Kwon and Park tightened
the speed limit for highly irreversible processes and obtained a tight
finite-time Landauer bound {[}164{]}. Van Vu and Saito derived the
finite-time quantum Landauer principle and the role of coherence in it
{[}165{]}. Rolandi and Perarnau-Llobet derived the finite-time Landauer
principle beyond weak coupling {[}166{]}. Kamijima, Funo and Sagawa
derived the finite-time trade-offs between measurement and feedback
{[}167{]}. Tkachenko showed for physical neural networks that
quasi-static inference needs no work, while finite-speed inference costs
at least a transport distance (preprint) {[}168{]}. Inference has no
floor that survives slow operation. The reset does. Proposition 1 prices
it.

\textbf{Holding a goal.} Tatikonda and Mitter {[}169{]} and Nair and
Evans {[}170{]} proved the data-rate theorem: holding an unstable linear
system needs a channel rate above the sum of the base-2 logarithms of
its unstable eigenvalue magnitudes. Nair, Fagnani, Zampieri and Evans
reviewed the field {[}171{]}. Sandberg, Delvenne, Newton and Mitter
found the maximum work a continuously measuring demon extracts in finite
time, and showed that any implementation needs an external power source
{[}172{]}. Horowitz and Sandberg showed that the information terms in
second-law inequalities form a hierarchy, each the minimum cost of
acquiring that information by a distinct measurement {[}173{]}.
Proposition 2 turns the bit count into a joule floor per step. It is a
joule-priced corollary of the data-rate theorem, not a new law.

\textbf{Quantum feedback and measurement.} The feedback second law holds
for quantum systems {[}3{]}. Jacobs extended it to weak measurements
{[}174{]}. Funo, Watanabe and Ueda derived integral fluctuation theorems
under quantum measurement and feedback {[}175{]}. Guryanova, Friis and
Huber proved that ideal projective measurements need infinite resources,
and gave the energy cost of a pointer that reproduces the system's
statistics {[}176{]}. Danageozian, Wilde and Buscemi bounded the
measurement heat of error identification by the Groenewold information
gain {[}177{]}. These price the measurement that writes a record.
Proposition 4 keeps the floor of the record itself.

\textbf{Groups.} Wolpert derived a strictly nonzero lower bound on the
minimal entropy production rate of interacting subsystems from the
constraints on which subsystems each one depends on {[}178{]}. Tasnim
and Wolpert extended the theory to subsystems that change state together
{[}179{]}. Wolpert related the entropy production of a whole Bayes net
to the precision of currents in its parts {[}180{]}. Rolandi, Abiuso and
Perarnau-Llobet showed that collective protocols can make dissipated
work grow sublinearly with the number of members {[}181{]}. Fujimoto and
Ito found by a game-theoretic method that the partial entropy
productions of interacting subsystems trade against each other, and that
the total is least when the subsystems share the penalty equally
{[}182{]}. Proposition 7 composes these joules with the group's coupled
bits.

\textbf{Rate ceilings.} Margolus and Levitin bounded orthogonal
transitions per second by energy {[}31{]}. Lloyd applied the limit to a
whole computer {[}183{]}. Deffner and Campbell reviewed the family of
quantum speed limits {[}184{]}.

\textbf{Living sensors.} Energy dissipation, adaptation speed and
accuracy are bound together in bacterial chemotaxis {[}185{]}. Learning
about the environment requires energy {[}186{]}, and inference accuracy
is tied to dissipation {[}187{]}. Sensing precision is limited by
receptors, time and energy, with an optimal allocation among them
{[}188{]}. No resource can be skipped. That fits Mixture of Limits.
Erasure inside adaptation has a cost {[}189{]}. Information acquired per
energy dissipated is the nearest existing quantity to \(v_J\) in living
sensors {[}190, 191{]}. Copying information costs energy {[}192{]}.
Writing the receipt has a price. Accurate and synchronized biochemical
clocks cost free energy {[}193, 194{]}, and so does reducing noise at
high sensitivity {[}195{]}. Holding a system away from equilibrium has a
minimum power {[}196{]}. That is the standing cost of an agent ready to
act. Bryant and Machta bounded the energy of sending a bit through a
cell's physical channels, in \(k_B T\) per bit, as a function of size,
distance and latency, and found it many orders of magnitude above unity
{[}197{]}. Tjalma and colleagues showed that the bits of the past that
best predict the future are prohibitively costly for cellular networks
{[}198{]}.

\textbf{Information gained by selection.} Selection also writes bits,
and its rate is bounded. Kimura measured the information that natural
selection accumulates in a genome by the substitutional load and, with
the load values he took as representative, estimated about 0.29 bits per
generation {[}341{]}. Worden bounded the rate at which selection can add
information to a genome in terms of the selection applied per generation
{[}342{]}. Hledík, Barton and Tkačik showed that, in stochastic
populations, the rate is bounded by the population size times the cost
of selection, across Wright-Fisher, Moran and diffusion models
{[}343{]}. Every such bound is a rate per generation. A generation
lengthens as an organism grows, so a carrier that learns only through
selection gains coupled bits ever more slowly per second. Nervous
systems write bits within a lifetime, and cultural records write them
across lifetimes. The kinematics of Section 2.2 apply to each carrier.
What changes from one carrier to the next is the time per written bit,
which is the reason agency moves from genomes to brains to tools when
the world changes faster than generations turn over.

Figure 3 orders the three carriers.

\includegraphics[width=0.95\linewidth,height=\textheight,keepaspectratio]{fig/carriers.png}

\emph{Figure 3. Where coupled bits are written. Qualitative diagram; the
one number is Kimura's estimate {[}341{]}. Bounds: Worden {[}342{]};
Hledík, Barton and Tkačik {[}343{]}. A
\href{https://research.openie.dev/figures/agency/carriers.png}{PNG copy}
is provided.}

\textbf{Information rates without joules.} Schreiber defined transfer
entropy, directed information flow per step {[}199{]}. It is per time,
not per joule, and it has no eligibility gate. Tishby and Polani charged
information per step in a Bellman recursion {[}200{]}. Landauer converts
those bits to joules. Stratonovich bounded the maximal gain from a given
amount of information {[}201{]}. That is a ceiling on value per coupled
bit.

\subsubsection{4.2 Demonstrated}\label{demonstrated-1}

\textbf{The TUR tested.} A two-qubit NMR experiment obeyed generalized
TURs and violated the specialized TUR where theory predicts {[}202{]}.
An optical tweezer information engine violated the original TUR near
maximal efficiency and obeyed generalized bounds with mutual information
{[}27{]}. Agents need the feedback form. A quantum dot Szilard engine
confirmed it independently {[}203{]}. Atomic-scale conductors examined
the TUR with measured current noise, the physics of every chip
{[}204{]}. Molecular motors were scored by precision per unit
dissipation {[}205{]}. Observed current fluctuations bound dissipation
from below {[}206{]}. A receipt of fluctuations is a lower-bound meter
for joules.

\textbf{Speed limits and finite-time costs tested.} Single atoms in an
optical lattice showed the Mandelstam-Tamm and Margolus-Levitin bounds
and the crossover between them {[}207{]}. The rate ceiling is measured.
A trapped colloid measured time and entropy trade-offs for fast thermal
transitions {[}208{]}. Brownian particles driven by optimal transport
protocols saturated the minimal finite-time dissipation bound {[}209{]}.
The kinematic optimum is reachable. A feedback trap measured a two-force
Brownian machine against linear response bounds {[}210{]}. A photonic
experiment tested geometric bounds on entropy production {[}211{]}. Dago
and Bellon measured erasure in a double-well memory at high speed and
traced the overhead above Landauer's bound to its sources {[}212{]}.

\textbf{Living sensors measured.} Mattingly, Kamino, Machta and Emonet
measured the rate at which E. coli acquires information during
chemotaxis, and showed that the cells climb gradients close to the limit
that rate sets {[}213{]}. Information limits behavior, measured.

\textbf{Heat engines at finite speed.} A single colloid ran a micrometer
Stirling engine {[}214{]}. A trapped bead ran a Brownian Carnot cycle
with efficiency at maximum power measured {[}215{]}. The energy cost of
choosing one of two options was measured {[}216{]}. That is the smallest
act. A colloidal engine ran on an active bacterial bath {[}217{]}, and
engineered noise gave near-Carnot efficiency at finite power by
shortening relaxation {[}218{]}. Engineering activity moves the bound.
An information engine flipped from refrigerator to heater with noise
{[}219{]}. The receipt must record the regime, not only means.

\textbf{Empowerment in simulation.} Klyubin, Polani and Nehaniv defined
empowerment, the channel capacity from actions to future sensors
{[}9{]}. It was computed for continuous control {[}220{]} and reviewed
{[}221{]}. It measures how much an agent can move the world, in bits. It
is not a rate per joule.

\subsubsection{4.3 Positions in the
literature}\label{positions-in-the-literature-1}

Kinematics is mostly theorem and experiment. The opinions here concern
which variable is right. Seifert holds that inference from fluctuations
is the route to hidden costs {[}137{]}. Rovelli and Wissner-Gross and
Freer both propose entropic kinematics {[}115, 118{]}. Kolchinsky shows
that dissipation alone does not bound how fast replicators grow or decay
{[}222{]}. Rates are bounded by dissipation together with activity. That
is the form of the speed limit, where \(\langle A\rangle\) sits beside
\(\Sigma\). Adlam, McQueen and Waegell argue that a purely quantum
system cannot be an agent, because building a world model and
deliberating need copying, which the no-cloning theorem forbids
(preprint) {[}223{]}. Clause 3 already asks for a classical record of
\(m\).

\subsection{5. Layer 3. Intelligence as the calculus of
agency}\label{layer-3.-intelligence-as-the-calculus-of-agency}

Intelligence is the derivative of gated agency in joules, read off a
completed run. This layer needs three things from the literature: theory
that ties learning and decision to joules, measurement of joules per
unit of useful work, and positions on what intelligence is.

\subsubsection{5.1 Proven}\label{proven-2}

\textbf{Learning and prediction priced in joules.} Goldt and Seifert
bounded the information a learning network acquires by the entropy it
produces {[}224{]}. Learning efficiency is at most one. This is
\(\iota\) for the learning part of an act, with a ceiling. Still showed
that optimal memory keeps only predictive information {[}225{]}. The
marginal joule spent on memory must buy prediction. Boyd, Crutchfield
and Gu proved that the agent that extracts the most work from data has
the maximum-likelihood model of that data {[}226{]}. Better models yield
more work per joule. Boyd, Crutchfield, Gu and Binder stated overfitting
and generalization as energetics {[}227{]}. Generalization appears as
joules. Ehrich, Still and Sivak priced the controller itself, beyond the
information it uses {[}228{]}. So \(J\) must be metered at the wall of
the whole system. Ehrich and Sivak gave the ledger of energy and
information flows in autonomous machines {[}229{]}. Wolpert reviewed
computation costs beyond Landauer {[}230{]}. Wolpert and colleagues set
out the open problems for the energy costs of computation in stochastic
thermodynamics {[}320{]}. Levy and Baxter showed that neural codes that
maximize bits per unit energy differ from codes that maximize bits
{[}231{]}. The right objective for a brain is information per joule.
Balasubramanian, Kimber and Berry set the objective for an exploratory
regime as transmission rate per unit power {[}232{]}. That is bits per
joule. Goldt and Seifert bounded the information a network learns about
a rule by the thermodynamic cost of learning, for batch and online
learning {[}233{]}. Fiderer, Barth, Smith and Briegel defined the work
capacity of an environment channel, the maximum rate at which any agent
can expect to extract work in a percept-action loop. Work-efficient
agents balance prediction against forgetting (preprint) {[}234{]}.
Elliott, Gu, Garner and Thompson showed that quantum adaptive agents can
need far less memory than memory-minimal classical agents {[}235{]}.
Less memory to reset is a lower floor to pay.

\textbf{Decisions and bounded optimality.} Russell and Subramanian
defined bounded optimality: the best program for a given machine and
environment {[}236{]}. Here the resource is joules, and the optimum is
\(\arg\min J\) subject to \(X \ge X_{\mathrm{req}}\). Russell and Wefald
priced each computation by its expected effect on the decision, net of
its cost, and stopped deliberation when no computation had positive net
value {[}316{]}. Horvitz framed the same trade for reasoning under
resource constraints {[}317{]}. The σ-law selector is that rule with the
cost in metered joules. Proposition 9 gives its exact form for a
cascade. Genewein and colleagues derived abstraction and hierarchy from
utility minus information cost {[}237{]}. Mixture of Limits gears are
such a hierarchy. Legg and Hutter defined intelligence as
complexity-weighted expected reward over all computable environments
{[}238{]}. It is an integral, but over reward, without energy, and
uncomputable. The derivative in joules replaces reward with confirmed
coupled bits and the measure with metered joules. Takahashi and Hayashi
give the closest existing calculus in joules {[}239{]}. They define two
ratios under explicit accounting conventions. Thermodynamic epiplexity
per joule is
\(I(W^{\mathrm{post}}; Z \mid W^{\mathrm{pre}})/E_{\mathrm{cons}}\): new
bits about an environment variable \(Z\) retained in the agent's state,
per joule consumed inside a stated boundary. Empowerment per joule is
the capacity of the channel from acts to observations per unit expected
energetic cost, the capacity-per-unit-cost of Verdú. They prove a
closed-cycle Landauer benchmark by joining a thermodynamic learning
inequality with data processing, and they show by a reversible copying
construction that, with an open boundary and fresh memory brought in
free, information gain and in-boundary dissipation decouple. They close
with a reporting checklist: boundary, energy balance, baseline,
coarse-graining, horizon and reset, throughput, estimator. This track
takes four things from that work. Their closed-cycle benchmark is K1 for
learning. Their decoupling construction is the reason for act-closed
metering in Definition 2 (Lemma 1). Their checklist is adopted in
Section 8.12. Their advice to report bits per joule beside bits per
second is the variable Prediction 1 optimizes. The definition differs
from their metrics in four exact ways. Empowerment per joule is a
capacity, a supremum over act distributions; the definition credits the
realized flow of the policy that ran (Definition 1), with capacity as
its bound. Their numbers are benchmark-relative efficiency reports, by
their own statement; this track ranks systems by \(J^*\) behind a
correctness gate (Proposition 8). They have no prediction fixed before
the act, no consequence clause and no refusal; the definition requires
all three before any bit counts. They give no kinematics of time;
Prediction 1 derives an optimal act time and a ceiling from the same
joule accounting. Ortega, Braun, Dyer, Kim and Tishby set out
information-theoretic bounded rationality {[}240{]}. Zénon, Solopchuk
and Pezzulo cast the cost of cognition as information {[}241{]}. Sims
applied rate-distortion theory to perception {[}242{]}. Their currency
is bits. Landauer converts bits to joules at the reset.

\subsubsection{5.2 Demonstrated}\label{demonstrated-2}

\textbf{Brains.} Laughlin, de Ruyter van Steveninck and Anderson
measured the energy cost per bit in blowfly photoreceptors and
interneurons {[}243{]}. Higher information rates cost more per bit. That
is a measured \(\iota\) in a living agent and a measured diminishing
return, \(\alpha_J < 0\) in sensing. Attwell and Laughlin built the
energy budget of grey matter by process {[}244{]}. Lennie showed that
energy limits how many cortical neurons can be active at once {[}245{]}.
Energy is the constraint, and the rest adapts. Energy per action
potential differs widely across neuron types {[}246{]}. Same function,
different joules. Neural design follows energy efficiency {[}247{]}.
Levy and Calvert audited the human cortex. Communication costs 35 times
as much as computation, and a neuron's computation sits a factor of
\(10^8\) from the best possible bits per joule {[}248{]}. That is a
distance to the floor. Harris and colleagues found that a visual pathway
synapse is sized to maximize information per unit energy, not
information {[}249{]}. The objective is bits per joule, measured at one
synapse. Padamsey, Katsanevaki, Dupuy and Rochefort found that food
restriction cut synaptic ATP (adenosine triphosphate) use in mouse
visual cortex by 29\% and broadened orientation tuning by 32\%
{[}250{]}. Fewer joules bought less precision. Plaçais and Preat showed
that starved flies switch off costly aversive long-term memory, and
forcing it back reduced survival {[}251{]}. Plaçais and colleagues found
that raised energy flux in the mushroom body triggers long-term memory
{[}252{]}. Mery and Kawecki found that forming long-term memory lowered
flies' resistance to desiccation {[}253{]}. Learning has a measured
metabolic price. Hechler, de Lange and Riedl measured lower cortical
oxygen use during confident prediction, by up to 12\% (preprint)
{[}254{]}. Prediction first can save joules. Malkin and colleagues gave
synapses an energy cost for reliability in trained networks and found
signatures of Bayesian inference in the optimum {[}255{]}. Li and van
Rossum showed that naive synaptic plasticity costs extreme energy and
that caching changes before consolidation saves it {[}256{]}. Lynn and
colleagues inferred broken detailed balance in the human brain from
neuroimaging, rising with physical and cognitive exertion {[}257{]}.
That entropy production is coarse-grained and informational. It is not a
joule meter.

\textbf{Machines.} Memory access costs far more energy than arithmetic
{[}258{]}. Parameter count matters only through the joules it makes you
move. Intelligence per watt measures task accuracy per unit power on
local accelerators {[}259{]}. Its numerator is benchmark accuracy, not
gated agency. MLPerf Power standardizes power measurement from
microwatts to megawatts, with each result gated by an accuracy target
{[}260{]}. That gate is the form of \(J^*\). FrugalGPT cascades language
models and stops when a scorer accepts the answer {[}321{]}. RouteLLM
learns to send each query to a strong or a weak model under a cost
target {[}322{]}. Their cost is dollars per query and their gate is a
quality threshold. Read with metered joules and \(C(z)\) as the gate,
both are cascades in the sense of Proposition 9. ML.ENERGY measures
inference energy automatically {[}261{]}. TokenPowerBench benchmarks the
power of large language model (LLM) inference {[}262{]}. Jin, Wei and
Brooks analyze the energy of test-time compute {[}263{]}.
General-purpose models cost much more energy per task than task-specific
ones {[}264{]}. The cheapest sufficient model wins in joules. That is
Mixture of Limits, measured. Samsi and colleagues measured the energy of
large language model inference on GPUs (graphics processing units)
{[}265{]}. Dillavou and colleagues demonstrated decentralized learning
in a physical circuit, with no processor {[}266{]}. Stern and colleagues
trained self-learning circuits for power-efficient solutions and
measured the trade between power and error {[}267{]}. Saggio and
colleagues demonstrated a quantum speed-up for learning agents in a
photonic experiment {[}268{]}. Fewer interactions to learn is fewer acts
to close.

\textbf{A kernel-reported acceptor.} Section 8.11 reports a
preregistered run of an acceptor agent on a laptop, with prediction
first, comparison, refusal and binding joule budgets. Falsifiers were
fixed before the run, and several fired: R-F1 (the selectors' joules on
shared tasks are not distinguishable), R-F2 (36 of 5,400 episodes
overspent: 9 for the σ-law arm and 7 for the fixed cascade at
\(B = 1.0\), and 20 for the goal-blind baseline at \(B = 0.5\)), R-F3,
R-F5 and R-F6 (Section 8.11). \(\iota\) was read off 4,940 receipts. The
median confirmed act gave \(1.5 \times 10^5\) bits per reported joule,
about \(2 \times 10^{15}\) times the joules per bit of the floor. Under
a binding budget the σ-law selector closed more tasks than the fixed
cascade, one Mixture of Limits instance. That comparison was chosen
after the run, so it is exploratory, and joules per closed task were
equal. The preregistered second run and a second machine repeated the
gain at the loosest budget only (Section 8.11.2). Integrated \(\iota\)
ranked a goal-blind agent first on a loose budget, while \(J^*\) ranked
it last.

\textbf{Groups simulated.} Crosato and colleagues computed the
thermodynamic quantities of simulated collective motion across its
critical point {[}269{]}. Chen and Prokopenko compared thermodynamic
efficiency with purely informational utilities as explanations of
collective behavior near criticality {[}270{]}. \(\iota_G\) in
Proposition 7 is the gated, metered form of such a ratio.

\textbf{Meters.} RAPL (Running Average Power Limit) is the set of on-die
energy counters in Intel and AMD processors. Its readings were validated
against external measurement {[}271{]}. DRAM (dynamic random-access
memory) readings were validated separately {[}272{]}. Meter quality
depends on the processor generation {[}273{]}. Apple's powermetrics
manual states that its average power values ``are estimated and may be
inaccurate'' and should not be used to compare devices {[}274{]}. On
Apple silicon, powermetrics readings are reported values. Measured
joules for that machine come from an external wall meter.

\textbf{Compression scored.} The Hutter Prize scores lossless
compression of the first gigabyte of an English Wikipedia dump, with a
prize fund of 500,000 euros and limits on runtime and memory {[}275{]}.
It bounds time and memory, never joules. A joule cap would make it a
calculus-of-agency benchmark for one task class.

\textbf{Estimates.} These sources are valuable, and their joule figures
are modeled, extrapolated or computed. They sit beside measured joules,
never in their place. Estimates are never \texttt{measured\_j}.
Strubell, Ganesh and McCallum sampled power for short runs and
extrapolated to full training of natural language processing models
{[}276{]}. Patterson and colleagues computed training energy from
reported hardware, runtime and data-center data {[}277{]}. Lacoste and
colleagues built a calculator from hardware type, runtime and region
{[}278{]}. The BLOOM footprint mixes metered and modeled inputs
{[}279{]}. Carbontracker reads on-device counters and predicts full-run
totals {[}280{]}. Its predictions are estimates. Elsworth and colleagues
report a median of 0.24 Wh per Gemini Apps text prompt, counting
accelerators, hosts, idle capacity and data-center overhead {[}281{]}.
It is an operator-reported figure, filed as \texttt{reported\_j}.

\subsubsection{5.3 Positions in the
literature}\label{positions-in-the-literature-2}

Each position is placed against the definition: intelligence is the
derivative of gated agency in joules, read off a completed run.

\begin{itemize}
\tightlist
\item
  Schmidhuber treats intrinsic reward as compression progress, a first
  derivative of how well an agent compresses its history {[}282{]}. It
  is the closest precedent for intelligence as a derivative read from a
  run. Its variable is compression, not joules.
\item
  Krakauer, Mitchell and Krakauer define intelligence as an emergent
  property expressed in ever cheaper and faster uses of emergent
  capabilities to solve problems, summarized as ``less is more''
  {[}339{]}. They conclude that large language models show emergent
  capability at most, not emergent intelligence. Their definition is the
  closest in words to \(\iota = dX/dJ\). It names energy as the
  denominator and does not meter it. The calculus supplies the meter,
  the gate that separates capability from intelligence (Section 1), and
  the derivative read off the run.
\item
  Chollet defines intelligence as skill-acquisition efficiency over
  priors and experience, with the ARC (Abstraction and Reasoning Corpus)
  benchmark {[}283{]}. It is a ratio. Its denominator is information and
  data, not joules. The ARC Prize leaderboard plots cost per task
  against score as a measure of efficiency {[}284{]}. The cost is in
  dollars. A joule meter turns it into \(J^*\).
\item
  Hernández-Orallo builds universal psychometrics across species and
  machines {[}285{]}. It evaluates by tasks. This track evaluates by
  joules to close tasks. Hernández-Orallo and Dowe's anytime test adapts
  to the time available {[}286{]}. The calculus adapts to the joules
  available.
\item
  Gershman, Horvitz and Tenenbaum define intelligence as expected
  utility net of computation cost {[}287{]}. This track makes the cost
  joules.
\item
  Lieder and Griffiths treat cognition as optimal use of limited
  resources {[}288{]}. It is close to Mixture of Limits. The resource is
  left abstract.
\item
  Balasubramanian reads the brain's design through its energy limits
  {[}289{]}.
\item
  Silver, Singh, Precup and Sutton hold that reward is enough {[}290{]}.
  Reward is chosen by a designer. Joules are read off a meter. Vamplew
  and colleagues reply that scalar reward cannot express multi-objective
  goals {[}291{]}. This track keeps one variable, joules, and moves
  multiplicity into the gate \(X \ge X_{\mathrm{req}}\).
\item
  Hafez and colleagues define agency and intelligence together
  {[}292{]}. It is not joule-native.
\item
  Karagoz makes energy self-sustainment the objective {[}293{]}. It has
  no prediction gate, no kinematics and no joule ranking across systems.
\item
  Chaisson proposes energy rate density, power per unit mass, as a
  complexity metric across cosmic evolution {[}294{]}. It is a rate in
  watts per kilogram. The calculus divides confirmed bits by joules.
\end{itemize}

The scope of \(\iota\) follows from its numerator. \(X\) counts
confirmed coupling to the world, so \(\iota\) measures how cheaply an
agent changes the world as it predicted. It does not measure
understanding. Krakauer, Mitchell and Krakauer tie intelligence to
compact, communicable representations {[}339{]}. A compact
representation enters the calculus through its effect: an agent that
holds one closes later tasks in fewer joules. Understanding that never
changes an act is outside what \(\iota\) reads.

\subsubsection{5.4 Where the calculus
stands}\label{where-the-calculus-stands}

\begin{itemize}
\tightlist
\item
  \textbf{Proven:} learning, memory, prediction and control each have a
  joule price with a ceiling.
\item
  \textbf{Demonstrated:} joules per bit are measured in brains. Joules
  per task are measured in machines. Meters can be validated.
\item
  \textbf{Proven here:} \(\iota\) is a bounded density on the joule axis
  under joule-proportional crediting (Proposition 5). Sampled \(\iota\)
  converges to it (Proposition 6). Groups compose by Proposition 7.
  Holding a goal sets an optimal act time and a ceiling on \(\iota\)
  (Prediction 1).
\item
  \textbf{Demonstrated here:} \(\iota\) and \(J^*\) read off the
  receipts of a kernel-reported acceptor, with falsifiers fixed in
  advance (Section 8.11). \(J^*\) ranks agents. Integrated \(\iota\)
  does not.
\item
  \textbf{Added by this track:} intelligence read as \(dX/dJ\) off a
  gated, completed run, with \(J^*\) as the ranking. Takahashi and
  Hayashi come closest, with bits-per-joule ratios under explicit
  accounting {[}239{]}. This track adds the eligibility gate, the
  realized flow in place of capacity, the gated ranking and its proof
  (Proposition 8), and the optimal act time (Prediction 1).
\end{itemize}

\subsection{6. Global findings, by
method}\label{global-findings-by-method}

The search ran in Russian, Chinese, Japanese, German, French, Spanish,
Portuguese, Korean, Italian, Polish, Hindi and Arabic, with English
follow-ups. A second pass searched Japanese, Chinese, Russian, German,
French, Spanish and Korean again and added Persian, Turkish and Hebrew.
Its results are folded into Sections 2 to 5 and 8 by what they prove,
measure or define. Results are organized by method and class of result.
Language is metadata only. No entry is ranked by where it came from.

\subsubsection{6.1 Theorems and reviews: information has a thermodynamic
price}\label{theorems-and-reviews-information-has-a-thermodynamic-price}

\begin{itemize}
\tightlist
\item
  Poplavskii gave an early systematic treatment of the energy cost of
  acquiring and processing information {[}295{]}. Written in Russian,
  with an English translation. Layer: definition.
\item
  Stratonovich proved value-of-information theorems {[}201{]}. English
  edition of a Russian monograph. Layer: kinematics.
\item
  Ito and Sagawa applied information thermodynamics on networks to E.
  coli chemotaxis. Transfer entropy bounds robustness, and
  information-thermodynamic efficiency is high where ordinary
  thermodynamic efficiency is low {[}296{]}. Written in Japanese.
  Layers: definition and kinematics.
\item
  Sun and Quan reviewed Maxwell's demon and the physical floor on
  dissipation in information processing {[}297{]}. Written in Chinese.
  Layer: definition.
\item
  Quan, Dong and Sun reviewed mesoscopic thermodynamics: the demon is
  consistent with the second law once erasure is counted, and
  power-efficiency constraints were tested {[}298{]}. Written in Chinese
  with an English abstract. Layers: definition and kinematics.
\item
  Parrondo's encyclopedic review gives the current account of
  information as a thermodynamic resource {[}299{]}. Layer: definition.
\item
  Strasberg's doctoral thesis tests common assumptions in information
  thermodynamics with physical models {[}300{]}. Written in English with
  a German abstract. Layer: definition.
\item
  Further theorems surfaced through searches in Japanese, Russian,
  Chinese, Korean, Portuguese, Italian and Polish: speed limits and TURs
  in chemical networks {[}127, 131, 132{]}, finite-time thermodynamics
  {[}159, 162, 163{]}, Langevin TURs {[}143, 144{]}, information bounds
  on entropy production {[}148, 149{]}, and single-shot work costs
  {[}38{]}.
\end{itemize}

\subsubsection{6.2 Experiments: the price
measured}\label{experiments-the-price-measured}

\begin{itemize}
\tightlist
\item
  Dago's doctoral thesis reports 1-bit erasure and writing near the
  minimal energy at high speed on an underdamped micro-cantilever
  {[}301{]}. Written in French. Layer: the definition's joule floor
  reached at kinematic speed.
\item
  Lagoin's doctoral thesis builds macroscopic Maxwell's demons,
  including a Szilard engine, from a vane in a granular gas {[}302{]}.
  Written in French. Layer: definition. The demon's accounting holds
  outside the microscopic regime.
\item
  Ciampini and colleagues used extractable work to witness quantum
  correlations in a photonic experiment {[}303{]}. Layer: definition.
\item
  Experiments surfaced through searches in Korean, Portuguese, Hindi,
  Italian, Spanish, Chinese and Japanese: information engines {[}27,
  80{]}, NMR demons {[}73, 74{]}, the TUR on qubits {[}202{]},
  active-bath engines {[}217, 218{]}, geometric bounds {[}211{]},
  symmetry breaking and Carnot cycles {[}215, 216{]}, trapped-ion demons
  {[}66, 88{]}, and feedback conversion and optimal transport {[}69, 76,
  209{]}.
\end{itemize}

\subsubsection{6.3 Theory of functional systems: prediction before the
act}\label{theory-of-functional-systems-prediction-before-the-act}

The Anokhin school states clauses 1 and 3 in physiological terms. A
result template, the acceptor of action results, is formed before the
act, compared after it, and the comparison reorganizes behavior. These
works were found through Russian-language searches.

\begin{itemize}
\tightlist
\item
  Anokhin's functional system {[}93, 94{]}.
\item
  Sudakov on the result as the system-forming factor and on
  probabilistic prediction of behavior {[}95, 304, 305, 306{]}.
\item
  Shvyrkov on goals organizing neuronal activity and on learning as
  selection of neuronal systems {[}96, 307{]}.
\item
  K. V. Anokhin on the brain's cognitive structure as a hypernetwork of
  functional-system elements {[}308{]}.
\item
  Saltykov and Grachev on the system-forming factor as anticipation
  {[}309{]}.
\item
  Vityaev on purposefulness formalized as rule learning that predicts
  results {[}310{]}.
\end{itemize}

\textbf{Implemented and run.} Cognitive architectures built on
functional-system theory form result templates, compare, and learn
subgoals, in simulation and on a physical robot {[}90, 91, 311, 312,
313{]}. None meters joules. Adding a meter to one of these agents would
test the full order of definition, kinematics and calculus on an
existing implementation. Differentiable probabilistic logic networks are
related calculus tooling, not built on functional systems, with no
eligibility gate and no joules {[}314{]}.

\subsubsection{6.4 Autonomy, organization and minimal
agency}\label{autonomy-organization-and-minimal-agency}

Found through Spanish, French and German searches. Agency is a
self-maintaining organization that regulates its interactions {[}103,
104, 105, 106, 107, 108{]}. The organization is a closure of constraints
on thermodynamic flows {[}109{]}. The autonomy school has stepped into
stochastic thermodynamics {[}110{]}. Layer: definition, as the
organization that must exist before any act counts.

\subsubsection{6.5 Measured joules for machine
intelligence}\label{measured-joules-for-machine-intelligence}

Found through Chinese-language searches on inference energy: automated
inference energy measurement {[}261{]}, per-token power benchmarks
{[}262{]}, and the energy of test-time reasoning {[}263{]}.

\subsubsection{6.6 Philosophy of agency}\label{philosophy-of-agency}

Rovelli's physics of agency surfaced through German {[}115{]}. Jaeger
and Azadi complete the set of position papers {[}116, 117{]}.

\subsubsection{6.7 Searches with no primary
research}\label{searches-with-no-primary-research}

Hindi and Arabic searches returned no primary research in those
languages. English follow-ups surfaced further work, listed above by
method. Persian, Turkish and Hebrew searches returned theses, reviews
and work already cited, with no new primary research.

\subsubsection{6.8 Convergence observations: information, energy and
prediction}\label{convergence-observations-information-energy-and-prediction}

Nine recent sources relate information, energy and prediction. Four were
read in full and are cited for what they state: Takahashi and Hayashi
{[}239{]}, Perrier {[}344{]}, Xu and Li {[}345{]} and Still, Sivak, Bell
and Crooks {[}54{]}. Five were read in part and are cited only for the
passages named: Wolpert {[}230{]}, Ortega and Braun {[}56{]}, Zhang,
Zhu, Zhang and Hu {[}348{]}, Fagan {[}346{]} and Kumar {[}347{]}. This
section records what each states and what it does not show. It is an
observation of overlapping questions, not a comparison of results and
not a claim about priority. Section and equation numbers are those of
the versions named. Quotations were checked by script against the saved
texts. Values marked ``recomputed'' are this paper's own arithmetic.

\textbf{In plain words}

\begin{enumerate}
\def\labelenumi{\arabic{enumi}.}
\tightlist
\item
  Several authors write the same kind of quantity: information gained or
  acted on, divided by energy spent (bits per joule), or its inverse
  (watts per unit of intelligence).
\item
  The firm results are lower bounds on heat. They say how little energy
  is physically required, not how much any real machine uses.
\item
  No source read here measured energy on real hardware. Where energy
  figures appear, they are estimated from FLOP counts or from model
  assumptions.
\end{enumerate}

\textbf{Units.} \(k_B = 1.380649\times10^{-23}\) J/K. At \(T = 300\) K,
\(k_B T\ln 2 = 2.871\times10^{-21}\) J per bit (recomputed), and its
inverse is \(3.483\times10^{20}\) bit/J (recomputed). With
\(\beta = 1/(k_B T)\), information in nats converts to joules by
multiplying by \(k_B T\), and in bits by \(k_B T\ln 2\).

\paragraph{Sources read in full}\label{sources-read-in-full}

\textbf{Takahashi and Hayashi, arXiv:2602.05463v2 (20 July 2026)
{[}239{]}.} The version read is the arXiv posting.

\begin{itemize}
\tightlist
\item
  Epiplexity per joule is \(\eta_E = \Delta I / E_{\mathrm{cons}}\) in
  bit/J, with
  \(\Delta I = I(W_{\mathrm{post}}; Z \mid W_{\mathrm{pre}})\) (Eqs 3
  and 10, Sec 2.2). The energy balance is
  \(E_{\mathrm{cons}} = Q_{\mathrm{diss}} + \Delta U_{\mathrm{sys}} + W_{\mathrm{out}} + \Delta E_{\mathrm{store}}\)
  (Eq 1).
\item
  Empowerment per joule is
  \(\eta_C^* = \sup_{p(a)} I(A_{0:\tau-1}; O_\tau) / \mathbb{E}[c(A)]\)
  in bit/J (Eq 18, Sec 3.2), with a cost-constrained curve (Eq 19).
  Empowerment is a channel capacity from actions to later observations
  (Eq 17, Sec 3.1).
\item
  The closed-cycle bound is
  \(Q_{\mathrm{diss}} \ge k_B T\ln 2\,\Delta I - T\,\Delta S_{\mathrm{sys}}\)
  (Eq 15), so \(\Delta I / Q_{\mathrm{diss}} \le 1/(k_B T\ln 2)\) when
  \(\Delta S_{\mathrm{sys}} = 0\) (Eq 16, Corollary 1). Eq 15 follows
  from Eq 14 by rearrangement (checked). Corollary 1 needs the
  bipartite-Markov and local-detailed-balance assumptions.
\item
  The authors give a bits-per-joule metric for learning, another for
  control, and a reporting checklist: boundary, baseline energy,
  horizon, estimator (Sec 5.1). Their Proposition 1 is a decoupling
  construction showing the bound fails if the accounting boundary
  excludes the cost of preparing fresh memory.
\item
  Scope as stated: ``the resulting numbers are benchmark-relative rather
  than universal''. On the ceiling: ``Landauer-scale limits are
  conceptual yardsticks far beyond current hardware.''
\item
  The Sec 5 worked example uses published perplexities, \(10^9\) tokens,
  and energy estimated from \(6PD\) FLOPs at an assumed
  \(2.5\times10^{11}\) FLOP/J. Recomputed:
  \(\Delta G \approx 4.51\times10^{8}\) bit, \(E\) (410M model)
  \(\approx 2.95\times10^{9}\) J, marginal efficiency \(\approx 0.106\)
  bit/J. The authors write: ``This number is convention-dependent, but
  it exercises the checklist.''
\item
  Not verified here: Lemma 1 restates an inequality from work not read.
  Eq 17 uses \(O_\tau\) while the Sec 3.2 text refers to
  \(I(A;S_\tau)\). Eq 20 is flagged by the authors: ``This is not
  claimed as a universal identity without assumptions''.
\end{itemize}

\textbf{Perrier, arXiv:2504.05328v2 (16 June 2025) {[}344{]}.}
Definitions only are cited.

\begin{itemize}
\tightlist
\item
  Watts per intelligence is
  \(\Phi(\Sigma) = P(\Sigma)/I(\Sigma) = E(\Sigma)/(\tau\, I(\Sigma))\)
  (Eq 7, Sec 3.3), where \(I\) is a task-weighted performance sum (Eq 2)
  or the Legg and Hutter sum (Eq 1). The assumptions are \(E = F N c\)
  with \(c = k_B T\ln 2\) and overhead \(F \ge 1\) (Eq 5), and
  \(I \le \alpha N\) with \(\alpha\) unspecified (Eq 6). Theorem 1 gives
  \(\Phi \ge cF/(\alpha\tau)\) (Eq 8).
\item
  The author states the premise openly: ``Our work assumes a correlation
  between irreversible computation to intelligence.'' The paper contains
  no measurement; it says only that ``In practice, one might measure the
  actual hardware usage''.
\item
  Errors found (my reading, not confirmed with the author): (1) Sec 3.1
  defines \(N\) per unit time, but Eqs 3 and 5 use \(E\) over
  \([0,\tau]\), which needs \(N\) as a total, and then
  \(cF/(\alpha\tau)\) falls to zero as \(\tau\) grows. (2) \(\alpha\) is
  free, so Theorem 1 gives no numerical floor. (3) The Theorem 2 proof
  yields a bound of the form \(\Delta i_K > -\log(1/\delta)\), which is
  negative for \(\delta < 1\), yet the text treats \(-\log(1/\delta)\)
  as positive. (4) Combining \(I \le \alpha\,\Delta i_K\) with
  \(P \ge \Delta i_K\, k_B T\ln 2/\tau\) gives
  \(I/P \le \alpha\tau/(k_B T\ln 2)\), while Eq 14 has no \(\alpha\), no
  \(k_B T\ln 2\) and the opposite dependence on \(\tau\). (5)
  \(\Delta i_K = K(y) - K(x)\) can be negative, and the direction of the
  Landauer attachment in Eq 13 is not argued; Theorem 3 inherits this.
  (6) The hardware ordering
  \(\Phi_{\mathrm{CPU}} > \Phi_{\mathrm{GPU}} > \Phi_{\mathrm{neuro}}\)
  is asserted with no values.
\end{itemize}

\textbf{Xu and Li, arXiv:2511.19156v4 (30 December 2025) {[}345{]}.}
Read in full for the saved text; figures were read by caption and
printed values only. Definitions only are cited.

\begin{itemize}
\tightlist
\item
  Derivation depth \(D_d\) (Eq 10) and derivation entropy
  \(H_{\mathrm{derive}} = D_d\lambda\) with \(\lambda = \ln 2\) nats in
  a stated classical regime (Eq 14, Assumption 3). Claimed bounds:
  \(E_{\mathrm{compute}} \ge D_d k_B T\ln 2\) (Proposition 4, Eqs 98 and
  99) and
  \(E_{\mathrm{storage}} \ge |S|\,k_B T\ln 2\,(t/t_{\mathrm{refresh}})\)
  (Proposition 5, Eq 100). A break-even query frequency between
  recomputing and storing is given in abstract cost units (Theorem 5, Eq
  32).
\item
  The authors call their cost an ``information-theoretic proxy for
  energy consumption''. The numerical validation (Sec VI) uses a
  synthetic 100K-atom knowledge base and the paper's own abstract costs,
  so it tests the model against itself. Application figures (1000 kWh,
  45 percent reduction) come without a measurement method and are not
  used here.
\item
  The abstract says ``Our findings suggest that the minimization of
  Derivation Entropy is a governing principle''. The text elsewhere
  declines to claim a sharp result: ``We deliberately refrain from
  formalizing'' the last observation as a phase-transition theorem.
\item
  Errors found (my reading, not confirmed with the authors): (1)
  Assumption 1 has \(E\tau \ge \kappa\) with \(\kappa\) an action in J
  s, while Corollary 1 sets \(\kappa_T = k_B T\ln 2\), an energy, and
  drops \(\tau\). (2) Eq 98 assumes every inference step erases a bit
  (Assumption 3), and Remark 15 of the same paper concedes reversible
  computation avoids this cost. (3) In Lemma 2 the contradiction
  argument negates the bound to \(K = o(n\log m)\) and claims this
  contradicts \(H \gg \log m\); it does not, since
  \(n\log m \gg \log m\). The lower bound is therefore unproved, and
  Theorems 4 to 6 and Corollary 7 depend on it. (4) Theorem 6 case (b)
  uses \(N = \Omega(L_S)\) where \(N = O(L_S)\) is needed. (5) The Eq
  135 critical storage with \(H(Q) = 10^{12}\) recomputes to
  \(1.25\times10^{10}\) (log base 2), \(1.80\times10^{10}\) (ln) or
  \(4.15\times10^{10}\) (log base 10), not the printed
  \(2.5\times10^{11}\). The Table II text (tightness 0.50 to 0.93)
  disagrees with the table (0.50 to 0.81, and 0.89 at high frequency).
\end{itemize}

\textbf{Still, Sivak, Bell and Crooks, arXiv:1203.3271v3 (5 October
2012) {[}54{]}.} The arXiv version was read; the journal version was
not.

\begin{itemize}
\tightlist
\item
  With system state \(s_t\) and driving signal \(x_t\),
  \(I_{\mathrm{mem}}(t) = I[s_t, x_t]\) and
  \(I_{\mathrm{pred}}(t) = I[s_t, x_{t+1}]\), in nats. Per step,
  \(\beta\langle W_{\mathrm{diss}}[x_t \to x_{t+1}]\rangle = I_{\mathrm{mem}}(t) - I_{\mathrm{pred}}(t)\)
  (Eq 14). In total,
  \(I_{\mathrm{mem}} - I_{\mathrm{pred}} \le \beta\langle W_{\mathrm{diss}}\rangle \le \beta\langle W_{\mathrm{ex}}\rangle\)
  (Eq 18). The refined Landauer bound is
  \(-\beta\langle Q\rangle \ge I_e + I_{\mathrm{mem}} - I_{\mathrm{pred}}\)
  (Eq 21).
\item
  Memory of the signal that does not help predict its next value costs
  dissipated work, at \(k_B T\) per nat. The authors' stated
  consequence: ``any system constructed to keep memory about its
  environment and to operate with maximal energetic efficiency has to be
  predictive''.
\item
  This is a lower bound on dissipation for systems that keep memory of a
  driving signal. It does not show that any real system attains it, and
  it is not a statement about general agents. There is no experiment and
  no joule data.
\item
  Checked: footnote 24 is correct; Eq 17 follows from Eqs 11, 14 and 15;
  Eq 18 follows from Eq 16; Eq 21 follows from Eqs 19 and 20.
  \(k_B T = 25.9\) meV at 300 K (recomputed; the paper gives 25 meV).
  Not verified: non-negativity of the per-step term for arbitrary signal
  processes is stated but not proved in the text, and it is not proved
  here.
\end{itemize}

\paragraph{Sources read in part}\label{sources-read-in-part}

Each entry is limited to the passage stated. No claim is made about the
unread parts.

\begin{itemize}
\tightlist
\item
  \textbf{Wolpert, arXiv:1905.05669v3 {[}230{]}.} Sec V (Eqs 29 to 44,
  Examples 1 to 6), Sec XVI, Appendix A and the opening of Appendix B
  were read. Sec V states \(Q(p_0) = \sigma(p_0) + S(p_0) - S(\pi p_0)\)
  (Eq 29) and the Landauer cost \(L(p,\pi) = S(p) - S(\pi p)\) (Eq 35),
  defined only for a fixed input distribution. About 70 percent of the
  115 pages was not read, and Proposition 1 rests on a cited result not
  read.
\item
  \textbf{Ortega and Braun, arXiv:1204.6481v2 {[}56{]}.} The main text
  (Sec 1 to 8) was read. The proofs of Theorems 5 to 7 and the
  Supplementary Material were not available. Cited for Eq 9,
  \(p(x) \propto p_0(x)\, e^{\beta U(x)}\), and for the one-molecule
  illustration of Sec 1.1 (cost \(k_B T\ln(1/p)\)). Printed errors were
  found in Eq 5 and Corollary 3. The stated optimum \(M = 35\) in Fig 4
  recomputes to 33 under my reading of the setup.
\item
  \textbf{Zhang, Zhu, Zhang and Hu, arXiv:2607.24287v1 {[}348{]}.} Sec
  4.1 to 4.6 (proofs of T1 to T6) were read. Cited for the statement
  that Axiom 1 ``is adopted as an empirical postulate rather than
  derived within the present framework'' and that its energy is
  dimensionless (nats). The proofs of T7 to T10 and Sec 4.7 and 4.8 were
  not read. One discrepancy in the T5 Euler-Lagrange equation is
  unresolved.
\item
  \textbf{Fagan, arXiv:2601.00021v2 {[}346{]}.} Sec 5.3 to 14 and the
  references were read; Sec 1 to 5.2 were read earlier and not re-read.
  Cited for Eq 116,
  \(\chi = \dot W_{\mathrm{causal}} / \dot I_{\mathrm{irr}}\), and the
  authors' statement that its proxies ``need not coincide numerically
  across systems''. The simulations (Sec 12) use proxies, not joules,
  and one curve in Sec 12.1 decreases with \(\lambda\) by construction.
\item
  \textbf{Kumar, Zenodo 18993403 {[}347{]}.} Appendix A.1 to A.28 and
  the references were read; the main text was read earlier and not
  re-read. Superscripts are lost in the saved text, so printed exponents
  in A.26 and A.28 are not cited. Cited for \(I \le P/(k_B T\ln 2)\) (Eq
  125), whose premise the paper states as conditional (``If generating
  useful knowledge requires irreversible information processing'').
  Recomputed: \(6.74\times10^{21}\) bit/s for 20 W at 310 K.
\end{itemize}

\paragraph{What the nine sources show
together}\label{what-the-nine-sources-show-together}

{\def\LTcaptype{none} % do not increment counter
\begin{longtable}[]{@{}
  >{\raggedright\arraybackslash}p{(\linewidth - 6\tabcolsep) * \real{0.2500}}
  >{\raggedright\arraybackslash}p{(\linewidth - 6\tabcolsep) * \real{0.2500}}
  >{\raggedright\arraybackslash}p{(\linewidth - 6\tabcolsep) * \real{0.2500}}
  >{\raggedright\arraybackslash}p{(\linewidth - 6\tabcolsep) * \real{0.2500}}@{}}
\toprule\noalign{}
\begin{minipage}[b]{\linewidth}\raggedright
Source
\end{minipage} & \begin{minipage}[b]{\linewidth}\raggedright
What it shows
\end{minipage} & \begin{minipage}[b]{\linewidth}\raggedright
What it does not show
\end{minipage} & \begin{minipage}[b]{\linewidth}\raggedright
Read status
\end{minipage} \\
\midrule\noalign{}
\endhead
\bottomrule\noalign{}
\endlastfoot
Takahashi and Hayashi & Two bit/J metrics (learning, control); a
closed-cycle Landauer-scale ceiling under stated assumptions; a
reporting checklist & A measured joule; an estimate of epiplexity or
empowerment for a real system; a test of the bound & Full (16 pp, no
appendices) \\
Perrier & A definition \(\Phi = P/I\) and a bound that holds only given
an assumed link \(I \le \alpha N\) & Measured energy; a numerical floor;
a correct Theorem 2 or 3 proof & Full (12 pp); definitions cited only \\
Xu and Li & Definitions of derivation depth and entropy; conditional
Landauer-type bounds; a synthetic simulation & Measured energy; a
correct proof of Lemma 2; independent validation & Full for the saved
text; definitions cited only \\
Still, Sivak, Bell and Crooks & A dissipation lower bound equal to
non-predictive memory,
\(\beta\langle W_{\mathrm{diss}}\rangle \ge I_{\mathrm{mem}} - I_{\mathrm{pred}}\)
& Attainment by real systems; general agents; a measured joule & Full (5
pp); journal version not read \\
\end{longtable}
}

Across all nine sources, no text read reports joules measured on real
hardware. Energy figures that appear are estimated from FLOP counts,
nominal power ratings or model assumptions. Where bounds appear, they
are lower bounds on dissipation under stated assumptions, or definitions
of ratios awaiting measurement.

\textbf{Limits of this section.} Quoted text was verified against
PDF-converted text, and two-column conversion can interleave lines. The
errors listed are from my reading and recomputation and have not been
sent to the authors. Version differences (v1 to v2) were compared by
structure and abstract only for Takahashi and Hayashi, Perrier and Xu
and Li.

\subsection{7. Open problems}\label{open-problems}

\subsubsection{7.1 What must be proven}\label{what-must-be-proven}

\textbf{P1. The run path is well defined.} Proven. Proposition 5 gives
\(\iota = dX/dJ\) as a density on the joule axis with
\(0 \le \iota \le 1/(k_B T\ln 2)\), under joule-proportional crediting.
Proposition 6 gives convergence of sampled \(\iota\) as the meter's
sampling refines, with the meter's relative error as the error term.

\textbf{P2. Agency-gated Landauer bound.} Under closed-cycle accounting,
each confirmed coupled bit costs at least \(k_B T \ln 2\), so
\(v_J \le 1/(k_B T \ln 2)\). In hand: measurement plus erasure work is
at least \(k_B T I\) {[}5{]}, and erasure costs at least \(k_B T \ln 2\)
per bit {[}6, 7{]}. Missing: a data-processing step showing that the
coupled bits credited in \(X\) cannot exceed the bits recorded in the
acceptor's memory during the comparison. Clause 3 already requires the
confirmation \(m\) to come from a physical record. Proposition 1 proves
the per-act floor with this step taken as assumption (R2).

\textbf{P3. Agency speed limit.} Proven for acts whose record meets (R1)
to (R4). Proposition 1 bounds the mean joules of act \(k\) below by
\(b_k k_B T\ln 2\) plus a finite-time term in \(\tau_k\),
\(\langle A\rangle_k\), the reset error \(\varepsilon\) and the
confirmed bits, through Fano's inequality. Prediction 1 adds the
standing joules of holding a goal and gives the optimal act time.

\textbf{P4. Precision bound for agents.} Proven. Proposition 3 bounds
the relative variance of confirmed agency by dynamical activity for the
count of confirmed acts, and by entropy production with the information
term under feedback.

\textbf{P5. Hindsight theorem.} \(\iota\) is a function of the completed
receipt only. No pre-run quantity fixes it. In hand: each \(m_k\)
compares an outcome to a prediction fixed before the act, so \(m_k\) is
unknown before act \(k\) closes. Missing: a construction of two worlds
that agree on everything a pre-run predictor sees and differ in some
\(m_k\). With P1, this makes ``intelligence is read off the completed
run'' a theorem.

\textbf{P6. Ranking survives meter error.} Proven. Proposition 8, part
3: if \((1+\epsilon)J^*_{S_1}(q) < (1-\epsilon)J^*_{S_2}(q)\), with
\(\epsilon\) the meter's relative error bound, then \(S_1 \succ_q S_2\)
under any reading within that error. Missing: validated \(\epsilon\) for
each meter {[}271{]}.

\textbf{P7. Mixture of Limits optimality.} Partly proven. Proposition 9
gives the exact conditions for independent gears and for two nested
gears, and the budget mechanism. The fixed cascade instance tries gears
from cheapest up, stops at the first gear whose output passes \(C(z)\),
and refuses with a receipt if the next gear would break the budget. It
achieves the joule minimum up to the cost of failed attempts. In hand:
if sufficiency is monotone in gear order and gear costs rise with order,
the first sufficient gear is the cheapest sufficient gear. Missing: a
regret bound when pass probabilities are estimated or sufficiency is not
monotone, the nested case with more than two gears, and the conditions
under which this rule and the σ-law selector,
\(a^* = \arg\max[H(a) - \lambda J(a)]\) subject to \(J \le B\), choose
the same act.

\textbf{P8. When agency compounds.} Conditions under which
\(\alpha_J > 0\) over a run, so that earlier receipts lower the joules
of later closures. Missing: a model of reuse in which stored receipts
cut later \(J\) by more than the storage and lookup joules. Still gives
the memory cost side {[}225{]}. Section 8.10 measures the template's
share of each act's joules.

\subsubsection{7.2 What must be
demonstrated}\label{what-must-be-demonstrated}

{\def\LTcaptype{none} % do not increment counter
\begin{longtable}[]{@{}
  >{\raggedright\arraybackslash}p{(\linewidth - 4\tabcolsep) * \real{0.3333}}
  >{\raggedright\arraybackslash}p{(\linewidth - 4\tabcolsep) * \real{0.3333}}
  >{\raggedright\arraybackslash}p{(\linewidth - 4\tabcolsep) * \real{0.3333}}@{}}
\toprule\noalign{}
\begin{minipage}[b]{\linewidth}\raggedright
ID
\end{minipage} & \begin{minipage}[b]{\linewidth}\raggedright
Claim
\end{minipage} & \begin{minipage}[b]{\linewidth}\raggedright
Evidence required
\end{minipage} \\
\midrule\noalign{}
\endhead
\bottomrule\noalign{}
\endlastfoot
D1 & \(J^*\) per task is measured, not estimated, on at least two
fabrics & \texttt{measured\_j} from RAPL on x86 Linux and an external
wall meter on Apple silicon \\
D2 & The σ-law selector closes tasks in fewer measured joules than the
fixed cascade under binding budgets & Paired runs with the same tasks,
gears, predicates, budgets and meters \\
D3 & Refuse with a receipt works & Zero overspends, refusals on
infeasible tasks, few refusals on feasible ones \\
D4 & \(\iota\) is read off receipts and is stable across seeds & Per-run
\(\iota_k\) series with spread within a preset bound \\
D5 & Estimates do not substitute for measurement & Rank agreement
between \texttt{est\_j} and \texttt{measured\_j}, reported, never
merged \\
D6 & The physical floor is respected and the gap is known & Measured
joules per confirmed bit divided by \(k_B T \ln 2\) \\
D7 & Prediction 1: the optimal act time and the ceiling of an agent that
holds a goal & A \(\tau\) sweep at two or more instability rates, with
act and standing joules metered separately and together \\
\end{longtable}
}

Section 8.11 reports the first run against D2 to D6, on reported joules.
Section 8.11.1 reports the second run on the same machine, with the
guard and refusal fixes, and Section 8.11.2 repeats it on a second
machine.

\subsection{8. Experiment plan}\label{experiment-plan}

Every table, receipt and figure in this plan follows one rule. Estimates
are never \texttt{measured\_j}. A modeled figure is labeled
\texttt{est\_j}. A metered figure is labeled \texttt{measured\_j} and
names its meter.

\subsubsection{8.1 Meters}\label{meters}

Acronyms: CPU (central processing unit), GPU (graphics processing unit),
ANE (Apple Neural Engine), AC (alternating current), NVML (NVIDIA
Management Library).

{\def\LTcaptype{none} % do not increment counter
\begin{longtable}[]{@{}
  >{\raggedright\arraybackslash}p{(\linewidth - 6\tabcolsep) * \real{0.2500}}
  >{\raggedright\arraybackslash}p{(\linewidth - 6\tabcolsep) * \real{0.2500}}
  >{\raggedright\arraybackslash}p{(\linewidth - 6\tabcolsep) * \real{0.2500}}
  >{\raggedright\arraybackslash}p{(\linewidth - 6\tabcolsep) * \real{0.2500}}@{}}
\toprule\noalign{}
\begin{minipage}[b]{\linewidth}\raggedright
Tier
\end{minipage} & \begin{minipage}[b]{\linewidth}\raggedright
Meter
\end{minipage} & \begin{minipage}[b]{\linewidth}\raggedright
What it reads
\end{minipage} & \begin{minipage}[b]{\linewidth}\raggedright
Label
\end{minipage} \\
\midrule\noalign{}
\endhead
\bottomrule\noalign{}
\endlastfoot
M0 & External wall-plug power meter, logging at 1 Hz or faster &
Whole-system AC power & \texttt{measured\_j}, wall \\
M1 & RAPL through Linux powercap, package, core, uncore and DRAM domains
where present & Cumulative on-die energy in microjoules &
\texttt{measured\_j}, rapl \\
M1g & NVML total energy counter on a discrete GPU & Cumulative GPU board
energy & \texttt{measured\_j}, nvml \\
M2 & Apple powermetrics, CPU and GPU samplers at 100 ms & CPU, GPU and
ANE power as reported by the system & \texttt{reported\_j},
powermetrics \\
\end{longtable}
}

\begin{itemize}
\tightlist
\item
  \textbf{RAPL.} Read the counter before and after each act. Handle
  wraparound. Validate against M0 on a calibration workload {[}271,
  272{]}.
\item
  \textbf{powermetrics.} Readings are filed as \texttt{reported\_j}.
  Measured joules on Apple silicon come from M0. Within one machine,
  \texttt{reported\_j} attributes joules to CPU, GPU and ANE.
\item
  \textbf{Idle baseline.} Record 120 s of idle power on every meter
  before each block. Report gross and idle-subtracted joules side by
  side. Never report only the subtracted figure.
\item
  \textbf{State control.} Fix the CPU frequency governor where the
  operating system allows. Log ambient temperature and thermal pressure.
  Close background jobs. Discard one warm-up run per arm.
\end{itemize}

\subsubsection{8.2 Tasks}\label{tasks}

Each task has a machine-checkable completeness predicate \(C(z)\): unit
tests for code, exact match for arithmetic, schema validation for
extraction, a verified final state for tool use. Two strata are fixed in
advance. Near-optimal tasks are those the cheapest gear usually closes.
Compounding tasks are sequences where earlier receipts can be reused.
\(X_{\mathrm{req}}\) is set by \(C(z)\). Passing is a gate, not a score.

\subsubsection{8.3 Arms}\label{arms}

All arms share the gears, the predicate, the budget, the meters and a
randomized interleaved run order.

\begin{enumerate}
\def\labelenumi{\arabic{enumi}.}
\tightlist
\item
  \textbf{σ-law selector:} \(a^* = \arg\max[H(a) - \lambda J(a)]\)
  subject to \(J \le B\), with refusal on a receipt.
\item
  \textbf{Fixed cascade (cheapest first):} the cheapest-sufficient rule,
  with gears tried in a fixed order from cheapest up. It is the matched
  comparison: same information, gears, budget and refusal right. Only
  the selection rule differs.
\item
  \textbf{Largest gear always:} reference ceiling.
\item
  \textbf{Smallest gear always:} reference floor, with refusal when it
  fails.
\end{enumerate}

Both selectors are Mixture of Limits instances. The philosophy is to use
the cheapest gear that is sufficient, because past the floor additional
joules do not purchase verifiable progress. The fixed cascade applies it
in a fixed order. The σ-law selector stops when the next act no longer
pays for its joules. That is a threshold-convergence rule, so it is a
Mixture of Limits instance too. The comparison is between two instances
of one philosophy. In code and data files the fixed cascade is the arm
\texttt{mol} and the σ-law selector is the arm \texttt{sigma}.

\subsubsection{8.4 Budgets, seeds and
statistics}\label{budgets-seeds-and-statistics}

\begin{itemize}
\tightlist
\item
  A pilot runs arm 2 without a budget and records the median measured
  joules per task, \(\tilde J_q\).
\item
  Budget levels are
  \(B \in \{0.25, 0.5, 0.75, 1.0\} \times \tilde J_q\), plus
  unconstrained. A level binds only if at least 20\% of tasks exceed it
  in the pilot. Levels that do not bind are reported and excluded from
  the D2 test.
\item
  At least 10 seeds per arm, task and budget. More if the pilot shows 10
  cannot resolve a 10\% difference in median \(J^*\) at 95\% confidence.
\item
  Report mean, median, standard deviation and a 95\% bootstrap
  confidence interval of \(J^*\) per task. Compare arms 1 and 2 with a
  Wilcoxon signed-rank test on per-task median \(J^*\), with effect
  size.
\item
  The preregistration names the primary endpoints. Secondary comparisons
  are corrected with the Holm step-down procedure. An analysis chosen
  after a run is labeled post hoc wherever it appears, in summaries
  included, and it is never the headline result. It becomes a primary
  endpoint only in the preregistration of the next run.
\end{itemize}

\subsubsection{8.5 Receipt schema}\label{receipt-schema}

The receipt is the record of the run. One record per act:

\begin{verbatim}
act_id, run_id, task_id, arm, seed, budget_j, gear
prediction_hash, prediction_ts        # R*, fixed before the act
outcome_hash, outcome_ts              # R
m                                     # 1, mu, or 0
coupled_bits_est                      # realized I(do(a); W'), bias-corrected, with interval
coupled_bits_bound                    # capacity bound, never credited as X
t_start, t_end
measured_j   {meter, domains, gross_j, idle_w, idle_subtracted_j}
reported_j   {meter, cpu_j, gpu_j, ane_j}    # powermetrics only
wall_j       {meter, gross_j, idle_w, idle_subtracted_j}
est_j        {model, inputs}                 # never merged with measured_j
cumulative_measured_j
refused, refuse_reason
passed_C
record_states, record_bits            # N and H of the act's record M_k
joule_share                           # P_k attribution when acts overlap
template_j   {meter, write_j, hold_j, reset_j}   # joules for R*
trajectory_j {tau, activity, work}    # physical runs only, Section 8.8
\end{verbatim}

\subsubsection{8.6 Metrics}\label{metrics}

\begin{itemize}
\tightlist
\item
  \(J^*\) per task: measured joules to pass \(C(z)\), from M0 where
  present, otherwise M1.
\item
  Closure rate per arm and budget.
\item
  Overspends: acts that exceeded \(B\). The target is zero.
\item
  Refusal quality: precision and recall of refusals against tasks shown
  infeasible within \(B\) by any arm.
\item
  Intelligence series: \(\iota_k = \Delta X_k/\Delta J_k\) per run, and
  its spread across seeds.
\item
  Acceleration sign: the sign of \(\alpha_J\) from regressing
  \(\iota_k\) on cumulative \(J\), per stratum.
\item
  Estimate check: Kendall's rank correlation between \texttt{est\_j} and
  \texttt{measured\_j}.
\item
  Distance to floor: measured joules per confirmed coupled bit divided
  by \(k_B T \ln 2\) at the logged temperature.
\item
  Act intelligence: \(\iota_k = b_k/J_k\) per act, which equals the
  density of Proposition 5 on that act.
\item
  Template share: \(J^{R^*}/J\) per act and the curve \(\iota(w)\)
  (Section 8.10).
\end{itemize}

\subsubsection{8.7 Falsifiers, fixed before any
run}\label{falsifiers-fixed-before-any-run}

{\def\LTcaptype{none} % do not increment counter
\begin{longtable}[]{@{}
  >{\raggedright\arraybackslash}p{(\linewidth - 4\tabcolsep) * \real{0.3333}}
  >{\raggedright\arraybackslash}p{(\linewidth - 4\tabcolsep) * \real{0.3333}}
  >{\raggedright\arraybackslash}p{(\linewidth - 4\tabcolsep) * \real{0.3333}}@{}}
\toprule\noalign{}
\begin{minipage}[b]{\linewidth}\raggedright
ID
\end{minipage} & \begin{minipage}[b]{\linewidth}\raggedright
Observation
\end{minipage} & \begin{minipage}[b]{\linewidth}\raggedright
Claim defeated
\end{minipage} \\
\midrule\noalign{}
\endhead
\bottomrule\noalign{}
\endlastfoot
F1 & The 95\% interval for the paired difference in median \(J^*\),
σ-law minus fixed cascade, includes or exceeds zero at every binding
budget & The σ-law selector beats its fair rival in joules \\
F2 & Any arm records an overspend & That arm's refusal implementation \\
F3 & Refusal recall on infeasible tasks is below 0.9, or more than 10\%
of refusals fall on tasks another arm closed within \(B\) & Refusal
correctness \\
F4 & Kendall's correlation between \texttt{est\_j} and
\texttt{measured\_j} is below 0.8 & Use of that estimator for any
ranking \\
F5 & The coefficient of variation of mean \(\iota\) across seeds exceeds
0.25 on a task & Reading \(\iota\) at that task's granularity \\
F6 & The fitted \(\alpha_J\) is not positive with 95\% confidence on the
compounding stratum & Compounding on those tasks \\
F7 & The mean measured joules per confirmed bit over a run fall below
\(k_B T \ln 2\), or any act of a macroscopic machine does & The
accounting itself: a meter or bookkeeping fault, fixed before any other
result is reported \\
F8 & M1 and M0 disagree on a calibration workload by more than the error
bound in P6 & That RAPL domain as a measured source on that machine \\
F9 & A physical act cycle shows work below the Proposition 1 floor at
its logged \(\tau\) and \(\langle A\rangle\) & The accounting of that
run: a missed jump channel, a miscalibrated potential or an unmetered
work source \\
F10 & A stationary run shows \(\mathrm{Var}(N)/\langle N\rangle^2\)
below \(1/\langle K\rangle\) beyond its confidence interval & The jump
count of that run: some jumps went uncounted \\
F11 & A controller holds the plant with measured joules per step below
\(k_B T\ln 2 \log_2\lvert\lambda\rvert\) & The meter or the plant model
of that run \\
F12 & The maximizer of measured \(\iota(\tau)\) lies outside
\(\tau^* \pm 2\delta\tau^*\) of Prediction 1, or the slope of
\(\ln\tau^*\) against \(\ln s\) lies outside \(-1/2\) by more than two
standard errors; or a run's \(\iota\) exceeds the ceiling of Prediction
1 & Additivity of act and standing joules on one meter, with \(\iota\)
as their density; for the ceiling, the accounting of that run \\
\end{longtable}
}

\subsubsection{8.8 Physical act cycles}\label{physical-act-cycles}

These runs test Propositions 1 to 4 at the physical floor. Work is
computed from measured trajectories and the calibrated potential, as in
{[}61, 67, 70{]}. It is filed as \texttt{trajectory\_j} and never merged
with \texttt{measured\_j}, \texttt{reported\_j} or \texttt{est\_j}.

\begin{itemize}
\tightlist
\item
  \textbf{Act cycle on a charge register (Propositions 1 and 4).} A
  single-electron box or quantum dot holds a one-bit record, as in
  {[}68, 70, 71{]}. Each cycle writes the comparison, acts, and resets
  the record in time \(\tau\). A charge detector counts every tunneling
  jump, so \(\langle A\rangle\) is measured. Sweep \(\tau\) over at
  least two decades. Result form: work per cycle against
  \(k_B T\ln 2 + k_B T/(2\tau\langle A\rangle)\), the floor of
  Proposition 1 for \(N = 2\) and \(b = 1\).
\item
  \textbf{Counting confirmed acts (Proposition 3).} On the same
  register, cycling in steady state, count confirmed acts \(N\) and all
  jumps \(K\) over repeated windows of equal length. Result form:
  \(\mathrm{Var}(N)/\langle N\rangle^2\) against \(1/\langle K\rangle\),
  and the spread of \(T_n\) against
  \(1/(\langle T_n\rangle\langle k\rangle)\). On a feedback-trap
  information engine as in {[}27{]}, run the forward and the backward
  experiment, compute \(\sigma_I\) from the measured trajectories, and
  test part 3.
\item
  \textbf{The optimal act time (Prediction 1).} This is the lab test.

  \begin{itemize}
  \tightlist
  \item
    \emph{Setup.} A colloidal bead in water at temperature \(T\) sits in
    a feedback trap {[}61{]}. Its position is read every \(\Delta t\),
    and a feedback force imposes a virtual potential \(U(x, t)\). Work
    is computed from the trajectory by stochastic energetics,
    \(W = \sum_t \partial_t U(x_t, t)\,\Delta t\) {[}319{]}. One
    coordinate runs the one-bit act cycle in a virtual double well,
    reset in time \(\tau\). The other is held at the top of a virtual
    inverted well, \(U = -\tfrac{1}{2}k x^2\), whose instability rate
    \(\Lambda\) is set by \(k\) and the bead's mobility.
  \item
    \emph{Blocks.} For each value of \(\Lambda\), three blocks run
    interleaved in randomized order. Block A: acts with no goal held, at
    nine or more log-spaced values of \(\tau\) over at least two
    decades; it fits \(J_0\) and \(a\). Block B: the goal held with no
    acts; it gives the standing power \(s\). Block C: both together, at
    the same values of \(\tau\); it gives \(\iota(\tau)\). At least 200
    cycles per point.
  \item
    \emph{Sweep of \(s\).} At least four values of \(\Lambda\) spanning
    at least one decade, so \(s\) spans at least one decade.
  \item
    \emph{Fit.} Fit \(J(\tau) = J_0 + a/\tau + s\tau\) to block C by
    weighted least squares. Read the measured \(\tau^*\) at the maximum
    of \(\iota(\tau)\) by a local quadratic fit in \(\ln\tau\). Compute
    \(\sqrt{a/s}\) and its band \(\delta\tau^*\) from blocks A and B.
    Regress \(\ln\tau^*\) on \(\ln s\) by least squares and report the
    slope with its standard error.
  \item
    \emph{Falsification.} The claim fails if the slope lies outside
    \(-1/2\) by more than two standard errors, or a measured \(\tau^*\)
    lies outside \(\sqrt{a/s} \pm 2\delta\tau^*\), or block C's joules
    per act differ from the sum of blocks A and B beyond their combined
    95\% interval. The last test is additivity on one meter, the
    physical content of the prediction.
  \item
    \emph{Power.} The slope's standard error must be below 0.1.
    Otherwise the run is reported as inconclusive and the sweep is
    extended. An inconclusive run is not a pass.
  \item
    \emph{Second substrate.} A single-electron box with a charge
    detector repeats the protocol, with \(\langle A\rangle\) counted
    jump by jump {[}70{]}. Agreement on two substrates is the test.
  \item
    All work is computed from trajectories and filed as
    \texttt{trajectory\_j}.
  \end{itemize}
\item
  \textbf{Holding an unstable plant (Proposition 2).} A digital
  controller holds a plant with one unstable mode of fixed \(|\lambda|\)
  per step, at \(f\) steps per second. Meter the controller with M0 and
  M1. Log the record's bits per step. Result form: measured joules per
  step divided by \(k_B T\ln 2\log_2|\lambda|\), reported as distance to
  floor. Sweep \(|\lambda|\) and \(f\).
\end{itemize}

\subsubsection{8.9 Living agents: measured
ι}\label{living-agents-measured-ux3b9}

Result form: \(\iota\) in a living preparation, as bits per decision
divided by joules per decision, both taken in the same preparation and
reported with distance to floor. A living \(\iota\) counts as measured
only when both terms are measured. A joule figure computed from
conductances or inferred from imaging is filed as \texttt{est\_j}.
Entropy production inferred from brain imaging is coarse-grained and
informational, so it is not a joule meter {[}257{]}.

\begin{itemize}
\tightlist
\item
  \textbf{Bacterial chemotaxis.} Information rate from single-cell
  behavior in a defined gradient, as Mattingly and colleagues measured
  it {[}213{]}. Joules from the oxygen consumption of the same culture,
  divided by cell number. The comparison floor is the cost of sending a
  bit through the cell's signaling channel, from Bryant and Machta
  {[}197{]}, beside \(k_B T\ln 2\).
\item
  \textbf{Mouse visual cortex.} Orientation information per trial from
  two-photon imaging, and synaptic ATP use from whole-cell recordings,
  as Padamsey and colleagues obtained both {[}250{]}. Food restriction
  moves the joules. The measured response of information to joules is
  the sign and size of \(dX/dJ\) in a cortex.
\item
  \textbf{Fly memory.} Energy flux in the mushroom body during long-term
  memory formation, imaged as Plaçais and colleagues did {[}252{]}, per
  learned association, with retention as the confirmation. Survival
  under starvation and desiccation is the outcome cost {[}251, 253{]}.
\end{itemize}

\subsubsection{8.10 The joule price of prediction
first}\label{the-joule-price-of-prediction-first}

Result form. Split each act's metered joules into the template's share
and the rest, \(J_k = J_k^{R^*} + J_k^{\mathrm{act}}\). \(J_k^{R^*}\)
covers writing, holding and resetting \(R^*(a)\). Resetting a template
with \(H(R^*)\) bits of entropy costs at least \(H(R^*)\,k_B T\ln 2\)
{[}6, 24{]}. Memory that does not predict costs dissipation {[}54{]}. A
template raises \(\iota\) when the act joules it saves exceed its own
joules. The measured result is the curve \(\iota(w)\) over template
width \(w\), the bits of \(R^*\), with its maximizer \(w^*\) and the
ratio \(\Delta J^{\mathrm{act}}/\Delta J^{R^*}\) at each step.

Protocol. Fix tasks, gears, predicates and budgets. Vary \(w\) over a
preset grid. An act may stop as soon as its outcome is resolved against
the template. Log \texttt{template\_j} and the act's remaining joules
for every act, on the same meters.

Evidence in hand: confident prediction lowered human cortical oxygen use
by up to 12\%, measured by metabolic imaging (preprint) {[}254{]}. In
cellular networks, the bits of the past that best predict the future are
prohibitively costly {[}198{]}. In percept-action loops, work-efficient
agents balance prediction against forgetting {[}234{]}.

\subsubsection{8.11 First run: a kernel-reported acceptor on a
laptop}\label{first-run-a-kernel-reported-acceptor-on-a-laptop}

This run applies Sections 8.1 to 8.7 to one task class on one machine.
Code, task sets, receipts and analysis are in the site repository under
\texttt{experiments/agency-run}. The preregistration and its falsifiers
were committed first (e830bff). The pilot and the locked calibration
came next (5828420), and then the main run (35c967a).

\textbf{Method.}

\begin{itemize}
\tightlist
\item
  \textbf{Agent and task.} A Python process closes Sudoku puzzles. The
  completeness predicate \(C(z)\) passes when the grid is full, every
  row, column and box holds 1 to 9, and every given is kept. A confirmed
  task credits \(b = n_{\mathrm{empty}}\log_2 9\) bits, the number of
  fills the act chooses among. That is the capacity bound of Definition
  1. The test set holds 30 puzzles, 10 per stratum (easy, medium, hard),
  generated from a fixed seed. A separate set of 30 calibrates the
  selectors.
\item
  \textbf{Gears.} g0 fills naked singles. g1 adds hidden singles. g2
  adds depth-first search.
\item
  \textbf{Arms.} The σ-law selector takes
  \(a^* = \arg\max[H(a) - \lambda J(a)]\), with \(H(a)\) the calibrated
  pass probability given the gears already tried, times \(b\), and
  \(\lambda = b/B\). The fixed cascade (cheapest first) takes the
  cheapest gear whose calibrated pass probability is at least 0.5. Both
  fix \(R^*\) before each act, hash it with SHA-256 (Secure Hash
  Algorithm, 256-bit) and timestamp it. Both compare after the act and
  refuse with a receipt when no act fits the budget. The goal-blind
  baseline draws one gear at random, with no prediction, no comparison,
  no escalation and no refusal. The budget cuts it off.
\item
  \textbf{Budgets and seeds.} Per task,
  \(B \in \{0.5, 1.0, 2.0\} \times J_{g2}(q)\), the pilot median of g2
  alone. Levels 0.5 and 1.0 bind: 98\% and 40\% of pilot runs exceed
  them. Level 2.0 does not (0.7\%). There are 20 seeds per condition, in
  a seeded interleaved order. In all, 5,400 episodes ran, with 4,940
  acts.
\item
  \textbf{Energy label.} Every joule is \texttt{reported\_j}: the
  kernel's per-process energy estimate for the agent
  (\texttt{ri\_energy\_nj} from \texttt{proc\_pid\_rusage}), read before
  and after each act, with no privileges. The process is the agent's
  wall. The Sudoku grid is held inside that process, so the world \(W'\)
  is not outside \(S\) as Definition 2 requires. This run tests the
  kinematics and the gated ranking, not the definition. The run that
  tests the definition holds the grid in a second process behind a
  socket, or on a microcontroller. \texttt{measured\_j} is empty,
  because no wall meter was attached. \texttt{est\_j} is cycles times
  \(4.20 \times 10^{-10}\) J, the pilot's median ratio.
\end{itemize}

\textbf{Results.} Mean ± standard deviation over 20 seeds. A run is one
arm at one budget level over all 30 tasks. CV is the coefficient of
variation, standard deviation over mean.

{\def\LTcaptype{none} % do not increment counter
\begin{longtable}[]{@{}
  >{\raggedright\arraybackslash}p{(\linewidth - 18\tabcolsep) * \real{0.1000}}
  >{\raggedright\arraybackslash}p{(\linewidth - 18\tabcolsep) * \real{0.1000}}
  >{\raggedright\arraybackslash}p{(\linewidth - 18\tabcolsep) * \real{0.1000}}
  >{\raggedright\arraybackslash}p{(\linewidth - 18\tabcolsep) * \real{0.1000}}
  >{\raggedright\arraybackslash}p{(\linewidth - 18\tabcolsep) * \real{0.1000}}
  >{\raggedright\arraybackslash}p{(\linewidth - 18\tabcolsep) * \real{0.1000}}
  >{\raggedright\arraybackslash}p{(\linewidth - 18\tabcolsep) * \real{0.1000}}
  >{\raggedright\arraybackslash}p{(\linewidth - 18\tabcolsep) * \real{0.1000}}
  >{\raggedright\arraybackslash}p{(\linewidth - 18\tabcolsep) * \real{0.1000}}
  >{\raggedright\arraybackslash}p{(\linewidth - 18\tabcolsep) * \real{0.1000}}@{}}
\toprule\noalign{}
\begin{minipage}[b]{\linewidth}\raggedright
Arm
\end{minipage} & \begin{minipage}[b]{\linewidth}\raggedright
\(B/J_{g2}\)
\end{minipage} & \begin{minipage}[b]{\linewidth}\raggedright
Tasks closed of 30
\end{minipage} & \begin{minipage}[b]{\linewidth}\raggedright
\(X\) (bits)
\end{minipage} & \begin{minipage}[b]{\linewidth}\raggedright
\(J\) (reported, mJ)
\end{minipage} & \begin{minipage}[b]{\linewidth}\raggedright
\(\iota\) (bits per reported joule)
\end{minipage} & \begin{minipage}[b]{\linewidth}\raggedright
CV of \(\iota\)
\end{minipage} & \begin{minipage}[b]{\linewidth}\raggedright
Refusals
\end{minipage} & \begin{minipage}[b]{\linewidth}\raggedright
Cuts
\end{minipage} & \begin{minipage}[b]{\linewidth}\raggedright
Overspends
\end{minipage} \\
\midrule\noalign{}
\endhead
\bottomrule\noalign{}
\endlastfoot
σ-law & 0.5 & 0.60 ± 0.50 & 108 ± 91 & 20.8 ± 0.3 & 5,254 ± 4,402 & 0.84
& 29.40 & 0 & 0 \\
σ-law & 1.0 & 9.10 ± 1.86 & 1,479 ± 316 & 43.8 ± 4.0 & 34,454 ± 10,007 &
0.29 & 20.90 & 0 & 9 \\
σ-law & 2.0 & 28.40 ± 0.60 & 4,823 ± 105 & 66.2 ± 5.8 & 73,345 ± 6,443 &
0.09 & 1.60 & 0 & 0 \\
Fixed cascade (cheapest first) & 0.5 & 0.70 ± 0.47 & 126 ± 85 & 18.8 ±
0.3 & 6,801 ± 4,569 & 0.67 & 29.30 & 0 & 0 \\
Fixed cascade (cheapest first) & 1.0 & 6.45 ± 1.50 & 1,047 ± 254 & 44.7
± 3.6 & 23,804 ± 7,067 & 0.30 & 23.55 & 0 & 7 \\
Fixed cascade (cheapest first) & 2.0 & 27.50 ± 0.51 & 4,668 ± 92 & 67.0
± 6.7 & 70,418 ± 7,766 & 0.11 & 2.50 & 0 & 0 \\
Goal-blind & 0.5 & 0.25 ± 0.44 & 45 ± 80 & 14.1 ± 3.1 & 2,775 ± 4,999 &
1.80 & 0 & 20.35 & 20 \\
Goal-blind & 1.0 & 3.85 ± 1.84 & 623 ± 301 & 26.6 ± 5.7 & 24,153 ±
13,364 & 0.55 & 0 & 17.45 & 0 \\
Goal-blind & 2.0 & 18.25 ± 2.71 & 3,034 ± 471 & 34.1 ± 9.7 & 92,724 ±
15,690 & 0.17 & 0 & 1.35 & 0 \\
\end{longtable}
}

\begin{itemize}
\tightlist
\item
  \textbf{Per act.} 1,902 acts were confirmed. The median act \(\iota\)
  was \(1.53 \times 10^5\) bits per reported joule, with interquartile
  range \(8.3 \times 10^4\) to \(4.7 \times 10^5\). Joules per confirmed
  bit were at least \(3.8 \times 10^{14}\) times \(k_B T\ln 2\) at 300
  K, with median \(2.3 \times 10^{15}\) times. The floor allows
  \(3.5 \times 10^{20}\) bits per joule.
\item
  \textbf{Joules per closed task.} The median over arms was 0.34 to 0.38
  mJ for easy tasks, 1.58 to 1.70 mJ for medium and 4.05 to 4.28 mJ for
  hard.
\item
  \textbf{Selectors, paired by seed (exploratory, post hoc).} This
  analysis was chosen after the run, so it is exploratory. At
  \(B = 1.0\,J_{g2}\) the σ-law selector closed 2.65 more tasks per run
  than the fixed cascade (95\% bootstrap interval 1.5 to 3.8). Its run
  \(\iota\) was higher by 10,650 bits per joule (4,747 to 16,592). At
  \(B = 2.0\,J_{g2}\) it closed 0.9 more (0.55 to 1.25). The σ-law rule
  goes straight to the gear that suffices, so the joules of failed light
  attempts stay in the budget. Joules per closed task were equal: on the
  4 tasks both arms closed at \(B = 1.0\,J_{g2}\), the paired interval
  ran from \(-7.9\) to \(+15.5\) µJ (R-F1).
\item
  \textbf{Package cross-check.} powermetrics, run non-interactively over
  30 s windows, read about 12.1 W of CPU power and 13.6 W of GPU power
  from other work on the machine while the agent slept. The agent added
  about 0.6 W on its own meter. Background CPU power moved by about 0.5
  W between windows, so idle-subtracted package joules came out
  negative. The per-process meter is the run's meter. D1 and D5 need a
  wall meter on an otherwise idle machine.
\end{itemize}

\textbf{Falsifier outcomes.} All were fixed before the run.

{\def\LTcaptype{none} % do not increment counter
\begin{longtable}[]{@{}
  >{\raggedright\arraybackslash}p{(\linewidth - 4\tabcolsep) * \real{0.3333}}
  >{\raggedright\arraybackslash}p{(\linewidth - 4\tabcolsep) * \real{0.3333}}
  >{\raggedright\arraybackslash}p{(\linewidth - 4\tabcolsep) * \real{0.3333}}@{}}
\toprule\noalign{}
\begin{minipage}[b]{\linewidth}\raggedright
ID
\end{minipage} & \begin{minipage}[b]{\linewidth}\raggedright
Observation
\end{minipage} & \begin{minipage}[b]{\linewidth}\raggedright
Outcome
\end{minipage} \\
\midrule\noalign{}
\endhead
\bottomrule\noalign{}
\endlastfoot
R-F1 & 95\% interval of the paired per-task difference in median
\(J^*\), σ-law minus fixed cascade, includes zero at every binding level
& Fired. At \(B = 1.0\), 4 tasks closed by both: interval \(-7.9\) to
\(+15.5\) µJ. At \(B = 0.5\), no task closed by both. On tasks both
close, the joules are equal \\
R-F2 & Any episode ends above its budget & Fired. 36 of 5,400. σ-law 9
and fixed cascade 7 at \(B = 1.0\), 0.07\% to 3.6\% over, when the final
comparison ran outside the guard. Goal-blind 20 at \(B = 0.5\), up to
33\% over, on budgets a few meter checks wide \\
R-F3 & Refusal recall below 0.9 on infeasible task levels, or over 10\%
of refusals on task levels some arm closes in half the seeds & Fired.
Recall 1.00 on 41 infeasible task levels. 13.1\% of refusals fell on
task levels some arm closed \\
R-F4 & Kendall's tau-b between \texttt{est\_j} and \texttt{reported\_j}
per act below 0.8 & Held. 0.96 on 4,000 acts \\
R-F5 & CV of run \(\iota\) above 0.25 in any condition & Fired at
binding levels (0.29 to 1.80). Held at \(B = 2.0\) (0.09, 0.11, 0.17) \\
R-F6 & Ranking by \(J^*\) differs from ranking by integrated \(\iota\) &
Fired. At \(B = 2.0\), \(J^*\) ranks σ-law, fixed cascade, goal-blind.
Integrated \(\iota\) ranks goal-blind first. At \(B = 1.0\) the last two
swap \\
R-F7 & A confirmed act below \(k_B T\ln 2\) per bit & Held. The lowest
act sat \(3.8 \times 10^{14}\) times above \\
R-F8 & Goal-blind run \(\iota\) at least the better selector's at a
binding level & Held. 2,775 against 6,801 at \(B = 0.5\). 24,153 against
34,454 at \(B = 1.0\) \\
\end{longtable}
}

\textbf{What the run settles.}

\begin{itemize}
\tightlist
\item
  \(\iota\) is read off receipts, act by act and run by run, as
  Propositions 5 and 6 require. On a loose budget it is stable across
  seeds.
\item
  \(J^*\) is the ranking, and integrated \(\iota\) is not. The
  goal-blind agent had the highest run \(\iota\) on the loose budget,
  closing 18 of 30 tasks against 28. Its failures were cheap. \(J^*\) is
  gated by \(C(z)\), so it ranks that agent last. This is why Section
  2.3 ranks by \(J^*\), with \(X_{\mathrm{req}}\) as a gate.
\item
  Exploratory, post hoc: under a binding budget the σ-law selector
  closed more tasks than the fixed cascade, and joules per closed task
  were equal. The gain lies in which tasks fit the budget, not in a
  cheaper closure of the same task. Section 8.11.2 states what the later
  runs replicate.
\item
  Two implementation fixes come before the next run. The comparison and
  the receipt write move inside the guarded region. A measured closing
  cost is charged before each act.
\end{itemize}

Figure 4 sets the gate beside the measured ranking at the loose budget.

\includegraphics[width=0.95\linewidth,height=\textheight,keepaspectratio]{fig/gate-then-rank.png}

\emph{Figure 4. Gate first, then rank by joules. Diagram with Run 1 data
at \(B = 2.0\,J_{g2}\), means over 20 seeds from the results table
above; integrated \(\iota\) alone would rank the goal-blind agent first
(R-F6). Source: \texttt{experiments/agency-run/results}. A
\href{https://research.openie.dev/figures/agency/gate-then-rank.png}{PNG
copy} is provided.}

\subsubsection{8.11.1 Second run: guarded acceptor, 40
seeds}\label{second-run-guarded-acceptor-40-seeds}

This run is the same task class, the same gears and the same three arms
as Section 8.11, on the same laptop. The preregistration for the second
run was committed first (610dd4c). The pilot and the locked calibration
came next (c88511f), and then the main run (920d7e2). Prediction 1 has
no software analogue in this run: the agent has no declared act cost
that falls with act time and no standing cost it controls, so the
preregistration states that the dose-response design is not run.

\textbf{Method.}

\begin{itemize}
\tightlist
\item
  \textbf{Guard.} Every joule of an episode is charged inside the guard:
  selection, hash, gear steps, comparison, receipt and final read. The
  opening check uses spent \(= 0\). The law arms refuse when the guard
  fails. The goal-blind arm is cut.
\item
  \textbf{Refusal threshold.} Both selectors use the \(r\)-quantile of
  calibration costs for budget-fit tests. \(r\) was locked from
  calibration replay before the main run: \(r = 0.1\), with recall 1.0
  and precision about 0.976 on the calibration split.
\item
  \textbf{Budgets and seeds.} Per task,
  \(B \in \{0.5, 1.0, 2.0\} \times J_{g2}(q)\), the pilot median of g2
  alone. Levels 0.5 and 1.0 bind: pilot fractions exceeding them are
  0.98 and 0.4. Level 2.0 does not (0.0067). There are 40 seeds per
  condition. In all, 10,800 episodes ran.
\item
  \textbf{Quiet window.} Quiet was not achieved. Probes ran until the
  limit. The last probe before the start, at 21:46 local time, recorded
  other-process CPU near 621\% of one core, with mediaanalysisd and
  editor helpers at the top. Baseline package power was about 21.4 W.
  Package power in the main window was about 22.7 W. The run records
  \texttt{quiet\_window\ =\ no}.
\item
  \textbf{Energy label.} Every joule is \texttt{reported\_j}: the
  kernel's per-process energy estimate for the agent
  (\texttt{ri\_energy\_nj} from \texttt{proc\_pid\_rusage}), with
  package figures from powermetrics where used for the cross-check.
  \texttt{est\_j} is cycles times \(1.783 \times 10^{-9}\) J, locked in
  the pilot. \texttt{measured\_j} is empty: no wall meter was attached.
\item
  \textbf{Package cross-check.} In the main window the process meter
  read 42.35 J. Net package joules were 121.81 J, a ratio of about 2.88
  to the process. Baseline package drift was about \(-0.082\) W. Because
  \texttt{quiet\_window\ =\ no}, the cross-check is declared not
  achievable. The per-process meter remains the run's meter.
\end{itemize}

\textbf{Results.} Mean ± standard deviation over 40 seeds. A run is one
arm at one budget level over all 30 tasks. CV is the coefficient of
variation, standard deviation over mean. \(J\) is reported in
millijoules for comparison with Section 8.11.

{\def\LTcaptype{none} % do not increment counter
\begin{longtable}[]{@{}
  >{\raggedright\arraybackslash}p{(\linewidth - 18\tabcolsep) * \real{0.1000}}
  >{\raggedright\arraybackslash}p{(\linewidth - 18\tabcolsep) * \real{0.1000}}
  >{\raggedright\arraybackslash}p{(\linewidth - 18\tabcolsep) * \real{0.1000}}
  >{\raggedright\arraybackslash}p{(\linewidth - 18\tabcolsep) * \real{0.1000}}
  >{\raggedright\arraybackslash}p{(\linewidth - 18\tabcolsep) * \real{0.1000}}
  >{\raggedright\arraybackslash}p{(\linewidth - 18\tabcolsep) * \real{0.1000}}
  >{\raggedright\arraybackslash}p{(\linewidth - 18\tabcolsep) * \real{0.1000}}
  >{\raggedright\arraybackslash}p{(\linewidth - 18\tabcolsep) * \real{0.1000}}
  >{\raggedright\arraybackslash}p{(\linewidth - 18\tabcolsep) * \real{0.1000}}
  >{\raggedright\arraybackslash}p{(\linewidth - 18\tabcolsep) * \real{0.1000}}@{}}
\toprule\noalign{}
\begin{minipage}[b]{\linewidth}\raggedright
Arm
\end{minipage} & \begin{minipage}[b]{\linewidth}\raggedright
\(B/J_{g2}\)
\end{minipage} & \begin{minipage}[b]{\linewidth}\raggedright
Tasks closed of 30
\end{minipage} & \begin{minipage}[b]{\linewidth}\raggedright
\(X\) (bits)
\end{minipage} & \begin{minipage}[b]{\linewidth}\raggedright
\(J\) (reported, mJ)
\end{minipage} & \begin{minipage}[b]{\linewidth}\raggedright
\(\iota\) (bits per reported joule)
\end{minipage} & \begin{minipage}[b]{\linewidth}\raggedright
CV of \(\iota\)
\end{minipage} & \begin{minipage}[b]{\linewidth}\raggedright
Refusals
\end{minipage} & \begin{minipage}[b]{\linewidth}\raggedright
Cuts
\end{minipage} & \begin{minipage}[b]{\linewidth}\raggedright
Overspends
\end{minipage} \\
\midrule\noalign{}
\endhead
\bottomrule\noalign{}
\endlastfoot
σ-law & 0.5 & 0.62 ± 0.49 & 113 ± 89 & 59.5 ± 2.0 & 1,942 ± 1,525 &
0.785 & 29.38 & 0 & 0 \\
σ-law & 1.0 & 2.90 ± 1.74 & 523 ± 314 & 133.7 ± 13.8 & 4,150 ± 2,892 &
0.697 & 27.10 & 0 & 0 \\
σ-law & 2.0 & 22.73 ± 1.30 & 3,980 ± 206 & 197.3 ± 18.2 & 20,359 ± 2,388
& 0.117 & 7.28 & 0 & 0 \\
Fixed cascade (cheapest first) & 0.5 & 0.57 ± 0.50 & 104 ± 90 & 60.1 ±
2.1 & 1,777 ± 1,548 & 0.871 & 29.43 & 0 & 0 \\
Fixed cascade (cheapest first) & 1.0 & 2.67 ± 1.25 & 483 ± 226 & 140.3 ±
13.5 & 3,593 ± 1,948 & 0.542 & 27.32 & 0 & 0 \\
Fixed cascade (cheapest first) & 2.0 & 19.93 ± 1.42 & 3,509 ± 227 &
201.2 ± 20.4 & 17,684 ± 2,909 & 0.164 & 10.07 & 0 & 0 \\
Goal-blind & 0.5 & 0.12 ± 0.33 & 23 ± 61 & 35.4 ± 11.6 & 513 ± 1,397 &
2.723 & 0 & 22.85 & 0 \\
Goal-blind & 1.0 & 1.00 ± 0.93 & 181 ± 169 & 72.0 ± 19.5 & 2,661 ± 2,579
& 0.969 & 0 & 20.00 & 0 \\
Goal-blind & 2.0 & 10.65 ± 2.27 & 1,859 ± 400 & 97.3 ± 30.1 & 20,302 ±
5,766 & 0.284 & 0 & 8.70 & 0 \\
\end{longtable}
}

\begin{itemize}
\tightlist
\item
  \textbf{Per act.} 7,960 acts ran. 2,448 were confirmed. Kendall's
  tau-b between \texttt{est\_j} and \texttt{reported\_j} was 0.942 on
  4,000 acts. The median act \(\iota\) was about \(3.11 \times 10^{4}\)
  bits per reported joule. Every confirmed act sat far above
  \(k_B T\ln 2\) per bit: the lowest was about \(6.4 \times 10^{14}\)
  times that floor.
\item
  \textbf{Refusal.} Recall was 1.0 on 49 infeasible pairs. There were
  5,223 refusals. The share on tasks closed by some arm was about 0.081.
\item
  \textbf{Overspends.} Zero episodes ended above budget. The maximum
  overspend fraction was \(-0.03684\): every episode finished under
  budget.
\item
  \textbf{Selectors, paired by seed.} This analysis was declared in the
  preregistration. At \(B = 2.0\,J_{g2}\) the σ-law selector closed 2.8
  more tasks per run than the fixed cascade (95\% bootstrap interval 2.3
  to 3.275). Its run \(\iota\) was higher by 2,674 bits per joule (1,831
  to 3,504). On 39 of 40 seeds the σ-law selector closed more. At
  \(B = 0.5\) and \(B = 1.0\) the intervals for the closed-task
  difference include zero. These results fire no falsifier.
\end{itemize}

\textbf{Falsifier outcomes.} All were fixed before the run. The table
states the second-run outcome beside the first-run outcome.

{\def\LTcaptype{none} % do not increment counter
\begin{longtable}[]{@{}
  >{\raggedright\arraybackslash}p{(\linewidth - 6\tabcolsep) * \real{0.2500}}
  >{\raggedright\arraybackslash}p{(\linewidth - 6\tabcolsep) * \real{0.2500}}
  >{\raggedright\arraybackslash}p{(\linewidth - 6\tabcolsep) * \real{0.2500}}
  >{\raggedright\arraybackslash}p{(\linewidth - 6\tabcolsep) * \real{0.2500}}@{}}
\toprule\noalign{}
\begin{minipage}[b]{\linewidth}\raggedright
ID
\end{minipage} & \begin{minipage}[b]{\linewidth}\raggedright
Observation
\end{minipage} & \begin{minipage}[b]{\linewidth}\raggedright
Run 2
\end{minipage} & \begin{minipage}[b]{\linewidth}\raggedright
Run 1
\end{minipage} \\
\midrule\noalign{}
\endhead
\bottomrule\noalign{}
\endlastfoot
R-F1 & At every binding level, the 95\% interval of the paired per-task
difference in median \(J^*\), σ-law minus fixed cascade, includes or
exceeds zero & Held, because no interval could be formed: at \(B = 1.0\)
no task closed by both, and at \(B = 0.5\) one did. Shared-task joules
show no σ-law saving & Fired \\
R-F2 & Any episode ends above its budget & Held. 0 of 10,800. The guard
charges comparison, receipt and final read & Fired \\
R-F3 & Refusal recall below 0.9 on infeasible pairs, or over 10\% of
refusals on pairs some arm closes & Held. Recall 1.0 on 49 infeasible
pairs. Share on tasks closed by some arm about 0.081. The \(r\)-quantile
refusal threshold is the change from Run 1 & Fired \\
R-F4 & Kendall's tau-b between \texttt{est\_j} and \texttt{reported\_j}
per act below 0.8 & Held. 0.942 on 4,000 acts & Held \\
R-F5 & CV of run \(\iota\) above 0.25 in any condition & Fired at
binding levels (0.54 to 2.72) and for the goal-blind arm at \(B = 2.0\)
(0.284). Held for both selectors at \(B = 2.0\) (0.117, 0.164) &
Fired \\
R-F6 & Ranking by \(J^*\) differs from ranking by integrated \(\iota\) &
Fired. At \(B = 2.0\), \(J^*\) ranks σ-law, fixed cascade, goal-blind.
Integrated \(\iota\) ranks σ-law, goal-blind, fixed cascade & Fired \\
R-F7 & A confirmed act below \(k_B T\ln 2\) per bit & Held. The lowest
act sat about \(6.4 \times 10^{14}\) times above & Held \\
R-F8 & Goal-blind run \(\iota\) at least the better selector's at a
binding level & Held. 513 against 1,942 at \(B = 0.5\). 2,661 against
4,150 at \(B = 1.0\) & Held \\
\end{longtable}
}

\textbf{What the run settles.}

\begin{itemize}
\tightlist
\item
  The guard closes R-F2. Zero overspends.
\item
  The refusal quantile closes R-F3. Recall stays 1.0 and the share of
  refusals on closable tasks falls under 10\%.
\item
  \(J^*\) remains the ranking. Integrated \(\iota\) still disagrees with
  it on the loose budget (R-F6).
\item
  Run \(\iota\) is still unstable across seeds at binding budgets
  (R-F5).
\item
  R-F1 held only because too few tasks were closed by both selectors to
  form an interval. No run shows the σ-law selector closing a shared
  task in fewer joules. The paired-by-seed gain at \(B = 2.0\) is more
  closed tasks, not cheaper closures of the same tasks.
\end{itemize}

\subsubsection{8.11.2 Second machine: a replication with a package
energy
counter}\label{second-machine-a-replication-with-a-package-energy-counter}

This run repeats the second-run protocol of Section 8.11.1 on a second
machine with a different energy counter. It is a replication of the
selector comparison, not a measurement of energy. Code, receipts and
summaries are in \texttt{experiments/agency-run}:
\texttt{PREREGISTRATION\_v2\_jetson.md} and its Amendment 1,
\texttt{results\_jetson/summary.md} (part A) and
\texttt{results\_jetson\_cyc/summary.md} (part B). The preregistration
addendum was committed first (cb059f0), then the meter backend and
harness (1938577), then Amendment 1 (7703822). The part A pilot
(6ddd086) and main run (a652503) came next, then the part B pilot
(3fde1fb) and main run (1bb2444).

\textbf{Method.}

\begin{itemize}
\tightlist
\item
  \textbf{Machine.} A Linux laptop with a 14-core x86 processor: 6
  performance cores with 2 threads each and 8 efficiency cores. Turbo
  was off. Every core was pinned to its base frequency: 2.6 GHz on
  performance cores and 1.9 GHz on efficiency cores. The agent ran
  pinned to one performance-core thread, one episode at a time. The
  median agent-core frequency during blocks was 2.6 GHz. The original
  settings were restored after the run.
\item
  \textbf{Unchanged.} Tasks, predicate, coupled bits, gears g0 to g2,
  the three arms, the guard, the refusal replay, the budget levels
  \(B \in \{0.5, 1.0, 2.0\} \times J_{g2}(q)\), 40 seeds and 10,800
  episodes per part, falsifiers R-F1 to R-F8 and the paired-by-seed
  analysis are those of Section 8.11.1.
\item
  \textbf{Energy counter.} RAPL (Running Average Power Limit) is the
  processor's on-chip package energy counter. Every interval was charged
  the package joules minus the idle baseline power of that block times
  the interval's seconds. RAPL is a model and sensor inside the chip,
  not a wall meter. It counts every process on the package.
\item
  \textbf{Quiet window.} Quiet was not achieved:
  \texttt{quiet\_window\ =\ no}. A CI (continuous integration) runner
  built Rust code on the machine throughout. The 40 probes before the
  start saw other processes using from 1.0 to 20.0 cores and package
  power from 2.3 to 28.7 W. The per-block gate was met before 180 of 360
  blocks in part A and 244 of 360 in part B. 98 and 208 blocks were
  quiet for their whole length.
\item
  \textbf{Part A.} The agent's meter is RAPL read before and after each
  act. Every figure is \texttt{reported\_j}.
\item
  \textbf{Part B.} Amendment 1 added this part before any pilot, because
  a harness test showed that acts last about 0.01 to 6 ms while RAPL
  updates about every 1 ms in steps of about 5 to 60 mJ. The agent's
  meter is \texttt{est\_j}: the agent thread's user cycles times \(c\).
  \(c\) was calibrated on RAPL in the part B pilot, net of idle, over 8
  rounds of 3 s, on the second attempt: \(c = 6.576 \times 10^{-10}\) J
  per cycle, from 41.11 J of net package energy over 62,516,635,701
  cycles. Every part B joule figure is \texttt{est\_j}. RAPL figures
  logged beside it stay \texttt{reported\_j}.
\item
  \textbf{Labels.} \texttt{measured\_j} is empty in both parts. No wall
  meter or shunt meter was attached.
\item
  \textbf{Arm names.} Arm names and their code fields follow Section
  8.3. The goal-blind baseline is \texttt{blind}.
\end{itemize}

\textbf{Part A results: RAPL read per act.} Mean ± standard deviation
over 40 seeds. \(J\) is \texttt{reported\_j} in joules, as in the
summary. Levels 0.5, 1.0 and 2.0 all bind: 0.647, 0.40 and 0.353 of
pilot g2 runs exceed them.

{\def\LTcaptype{none} % do not increment counter
\begin{longtable}[]{@{}
  >{\raggedright\arraybackslash}p{(\linewidth - 18\tabcolsep) * \real{0.1000}}
  >{\raggedright\arraybackslash}p{(\linewidth - 18\tabcolsep) * \real{0.1000}}
  >{\raggedright\arraybackslash}p{(\linewidth - 18\tabcolsep) * \real{0.1000}}
  >{\raggedright\arraybackslash}p{(\linewidth - 18\tabcolsep) * \real{0.1000}}
  >{\raggedright\arraybackslash}p{(\linewidth - 18\tabcolsep) * \real{0.1000}}
  >{\raggedright\arraybackslash}p{(\linewidth - 18\tabcolsep) * \real{0.1000}}
  >{\raggedright\arraybackslash}p{(\linewidth - 18\tabcolsep) * \real{0.1000}}
  >{\raggedright\arraybackslash}p{(\linewidth - 18\tabcolsep) * \real{0.1000}}
  >{\raggedright\arraybackslash}p{(\linewidth - 18\tabcolsep) * \real{0.1000}}
  >{\raggedright\arraybackslash}p{(\linewidth - 18\tabcolsep) * \real{0.1000}}@{}}
\toprule\noalign{}
\begin{minipage}[b]{\linewidth}\raggedright
Arm
\end{minipage} & \begin{minipage}[b]{\linewidth}\raggedright
\(B/J_{g2}\)
\end{minipage} & \begin{minipage}[b]{\linewidth}\raggedright
Tasks closed of 30
\end{minipage} & \begin{minipage}[b]{\linewidth}\raggedright
\(X\) (bits)
\end{minipage} & \begin{minipage}[b]{\linewidth}\raggedright
\(J\) (reported, J)
\end{minipage} & \begin{minipage}[b]{\linewidth}\raggedright
\(\iota\) (bits per reported joule)
\end{minipage} & \begin{minipage}[b]{\linewidth}\raggedright
CV of \(\iota\)
\end{minipage} & \begin{minipage}[b]{\linewidth}\raggedright
Refusals
\end{minipage} & \begin{minipage}[b]{\linewidth}\raggedright
Cuts
\end{minipage} & \begin{minipage}[b]{\linewidth}\raggedright
Overspends
\end{minipage} \\
\midrule\noalign{}
\endhead
\bottomrule\noalign{}
\endlastfoot
σ-law & 0.5 & 0.00 & 0 & 0.0018 & 0 & n/a & 30.00 & 0.00 & 323 \\
σ-law & 1.0 & 0.00 & 0 & 0.0002 & 0 & n/a & 30.00 & 0.00 & 322 \\
σ-law & 2.0 & 1.00 ± 0.39 & 181 ± 71 & 0.0018 ± 0.0336 & 14,427 ± 18,680
& 1.295 & 29.00 & 0.00 & 322 \\
Fixed cascade (cheapest first) & 0.5 & 0.00 & 0 & -0.0008 & 0 & n/a &
30.00 & 0.00 & 326 \\
Fixed cascade (cheapest first) & 1.0 & 0.00 & 0 & -0.0005 & 0 & n/a &
30.00 & 0.00 & 324 \\
Fixed cascade (cheapest first) & 2.0 & 1.00 ± 0.51 & 181 ± 91 & 0.0107 ±
0.0286 & 15,451 ± 23,767 & 1.538 & 29.00 & 0.00 & 320 \\
Goal-blind & 0.5 & 0.00 & 0 & 0.0026 & 0 & n/a & 0.00 & 30.00 & 327 \\
Goal-blind & 1.0 & 0.00 & 0 & 0.0009 & 0 & n/a & 0.00 & 30.00 & 323 \\
Goal-blind & 2.0 & 0.28 ± 0.45 & 50 ± 82 & 0.0019 ± 0.0180 & 3,801 ±
13,749 & 3.618 & 0.00 & 28.90 & 321 \\
\end{longtable}
}

\begin{itemize}
\tightlist
\item
  \textbf{Counter resolution.} The median episode lasted about 0.04 ms.
  10,075 of 10,800 episodes saw no change in the package counter. 10,549
  had \texttt{reported\_j} at or below zero after idle subtraction.
\item
  \textbf{Budgets.} Pilot medians of \(J_{g2}\) ranged from \(-0.0110\)
  to \(0.0914\) J, median \(0.0097\) J. The guard's open threshold was
  0.117 J. Budgets fell below it for 30, 30 and 28 of 30 tasks at the
  three levels, so the law arms refused nearly every task.
\item
  \textbf{Overspends.} Every condition recorded 320 to 327 overspends.
  The maximum overspend fraction was 10.86. A charge that reads zero for
  most acts and one whole counter step for a few cannot hold a budget of
  a few millijoules.
\item
  \textbf{Per act.} 137 acts ran. 91 were confirmed. 44 acts read zero.
  26 confirmed acts read zero or less. Kendall's tau-b between
  \texttt{est\_j} and \texttt{reported\_j} was 0.368 on 93 acts.
\item
  \textbf{Cross-check.} Block package net was 3.445 J against 2.737 J
  summed over episodes, a ratio of 0.794. Idle package power was 4.01 W
  before the run and 5.55 W after. The cross-check is declared not
  achievable.
\item
  \textbf{Selectors, paired by seed.} At \(B = 2.0\,J_{g2}\) the
  closed-task difference was 0 (95\% bootstrap interval \(-0.15\) to
  \(0.15\)). The run \(\iota\) difference was 1,328 bits per reported
  joule (\(-8,971\) to \(12,581\)). 30 of 40 seeds tied. At 0.5 and 1.0
  no arm closed a task.
\end{itemize}

\textbf{Part B results: cycle count times a RAPL-calibrated constant.}
Mean ± standard deviation over 40 seeds. \(J\) is \texttt{est\_j} in
joules, as in the summary. Levels 0.5 and 1.0 bind: 0.98 and 0.40 of
pilot g2 runs exceed them. Level 2.0 does not (0.0).

{\def\LTcaptype{none} % do not increment counter
\begin{longtable}[]{@{}
  >{\raggedright\arraybackslash}p{(\linewidth - 18\tabcolsep) * \real{0.1000}}
  >{\raggedright\arraybackslash}p{(\linewidth - 18\tabcolsep) * \real{0.1000}}
  >{\raggedright\arraybackslash}p{(\linewidth - 18\tabcolsep) * \real{0.1000}}
  >{\raggedright\arraybackslash}p{(\linewidth - 18\tabcolsep) * \real{0.1000}}
  >{\raggedright\arraybackslash}p{(\linewidth - 18\tabcolsep) * \real{0.1000}}
  >{\raggedright\arraybackslash}p{(\linewidth - 18\tabcolsep) * \real{0.1000}}
  >{\raggedright\arraybackslash}p{(\linewidth - 18\tabcolsep) * \real{0.1000}}
  >{\raggedright\arraybackslash}p{(\linewidth - 18\tabcolsep) * \real{0.1000}}
  >{\raggedright\arraybackslash}p{(\linewidth - 18\tabcolsep) * \real{0.1000}}
  >{\raggedright\arraybackslash}p{(\linewidth - 18\tabcolsep) * \real{0.1000}}@{}}
\toprule\noalign{}
\begin{minipage}[b]{\linewidth}\raggedright
Arm
\end{minipage} & \begin{minipage}[b]{\linewidth}\raggedright
\(B/J_{g2}\)
\end{minipage} & \begin{minipage}[b]{\linewidth}\raggedright
Tasks closed of 30
\end{minipage} & \begin{minipage}[b]{\linewidth}\raggedright
\(X\) (bits)
\end{minipage} & \begin{minipage}[b]{\linewidth}\raggedright
\(J\) (estimated, J)
\end{minipage} & \begin{minipage}[b]{\linewidth}\raggedright
\(\iota\) (bits per estimated joule)
\end{minipage} & \begin{minipage}[b]{\linewidth}\raggedright
CV of \(\iota\)
\end{minipage} & \begin{minipage}[b]{\linewidth}\raggedright
Refusals
\end{minipage} & \begin{minipage}[b]{\linewidth}\raggedright
Cuts
\end{minipage} & \begin{minipage}[b]{\linewidth}\raggedright
Overspends
\end{minipage} \\
\midrule\noalign{}
\endhead
\bottomrule\noalign{}
\endlastfoot
σ-law & 0.5 & 1.10 ± 0.98 & 200 ± 178 & 0.0345 ± 0.0017 & 6,027 ± 5,458
& 0.906 & 28.90 & 0.00 & 0 \\
σ-law & 1.0 & 8.75 ± 3.14 & 1,574 ± 565 & 0.0805 ± 0.0098 & 20,519 ±
8,516 & 0.415 & 21.25 & 0.00 & 0 \\
σ-law & 2.0 & 28.60 ± 0.98 & 4,882 ± 154 & 0.1008 ± 0.0234 & 50,667 ±
10,088 & 0.199 & 1.40 & 0.00 & 0 \\
Fixed cascade (cheapest first) & 0.5 & 1.15 ± 0.92 & 208 ± 167 & 0.0366
± 0.0013 & 5,833 ± 4,757 & 0.816 & 28.85 & 0.00 & 0 \\
Fixed cascade (cheapest first) & 1.0 & 9.22 ± 3.07 & 1,660 ± 551 &
0.0813 ± 0.0106 & 21,412 ± 8,130 & 0.380 & 20.77 & 0.00 & 0 \\
Fixed cascade (cheapest first) & 2.0 & 27.77 ± 2.30 & 4,745 ± 377 &
0.1099 ± 0.0339 & 47,108 ± 13,216 & 0.281 & 2.23 & 0.00 & 0 \\
Goal-blind & 0.5 & 0.20 ± 0.46 & 36 ± 84 & 0.0210 ± 0.0055 & 1,722 ±
3,927 & 2.280 & 0.00 & 21.50 & 0 \\
Goal-blind & 1.0 & 3.73 ± 2.12 & 667 ± 381 & 0.0436 ± 0.0140 & 16,393 ±
9,246 & 0.564 & 0.00 & 16.68 & 0 \\
Goal-blind & 2.0 & 15.65 ± 2.50 & 2,662 ± 433 & 0.0499 ± 0.0160 & 57,266
± 14,596 & 0.255 & 0.00 & 4.17 & 0 \\
\end{longtable}
}

\begin{itemize}
\tightlist
\item
  \textbf{Per act.} 7,800 acts ran. 3,847 were confirmed. The median act
  \(\iota\) was 69,892 bits per estimated joule, with interquartile
  range 35,759 to 146,697. Every confirmed act sat far above
  \(k_B T\ln 2\) per bit: the lowest was about \(3.6 \times 10^{14}\)
  times that floor, the median about \(5.0 \times 10^{15}\) times.
\item
  \textbf{Refusal.} Recall was 1.0 on 45 infeasible pairs. There were
  4,136 refusals. The share on tasks closed by some arm was about 0.071.
\item
  \textbf{Overspends.} Zero episodes ended above budget. The maximum
  overspend fraction was \(-0.03977\).
\item
  \textbf{Shared tasks.} At \(B = 1.0\), 9 tasks closed by both
  selectors. The paired per-task difference in median \(J^*\), σ-law
  minus fixed cascade, had mean \(1.02 \times 10^{-4}\) J (95\% interval
  \(-5.23 \times 10^{-4}\) to \(6.03 \times 10^{-4}\), Wilcoxon
  \(p = 0.214\)). At \(B = 0.5\) one task closed by both, with
  difference \(-6.68 \times 10^{-6}\) J. At \(B = 2.0\), 29 tasks closed
  by both, with mean difference \(-2.65 \times 10^{-4}\) J
  (\(-6.63 \times 10^{-4}\) to \(5.67 \times 10^{-5}\), \(p = 0.358\)).
\item
  \textbf{Selectors, paired by seed.} At \(B = 2.0\,J_{g2}\) the σ-law
  selector closed 0.825 more tasks per run than the fixed cascade (95\%
  bootstrap interval 0.125 to 1.575). Its run \(\iota\) was higher by
  3,559 bits per estimated joule (\(-1,068\) to \(8,320\)). 9 of 40
  seeds favored the σ-law selector and 27 tied. At \(B = 0.5\) the
  closed-task difference was \(-0.05\) (\(-0.45\) to \(0.35\)). At
  \(B = 1.0\) it was \(-0.475\) (\(-1.8\) to \(0.775\)).
\item
  \textbf{The estimate against the counter.} Kendall's tau-b between
  each episode's \texttt{est\_j} and the same episode's RAPL package net
  was \(-0.037\) on 4,000 episodes. 5,258 of 10,800 episodes saw no
  change in the package counter. Summed over blocks, \texttt{est\_j} was
  23.23 J against 17.14 J of RAPL package net, a ratio of 1.355. Idle
  package power was 9.37 W before the run and 3.17 W after. The
  cross-check is declared not achievable.
\end{itemize}

\textbf{Falsifier outcomes.} All were fixed before the run. Part B
replaces R-F4 with the episode test against RAPL, as Amendment 1 states.
In Run 2 and in the two parts here, ``Held'' means the observation did
not occur.

{\def\LTcaptype{none} % do not increment counter
\begin{longtable}[]{@{}
  >{\raggedright\arraybackslash}p{(\linewidth - 8\tabcolsep) * \real{0.2000}}
  >{\raggedright\arraybackslash}p{(\linewidth - 8\tabcolsep) * \real{0.2000}}
  >{\raggedright\arraybackslash}p{(\linewidth - 8\tabcolsep) * \real{0.2000}}
  >{\raggedright\arraybackslash}p{(\linewidth - 8\tabcolsep) * \real{0.2000}}
  >{\raggedright\arraybackslash}p{(\linewidth - 8\tabcolsep) * \real{0.2000}}@{}}
\toprule\noalign{}
\begin{minipage}[b]{\linewidth}\raggedright
ID
\end{minipage} & \begin{minipage}[b]{\linewidth}\raggedright
Observation
\end{minipage} & \begin{minipage}[b]{\linewidth}\raggedright
Part A (RAPL per act)
\end{minipage} & \begin{minipage}[b]{\linewidth}\raggedright
Part B (cycle model)
\end{minipage} & \begin{minipage}[b]{\linewidth}\raggedright
Run 2
\end{minipage} \\
\midrule\noalign{}
\endhead
\bottomrule\noalign{}
\endlastfoot
R-F1 & At every binding level, the 95\% interval of the paired per-task
difference in median \(J^*\), σ-law minus fixed cascade, includes or
exceeds zero & Held, because no interval could be formed: one task
closed by both, at \(B = 2.0\), with difference \(-0.0052\) J, and none
at the other levels & Held, because the observation needs every binding
level and \(B = 0.5\) has a single shared task (\(-6.68 \times 10^{-6}\)
J) and no interval. At \(B = 1.0\) the interval includes zero & Held \\
R-F2 & Any episode ends above its budget & Fired. 320 to 327 per
condition & Held. 0 of 10,800 & Held \\
R-F3 & Refusal recall below 0.9 on infeasible pairs, or over 10\% of
refusals on pairs some arm closes & Held. Recall 1.0 on 88 pairs. Share
0.0011 & Held. Recall 1.0 on 45 pairs. Share 0.071 & Held \\
R-F4 & Part A: Kendall's tau-b between \texttt{est\_j} and
\texttt{reported\_j} per act below 0.8. Part B: tau-b between episode
\texttt{est\_j} and episode RAPL net below 0.8 & Fired. 0.368 on 93 acts
& Fired. \(-0.037\) on 4,000 episodes & Held \\
R-F5 & CV of run \(\iota\) above 0.25 in any condition & Fired (1.295 to
3.618) & Fired. 0.380 to 2.280 at binding levels. At \(B = 2.0\): σ-law
0.199, fixed cascade 0.281, goal-blind 0.255 & Fired \\
R-F6 & Ranking by \(J^*\) differs from ranking by integrated \(\iota\) &
Fired. At \(B = 2.0\), \(J^*\) ranks σ-law, fixed cascade, goal-blind.
Integrated \(\iota\) ranks fixed cascade first & Fired. At \(B = 2.0\),
\(J^*\) ranks σ-law, fixed cascade, goal-blind. Integrated \(\iota\)
ranks goal-blind, σ-law, fixed cascade & Fired \\
R-F7 & A confirmed act below \(k_B T\ln 2\) per bit & Fired. 26
confirmed acts read zero or less, which places them below the floor.
That is counter quantization & Held. The lowest act sat about
\(3.6 \times 10^{14}\) times above & Held \\
R-F8 & Goal-blind run \(\iota\) at least the better selector's at a
binding level & Fired. At \(B = 0.5\) and \(1.0\) every arm reads 0 &
Held. 1,722 against 6,027 at \(B = 0.5\). 16,393 against 21,412 at
\(B = 1.0\) & Held \\
\end{longtable}
}

\textbf{What the replication settles.}

\begin{itemize}
\tightlist
\item
  Per-act RAPL cannot meter acts shorter than one millisecond. Most
  episodes saw no counter update, so the guard refused or overspent from
  quantization alone. Part A is a meter-resolution result. It says
  nothing about the selectors.
\item
  With a cycle meter, the selection result has the same sign as Run 2 at
  \(B = 2.0\): the σ-law selector closed more tasks per run. The gain is
  smaller (0.825 against 2.8), and level 2.0 did not bind on this
  machine. At the binding levels the intervals include zero.
\item
  The cycle model did not track the package counter episode by episode
  (tau-b \(-0.037\)), and its block total was 1.355 times the counter's
  net. Part B joules rank acts by cycles. They are not energy readings.
\item
  Shared-task joules show no σ-law saving here either. Where the σ-law
  selector gains, it gains in tasks closed, not in cheaper closure of
  the same task.
\item
  No figure in this section is \texttt{measured\_j}.
\end{itemize}

Figure 5 sets the closed-task results of the three runs side by side.

\includegraphics[width=0.95\linewidth,height=\textheight,keepaspectratio]{fig/run-results.png}

\emph{Figure 5. Tasks closed of 30 in three preregistered runs. Data
figure: mean ± one standard deviation over seeds, from the results
tables of Sections 8.11, 8.11.1 and 8.11.2 (part B). Sources:
\texttt{experiments/agency-run/results}, \texttt{results2} and
\texttt{results\_jetson\_cyc}. A
\href{https://research.openie.dev/figures/agency/run-results.png}{PNG
copy} is provided.}

\textbf{What replicates, and what does not.} The run 1 selector result
was exploratory. Two later runs tested it under preregistration.

{\def\LTcaptype{none} % do not increment counter
\begin{longtable}[]{@{}
  >{\raggedright\arraybackslash}p{(\linewidth - 4\tabcolsep) * \real{0.3333}}
  >{\raggedright\arraybackslash}p{(\linewidth - 4\tabcolsep) * \real{0.3333}}
  >{\raggedright\arraybackslash}p{(\linewidth - 4\tabcolsep) * \real{0.3333}}@{}}
\toprule\noalign{}
\begin{minipage}[b]{\linewidth}\raggedright
Claim from run 1 (exploratory)
\end{minipage} & \begin{minipage}[b]{\linewidth}\raggedright
Run 2 (preregistered)
\end{minipage} & \begin{minipage}[b]{\linewidth}\raggedright
Second machine, part B (preregistered)
\end{minipage} \\
\midrule\noalign{}
\endhead
\bottomrule\noalign{}
\endlastfoot
More tasks closed at \(B = 1.0\,J_{g2}\): +2.65 (1.5 to 3.8) & Not
replicated. The interval includes zero & Not replicated. \(-0.475\)
(\(-1.8\) to \(0.775\)) \\
More tasks closed at \(B = 2.0\,J_{g2}\): +0.9 (0.55 to 1.25) &
Replicated. +2.8 (2.3 to 3.275) & Replicated. +0.825 (0.125 to 1.575) \\
Joules per closed task & Equal. No shared-task saving & Equal.
Shared-task intervals include zero \\
\end{longtable}
}

The replicated result is narrow. At the loosest budget level the σ-law
selector closes more tasks than the fixed cascade, on two machines. No
run shows fewer joules for the same closed task, and the binding-budget
gain of run 1 did not replicate. Part A is a meter-resolution result and
tests neither claim.

\textbf{Measurement protocol for the next run.} Both parts fail at the
meter, so the protocol changes before the selectors are tested again.

\begin{enumerate}
\def\labelenumi{\arabic{enumi}.}
\tightlist
\item
  \textbf{Idle machine.} Stop nothing that belongs to someone else. Wait
  for a quiet window, or schedule one: other processes at or below 0.5
  cores for the whole run, with \texttt{quiet\_window\ =\ yes} recorded.
  If quiet is not reached, the run does not start.
\item
  \textbf{Wall or shunt meter.} Attach a wall power meter, or a shunt
  meter on the supply rail, sampled faster than one act. Only its
  readings fill \texttt{measured\_j}.
\item
  \textbf{Resolution check first.} Before the pilot, record the meter's
  update interval and step size against the median act length. If the
  meter cannot resolve one act, meter blocks of acts and attribute by
  count, and say so in the preregistration.
\item
  \textbf{Calibrate the estimate on the meter.} Fit \texttt{est\_j} on
  the wall or shunt meter at block level, and require R-F4 to hold
  before \texttt{est\_j} drives the guard.
\item
  \textbf{Pin and record.} Pin frequency, record idle baselines before
  and after, report drift, and keep the package counter as a labeled
  cross-check.
\item
  \textbf{Task closure is primary.} The next run's preregistration makes
  task closure the primary endpoint: the paired-by-seed difference in
  tasks closed, σ-law minus fixed cascade, at each budget level, with
  the Holm step-down correction across levels. Shared-task joules (R-F1)
  are secondary.
\end{enumerate}

\subsubsection{8.12 Reporting}\label{reporting}

Publish all receipts, raw meter traces, idle baselines and analysis code
with the result. Report measured, reported and estimated joules in
separate columns. Never sum across tiers. State which falsifiers fired,
with the same prominence as those that passed. Every run also answers
the checklist of Takahashi and Hayashi {[}239{]}: the accounting
boundary; the energy balance terms (stored energy and exported work, or
a statement that they are negligible); the baseline and whether joules
are total or incremental; the coarse-graining of outcomes; the horizon
and reset protocol; wall-clock time and mean power beside bits per
joule; and the estimator of coupled bits with its bias correction.

\subsection{References}\label{references}

Entries marked ``preprint'' are cited by their arXiv posting. Entries
published in a language other than English say so.

\begin{enumerate}
\def\labelenumi{\arabic{enumi}.}
\tightlist
\item
  Pearl, J. (1995). Causal diagrams for empirical research.
  \emph{Biometrika} 82:669-688. https://doi.org/10.1093/biomet/82.4.669
\item
  Everitt, T., Hutter, M., Kumar, R. and Krakovna, V. (2019). Reward
  tampering problems and solutions in reinforcement learning: a causal
  influence diagram perspective. arXiv:1908.04734 (preprint).
\item
  Sagawa, T. and Ueda, M. (2008). Second law of thermodynamics with
  discrete quantum feedback control. \emph{Physical Review Letters}
  100:080403. https://doi.org/10.1103/PhysRevLett.100.080403
\item
  Sagawa, T. and Ueda, M. (2010). Generalized Jarzynski equality under
  nonequilibrium feedback control. \emph{Physical Review Letters}
  104:090602. https://doi.org/10.1103/PhysRevLett.104.090602
\item
  Sagawa, T. and Ueda, M. (2009). Minimal energy cost for thermodynamic
  information processing: measurement and information erasure.
  \emph{Physical Review Letters} 102:250602.
  https://doi.org/10.1103/PhysRevLett.102.250602
\item
  Landauer, R. (1961). Irreversibility and heat generation in the
  computing process. \emph{IBM Journal of Research and Development}
  5:183-191. https://doi.org/10.1147/rd.53.0183
\item
  Bennett, C. H. (1982). The thermodynamics of computation: a review.
  \emph{International Journal of Theoretical Physics} 21:905-940.
  https://doi.org/10.1007/BF02084158
\item
  Ay, N. and Polani, D. (2008). Information flows in causal networks.
  \emph{Advances in Complex Systems} 11:17-41.
  https://doi.org/10.1142/S0219525908001465
\item
  Klyubin, A. S., Polani, D. and Nehaniv, C. L. (2005). Empowerment: a
  universal agent-centric measure of control. \emph{2005 IEEE Congress
  on Evolutionary Computation} 1:128-135.
  https://doi.org/10.1109/CEC.2005.1554676
\item
  Kenton, Z., Kumar, R., Farquhar, S., Richens, J. et al.~(2023).
  Discovering agents. \emph{Artificial Intelligence} 322:103963.
  https://doi.org/10.1016/j.artint.2023.103963
\item
  James, R. G., Barnett, N. and Crutchfield, J. P. (2016). Information
  flows? A critique of transfer entropies. \emph{Physical Review
  Letters} 116:238701. https://doi.org/10.1103/PhysRevLett.116.238701
\item
  Lizier, J. T. and Prokopenko, M. (2010). Differentiating information
  transfer and causal effect. \emph{The European Physical Journal B}
  73:605-615. https://doi.org/10.1140/epjb/e2010-00034-5
\item
  Janzing, D., Balduzzi, D., Grosse-Wentrup, M. and Schölkopf, B.
  (2013). Quantifying causal influences. \emph{The Annals of Statistics}
  41. https://doi.org/10.1214/13-AOS1145
\item
  Albantakis, L., Marshall, W., Hoel, E. and Tononi, G. (2019). What
  caused what? A quantitative account of actual causation using
  dynamical causal networks. \emph{Entropy} 21:459.
  https://doi.org/10.3390/e21050459
\item
  Aguilera, M., Millidge, B., Tschantz, A. and Buckley, C. L. (2022).
  How particular is the physics of the free energy principle?
  \emph{Physics of Life Reviews} 40:24-50.
  https://doi.org/10.1016/j.plrev.2021.11.001
\item
  Biehl, M., Pollock, F. A. and Kanai, R. (2021). A technical critique
  of some parts of the free energy principle. \emph{Entropy} 23:293.
  https://doi.org/10.3390/e23030293
\item
  Bruineberg, J., Dołęga, K., Dewhurst, J. and Baltieri, M. (2022). The
  emperor's new Markov blankets. \emph{Behavioral and Brain Sciences}
  45:e183. https://doi.org/10.1017/S0140525X21002351
\item
  Orseau, L., McGregor McGill, S. and Legg, S. (2018). Agents and
  devices: a relative definition of agency. arXiv:1805.12387 (preprint).
\item
  Albantakis, L., Barbosa, L., Findlay, G., Grasso, M. et al.~(2023).
  Integrated information theory (IIT) 4.0: formulating the properties of
  phenomenal existence in physical terms. \emph{PLOS Computational
  Biology} 19:e1011465. https://doi.org/10.1371/journal.pcbi.1011465
\item
  Krakauer, D., Bertschinger, N., Olbrich, E., Flack, J. C. and Ay, N.
  (2020). The information theory of individuality. \emph{Theory in
  Biosciences} 139:209-223. https://doi.org/10.1007/s12064-020-00313-7
\item
  Hoel, E. P., Albantakis, L. and Tononi, G. (2013). Quantifying causal
  emergence shows that macro can beat micro. \emph{Proceedings of the
  National Academy of Sciences} 110:19790-19795.
  https://doi.org/10.1073/pnas.1314922110
\item
  Rosas, F. E., Mediano, P. A. M., Biehl, M., Chandaria, S. and Polani,
  D. (2020). Causal blankets: theory and algorithmic framework. In
  \emph{Active Inference: IWAI 2020}, Communications in Computer and
  Information Science, Springer, pp.~187-198.
  https://doi.org/10.1007/978-3-030-64919-7\_19
\item
  Shiraishi, N., Funo, K. and Saito, K. (2018). Speed limit for
  classical stochastic processes. \emph{Physical Review Letters}
  121:070601. https://doi.org/10.1103/PhysRevLett.121.070601
\item
  Esposito, M. and Van den Broeck, C. (2011). Second law and Landauer
  principle far from equilibrium. \emph{EPL} 95:40004.
  https://doi.org/10.1209/0295-5075/95/40004
\item
  Barato, A. C. and Seifert, U. (2015). Thermodynamic uncertainty
  relation for biomolecular processes. \emph{Physical Review Letters}
  114:158101. https://doi.org/10.1103/PhysRevLett.114.158101
\item
  Gingrich, T. R., Horowitz, J. M., Perunov, N. and England, J. L.
  (2016). Dissipation bounds all steady-state current fluctuations.
  \emph{Physical Review Letters} 116:120601.
  https://doi.org/10.1103/PhysRevLett.116.120601
\item
  Paneru, G., Dutta, S., Tlusty, T. and Pak, H. K. (2020). Reaching and
  violating thermodynamic uncertainty bounds in information engines.
  \emph{Physical Review E} 102:032126.
  https://doi.org/10.1103/PhysRevE.102.032126
\item
  Potts, P. P. and Samuelsson, P. (2019). Thermodynamic uncertainty
  relations including measurement and feedback. \emph{Physical Review E}
  100:052137. https://doi.org/10.1103/PhysRevE.100.052137
\item
  Proesmans, K., Ehrich, J. and Bechhoefer, J. (2020). Finite-time
  Landauer principle. \emph{Physical Review Letters} 125:100602.
  https://doi.org/10.1103/PhysRevLett.125.100602
\item
  Sivak, D. A. and Crooks, G. E. (2012). Thermodynamic metrics and
  optimal paths. \emph{Physical Review Letters} 108:190602.
  https://doi.org/10.1103/PhysRevLett.108.190602
\item
  Margolus, N. and Levitin, L. B. (1998). The maximum speed of dynamical
  evolution. \emph{Physica D} 120:188-195.
  https://doi.org/10.1016/S0167-2789(98)00054-2
\item
  Touchette, H. and Lloyd, S. (2000). Information-theoretic limits of
  control. \emph{Physical Review Letters} 84:1156-1159.
  https://doi.org/10.1103/PhysRevLett.84.1156
\item
  Touchette, H. and Lloyd, S. (2004). Information-theoretic approach to
  the study of control systems. \emph{Physica A} 331:140-172.
  https://doi.org/10.1016/j.physa.2003.09.007
\item
  Garrahan, J. P. (2017). Simple bounds on fluctuations and uncertainty
  relations for first-passage times of counting observables.
  \emph{Physical Review E} 95:032134.
  https://doi.org/10.1103/PhysRevE.95.032134
\item
  Bennett, C. H. (1973). Logical reversibility of computation. \emph{IBM
  Journal of Research and Development} 17:525-532.
  https://doi.org/10.1147/rd.176.0525
\item
  Sagawa, T. and Ueda, M. (2012). Fluctuation theorem with information
  exchange: role of correlations in stochastic thermodynamics.
  \emph{Physical Review Letters} 109:180602.
  https://doi.org/10.1103/PhysRevLett.109.180602
\item
  Hartle, H., Wolpert, D., Stier, A., Kempes, C. P. and Manzano, G.
  (2024). Work extraction with feedback control using limited resources.
  arXiv:2407.05507 (preprint).
\item
  Horodecki, M. and Oppenheim, J. (2013). Fundamental limitations for
  quantum and nanoscale thermodynamics. \emph{Nature Communications}
  4:2059. https://doi.org/10.1038/ncomms3059
\item
  Parrondo, J. M. R., Horowitz, J. M. and Sagawa, T. (2015).
  Thermodynamics of information. \emph{Nature Physics} 11:131-139.
  https://doi.org/10.1038/nphys3230
\item
  Seifert, U. (2012). Stochastic thermodynamics, fluctuation theorems
  and molecular machines. \emph{Reports on Progress in Physics}
  75:126001. https://doi.org/10.1088/0034-4885/75/12/126001
\item
  Conant, R. C. and Ashby, W. R. (1970). Every good regulator of a
  system must be a model of that system. \emph{International Journal of
  Systems Science} 1:89-97. https://doi.org/10.1080/00207727008920220
\item
  Richens, J. and Everitt, T. (2024). Robust agents learn causal world
  models. \emph{International Conference on Learning Representations
  (ICLR 2024)}, oral, honorable mention. arXiv:2402.10877.
\item
  Richens, J., Abel, D., Bellot, A. and Everitt, T. (2025). General
  agents contain world models. arXiv:2506.01622 (preprint).
\item
  Horowitz, J. M. and Esposito, M. (2014). Thermodynamics with
  continuous information flow. \emph{Physical Review X} 4:031015.
  https://doi.org/10.1103/PhysRevX.4.031015
\item
  Ito, S. and Sagawa, T. (2013). Information thermodynamics on causal
  networks. \emph{Physical Review Letters} 111:180603.
  https://doi.org/10.1103/PhysRevLett.111.180603
\item
  Hartich, D., Barato, A. C. and Seifert, U. (2014). Stochastic
  thermodynamics of bipartite systems: transfer entropy inequalities and
  a Maxwell's demon interpretation. \emph{Journal of Statistical
  Mechanics} P02016. https://doi.org/10.1088/1742-5468/2014/02/P02016
\item
  Allahverdyan, A. E., Janzing, D. and Mahler, G. (2009). Thermodynamic
  efficiency of information and heat flow. \emph{Journal of Statistical
  Mechanics} P09011. https://doi.org/10.1088/1742-5468/2009/09/P09011
\item
  Esposito, M. and Schaller, G. (2012). Stochastic thermodynamics for
  ``Maxwell demon'' feedbacks. \emph{EPL} 99:30003.
  https://doi.org/10.1209/0295-5075/99/30003
\item
  Mandal, D. and Jarzynski, C. (2012). Work and information processing
  in a solvable model of Maxwell's demon. \emph{PNAS} 109:11641-11645.
  https://doi.org/10.1073/pnas.1204263109
\item
  Boyd, A. B., Mandal, D. and Crutchfield, J. P. (2016). Identifying
  functional thermodynamics in autonomous Maxwellian ratchets. \emph{New
  Journal of Physics} 18:023049.
  https://doi.org/10.1088/1367-2630/18/2/023049
\item
  Boyd, A. B., Mandal, D. and Crutchfield, J. P. (2018). Thermodynamics
  of modularity: structural costs beyond the Landauer bound.
  \emph{Physical Review X} 8:031036.
  https://doi.org/10.1103/PhysRevX.8.031036
\item
  Wolpert, D. H. and Kolchinsky, A. (2020). Thermodynamics of computing
  with circuits. \emph{New Journal of Physics} 22:063047.
  https://doi.org/10.1088/1367-2630/ab82b8
\item
  Kolchinsky, A. and Wolpert, D. H. (2017). Dependence of dissipation on
  the initial distribution over states. \emph{Journal of Statistical
  Mechanics} 083202. https://doi.org/10.1088/1742-5468/aa7ee1
\item
  Still, S., Sivak, D. A., Bell, A. J. and Crooks, G. E. (2012).
  Thermodynamics of prediction. \emph{Physical Review Letters}
  109:120604. https://doi.org/10.1103/PhysRevLett.109.120604
\item
  Kolchinsky, A. and Wolpert, D. H. (2018). Semantic information,
  autonomous agency and non-equilibrium statistical physics.
  \emph{Interface Focus} 8:20180041.
  https://doi.org/10.1098/rsfs.2018.0041
\item
  Ortega, P. A. and Braun, D. A. (2013). Thermodynamics as a theory of
  decision-making with information-processing costs. \emph{Proceedings
  of the Royal Society A} 469:20120683.
  https://doi.org/10.1098/rspa.2012.0683
\item
  Kolchinsky, A. and Wolpert, D. H. (2020). Thermodynamic costs of
  Turing machines. \emph{Physical Review Research} 2:033312.
  https://doi.org/10.1103/PhysRevResearch.2.033312
\item
  Kolchinsky, A. and Wolpert, D. H. (2021). Work, entropy production,
  and thermodynamics of information under protocol constraints.
  \emph{Physical Review X} 11:041024.
  https://doi.org/10.1103/PhysRevX.11.041024
\item
  Manzano, G., Kardeş, G., Roldán, É. and Wolpert, D. H. (2024).
  Thermodynamics of computations with absolute irreversibility,
  unidirectional transitions, and stochastic computation times.
  \emph{Physical Review X} 14:021026.
  https://doi.org/10.1103/PhysRevX.14.021026
\item
  Bérut, A., Arakelyan, A., Petrosyan, A., Ciliberto, S.,
  Dillenschneider, R. and Lutz, E. (2012). Experimental verification of
  Landauer's principle linking information and thermodynamics.
  \emph{Nature} 483:187-189. https://doi.org/10.1038/nature10872
\item
  Jun, Y., Gavrilov, M. and Bechhoefer, J. (2014). High-precision test
  of Landauer's principle in a feedback trap. \emph{Physical Review
  Letters} 113:190601. https://doi.org/10.1103/PhysRevLett.113.190601
\item
  Orlov, A. O., Lent, C. S., Thorpe, C. C., Boechler, G. P. and Snider,
  G. L. (2012). Experimental test of Landauer's principle at the
  sub-k\_BT level. \emph{Japanese Journal of Applied Physics} 51:06FE10.
  https://doi.org/10.1143/JJAP.51.06FE10
\item
  Hong, J., Lambson, B., Dhuey, S. and Bokor, J. (2016). Experimental
  test of Landauer's principle in single-bit operations on nanomagnetic
  memory bits. \emph{Science Advances} 2:e1501492.
  https://doi.org/10.1126/sciadv.1501492
\item
  Gavrilov, M. and Bechhoefer, J. (2016). Erasure without work in an
  asymmetric double-well potential. \emph{Physical Review Letters}
  117:200601. https://doi.org/10.1103/PhysRevLett.117.200601
\item
  Gavrilov, M., Chétrite, R. and Bechhoefer, J. (2017). Direct
  measurement of weakly nonequilibrium system entropy is consistent with
  Gibbs-Shannon form. \emph{PNAS} 114:11097-11102.
  https://doi.org/10.1073/pnas.1708689114
\item
  Yan, L. L., Xiong, T. P., Rehan, K., Zhou, F., Liang, D. F. et
  al.~(2018). Single-atom demonstration of the quantum Landauer
  principle. \emph{Physical Review Letters} 120:210601.
  https://doi.org/10.1103/PhysRevLett.120.210601
\item
  Dago, S., Pereda, J., Barros, N., Ciliberto, S. and Bellon, L. (2021).
  Information and thermodynamics: fast and precise approach to
  Landauer's bound in an underdamped micromechanical oscillator.
  \emph{Physical Review Letters} 126:170601.
  https://doi.org/10.1103/PhysRevLett.126.170601
\item
  Scandi, M., Barker, D., Lehmann, S., Dick, K. A., Maisi, V. F. and
  Perarnau-Llobet, M. (2022). Minimally dissipative information erasure
  in a quantum dot via thermodynamic length. \emph{Physical Review
  Letters} 129:270601. https://doi.org/10.1103/PhysRevLett.129.270601
\item
  Toyabe, S., Sagawa, T., Ueda, M., Muneyuki, E. and Sano, M. (2010).
  Experimental demonstration of information-to-energy conversion and
  validation of the generalized Jarzynski equality. \emph{Nature
  Physics} 6:988-992. https://doi.org/10.1038/nphys1821
\item
  Koski, J. V., Maisi, V. F., Pekola, J. P. and Averin, D. V. (2014).
  Experimental realization of a Szilard engine with a single electron.
  \emph{PNAS} 111:13786-13789. https://doi.org/10.1073/pnas.1406966111
\item
  Koski, J. V., Maisi, V. F., Sagawa, T. and Pekola, J. P. (2014).
  Experimental observation of the role of mutual information in the
  nonequilibrium dynamics of a Maxwell demon. \emph{Physical Review
  Letters} 113:030601. https://doi.org/10.1103/PhysRevLett.113.030601
\item
  Koski, J. V., Kutvonen, A., Khaymovich, I. M., Ala-Nissila, T. and
  Pekola, J. P. (2015). On-chip Maxwell's demon as an
  information-powered refrigerator. \emph{Physical Review Letters}
  115:260602. https://doi.org/10.1103/PhysRevLett.115.260602
\item
  Camati, P. A., Peterson, J. P. S., Batalhão, T. B., Micadei, K.,
  Souza, A. M. et al.~(2016). Experimental rectification of entropy
  production by Maxwell's demon in a quantum system. \emph{Physical
  Review Letters} 117:240502.
  https://doi.org/10.1103/PhysRevLett.117.240502
\item
  Peterson, J. P. S., Sarthour, R. S., Souza, A. M., Oliveira, I. S.,
  Goold, J. et al.~(2016). Experimental demonstration of information to
  energy conversion in a quantum system at the Landauer limit.
  \emph{Proceedings of the Royal Society A} 472:20150813.
  https://doi.org/10.1098/rspa.2015.0813
\item
  Cottet, N., Jezouin, S., Bretheau, L., Campagne-Ibarcq, P., Ficheux,
  Q. et al.~(2017). Observing a quantum Maxwell demon at work.
  \emph{PNAS} 114:7561-7564. https://doi.org/10.1073/pnas.1704827114
\item
  Masuyama, Y., Funo, K., Murashita, Y., Noguchi, A., Kono, S. et
  al.~(2018). Information-to-work conversion by Maxwell's demon in a
  superconducting circuit quantum electrodynamical system. \emph{Nature
  Communications} 9:1291. https://doi.org/10.1038/s41467-018-03686-y
\item
  Naghiloo, M., Alonso, J. J., Romito, A., Lutz, E. and Murch, K. W.
  (2018). Information gain and loss for a quantum Maxwell's demon.
  \emph{Physical Review Letters} 121:030604.
  https://doi.org/10.1103/PhysRevLett.121.030604
\item
  Kumar, A., Wu, T.-Y., Giraldo, F. and Weiss, D. S. (2018). Sorting
  ultracold atoms in a three-dimensional optical lattice in a
  realization of Maxwell's demon. \emph{Nature} 561:83-87.
  https://doi.org/10.1038/s41586-018-0458-7
\item
  Linpeng, X., Bresque, L., Maffei, M., Jordan, A. N. et al.~(2022).
  Energetic cost of measurements using quantum, coherent, and thermal
  light. \emph{Physical Review Letters} 128:220506.
  https://doi.org/10.1103/PhysRevLett.128.220506
\item
  Paneru, G., Lee, D. Y., Tlusty, T. and Pak, H. K. (2018). Lossless
  Brownian information engine. \emph{Physical Review Letters}
  120:020601. https://doi.org/10.1103/PhysRevLett.120.020601
\item
  Admon, T., Rahav, S. and Roichman, Y. (2018). Experimental realization
  of an information machine with tunable temporal correlations.
  \emph{Physical Review Letters} 121:180601.
  https://doi.org/10.1103/PhysRevLett.121.180601
\item
  Ribezzi-Crivellari, M. and Ritort, F. (2019). Large work extraction
  and the Landauer limit in a continuous Maxwell demon. \emph{Nature
  Physics} 15:660-664. https://doi.org/10.1038/s41567-019-0481-0
\item
  Debiossac, M., Grass, D., Alonso, J. J., Lutz, E. and Kiesel, N.
  (2020). Thermodynamics of continuous non-Markovian feedback control.
  \emph{Nature Communications} 11:1360.
  https://doi.org/10.1038/s41467-020-15148-5
\item
  Saha, T. K., Lucero, J. N. E., Ehrich, J., Sivak, D. A. and
  Bechhoefer, J. (2021). Maximizing power and velocity of an information
  engine. \emph{PNAS} 118:e2023356118.
  https://doi.org/10.1073/pnas.2023356118
\item
  Saha, T. K., Lucero, J. N. E., Ehrich, J., Sivak, D. A. and
  Bechhoefer, J. (2022). Bayesian information engine that optimally
  exploits noisy measurements. \emph{Physical Review Letters}
  129:130601. https://doi.org/10.1103/PhysRevLett.129.130601
\item
  Saha, T. K., Ehrich, J., Gavrilov, M., Still, S., Sivak, D. A. and
  Bechhoefer, J. (2023). Information engine in a nonequilibrium bath.
  \emph{Physical Review Letters} 131:057101.
  https://doi.org/10.1103/PhysRevLett.131.057101
\item
  Barker, D., Scandi, M., Lehmann, S., Thelander, C., Dick, K. A.,
  Perarnau-Llobet, M. and Maisi, V. F. (2022). Experimental verification
  of the work fluctuation-dissipation relation for information-to-work
  conversion. \emph{Physical Review Letters} 128:040602.
  https://doi.org/10.1103/PhysRevLett.128.040602
\item
  Yan, L.-L., Bu, J.-T., Zeng, Q., Zhang, K., Cui, K.-F. et al.~(2024).
  Experimental verification of demon-involved fluctuation theorems.
  \emph{Physical Review Letters} 133:090402.
  https://doi.org/10.1103/PhysRevLett.133.090402
\item
  Zeng, Q. and Wang, J. (2021). New fluctuation theorems on Maxwell's
  demon. \emph{Science Advances} 7:eabf1807.
  https://doi.org/10.1126/sciadv.abf1807
\item
  Demin, A. V. and Vityaev, E. E. (2014). Learning in a virtual model of
  the C. elegans nematode for locomotion and chemotaxis.
  \emph{Biologically Inspired Cognitive Architectures} 7:9-14.
  https://doi.org/10.1016/j.bica.2013.11.005
\item
  Putintsev, N. I., Isupov, O. V. and Vityaev, E. E. (2015). Adaptive
  control system for a mobile agent in a physical environment based on
  functional systems theory. \emph{Russian Journal of Genetics: Applied
  Research} 5:601-608. https://doi.org/10.1134/S2079059715060131
\item
  Rosenblueth, A., Wiener, N. and Bigelow, J. (1943). Behavior, purpose
  and teleology. \emph{Philosophy of Science} 10:18-24.
  https://doi.org/10.1086/286788
\item
  Anokhin, P. K. (1968). The functional system as a unit of organism
  integrative activity. In \emph{Systems Theory and Biology}, Springer,
  pp.~376-403. https://doi.org/10.1007/978-3-642-88343-9\_15
\item
  Anokhin, P. K. (1974). The biological roots of the conditioned reflex.
  In \emph{Biology and Neurophysiology of the Conditioned Reflex and Its
  Role in Adaptive Behavior}, Pergamon, pp.~1-24.
  https://doi.org/10.1016/B978-0-08-021516-7.50008-2
\item
  Sudakov, K. V. (1997). The theory of functional systems: general
  postulates and principles of dynamic organization. \emph{Integrative
  Physiological and Behavioral Science} 32:392-414.
  https://doi.org/10.1007/BF02688634
\item
  Shvyrkov, V. B. (1980). Goal as a system-forming factor in behavior
  and learning. In \emph{Neural Mechanisms of Goal-directed Behavior and
  Learning}, Academic Press, pp.~199-219.
  https://doi.org/10.1016/B978-0-12-688980-2.50018-7
\item
  Rosen, R. (2012). \emph{Anticipatory Systems: Philosophical,
  Mathematical, and Methodological Foundations}, 2nd ed.~Springer.
  https://doi.org/10.1007/978-1-4614-1269-4
\item
  Sterling, P. (2012). Allostasis: a model of predictive regulation.
  \emph{Physiology \& Behavior} 106:5-15.
  https://doi.org/10.1016/j.physbeh.2011.06.004
\item
  Virgo, N., Biehl, M. and McGregor, S. (2021). Interpreting dynamical
  systems as Bayesian reasoners. In \emph{Machine Learning and
  Principles and Practice of Knowledge Discovery in Databases} (ECML
  PKDD 2021 workshops), Communications in Computer and Information
  Science, Springer, pp.~726-762.
  https://doi.org/10.1007/978-3-030-93736-2\_52
\item
  Biehl, M. and Virgo, N. (2023). Interpreting systems as solving
  POMDPs: a step towards a formal understanding of agency. In
  \emph{Active Inference: IWAI 2022}, Communications in Computer and
  Information Science 1721, Springer, pp.~16-31.
  https://doi.org/10.1007/978-3-031-28719-0\_2
\item
  Maturana, H. R. and Varela, F. J. (1980). \emph{Autopoiesis and
  Cognition: The Realization of the Living}. Reidel.
  https://doi.org/10.1007/978-94-009-8947-4
\item
  Kauffman, S. (2003). Molecular autonomous agents. \emph{Philosophical
  Transactions of the Royal Society A} 361:1089-1099.
  https://doi.org/10.1098/rsta.2003.1186
\item
  Di Paolo, E. A. (2005). Autopoiesis, adaptivity, teleology, agency.
  \emph{Phenomenology and the Cognitive Sciences} 4:429-452.
  https://doi.org/10.1007/s11097-005-9002-y
\item
  Barandiaran, X. and Moreno, A. (2006). On what makes certain dynamical
  systems cognitive: a minimally cognitive organization program.
  \emph{Adaptive Behavior} 14:171-185.
  https://doi.org/10.1177/105971230601400208
\item
  Barandiaran, X. and Moreno, A. (2008). Adaptivity: from metabolism to
  behavior. \emph{Adaptive Behavior} 16:325-344.
  https://doi.org/10.1177/1059712308093868
\item
  Barandiaran, X. E., Di Paolo, E. and Rohde, M. (2009). Defining
  agency: individuality, normativity, asymmetry, and spatio-temporality
  in action. \emph{Adaptive Behavior} 17:367-386.
  https://doi.org/10.1177/1059712309343819
\item
  Ruiz-Mirazo, K. and Moreno, A. (2011). Autonomy in evolution: from
  minimal to complex life. \emph{Synthese} 185:21-52.
  https://doi.org/10.1007/s11229-011-9874-z
\item
  Moreno, A. and Mossio, M. (2015). \emph{Biological Autonomy: A
  Philosophical and Theoretical Enquiry}. Springer.
  https://doi.org/10.1007/978-94-017-9837-2
\item
  Montévil, M. and Mossio, M. (2015). Biological organisation as closure
  of constraints. \emph{Journal of Theoretical Biology} 372:179-191.
  https://doi.org/10.1016/j.jtbi.2015.02.029
\item
  Aguilera, M. and Barandiaran, X. E. (2024). Thermina: a minimal model
  of autonomous agency from the lens of stochastic thermodynamics.
  \emph{ALIFE 2024: Proceedings of the 2024 Artificial Life Conference}.
  https://doi.org/10.1162/isal\_a\_00826
\item
  Barandiaran, X. E. and Almendros, L. S. (2024). Transforming agency:
  on the mode of existence of large language models. arXiv:2407.10735
  (preprint).
\item
  Levin, M. (2022). Technological approach to mind everywhere: an
  experimentally-grounded framework for understanding diverse bodies and
  minds. \emph{Frontiers in Systems Neuroscience} 16:768201.
  https://doi.org/10.3389/fnsys.2022.768201
\item
  Friston, K. (2010). The free-energy principle: a unified brain theory?
  \emph{Nature Reviews Neuroscience} 11:127-138.
  https://doi.org/10.1038/nrn2787
\item
  Friston, K. (2013). Life as we know it. \emph{Journal of the Royal
  Society Interface} 10:20130475. https://doi.org/10.1098/rsif.2013.0475
\item
  Rovelli, C. (2020). Agency in physics. arXiv:2007.05300 (preprint).
\item
  Jaeger, J. (2023). Artificial intelligence is algorithmic mimicry: why
  artificial ``agents'' are not (and won't be) proper agents.
  arXiv:2307.07515 (preprint).
\item
  Azadi, P. (2025). Computational irreducibility as the foundation of
  agency. arXiv:2505.04646 (preprint).
\item
  Wissner-Gross, A. D. and Freer, C. E. (2013). Causal entropic forces.
  \emph{Physical Review Letters} 110:168702.
  https://doi.org/10.1103/PhysRevLett.110.168702
\item
  Kappen, H. J. (2013). Comment: causal entropic forces. arXiv:1312.4185
  (preprint).
\item
  Fields, C., Goldstein, A. and Sandved-Smith, L. (2024). Making the
  thermodynamic cost of active inference explicit. \emph{Entropy}
  26:622. https://doi.org/10.3390/e26080622
\item
  Bartlett, S., Eckford, A. W., Egbert, M., Lingam, M. et al.~(2025).
  Physics of life: exploring information as a distinctive feature of
  living systems. \emph{PRX Life} 3:037003.
  https://doi.org/10.1103/rsx4-8x5f
\item
  Funo, K., Shiraishi, N. and Saito, K. (2019). Speed limit for open
  quantum systems. \emph{New Journal of Physics} 21:013006.
  https://doi.org/10.1088/1367-2630/aaf9f5
\item
  Shiraishi, N. and Saito, K. (2021). Speed limit for open systems
  coupled to general environments. \emph{Physical Review Research}
  3:023074. https://doi.org/10.1103/PhysRevResearch.3.023074
\item
  Vo, V. T., Van Vu, T. and Hasegawa, Y. (2020). Unified approach to
  classical speed limit and thermodynamic uncertainty relation.
  \emph{Physical Review E} 102:062132.
  https://doi.org/10.1103/PhysRevE.102.062132
\item
  Falasco, G. and Esposito, M. (2020). Dissipation-time uncertainty
  relation. \emph{Physical Review Letters} 125:120604.
  https://doi.org/10.1103/PhysRevLett.125.120604
\item
  Ito, S. (2018). Stochastic thermodynamic interpretation of information
  geometry. \emph{Physical Review Letters} 121:030605.
  https://doi.org/10.1103/PhysRevLett.121.030605
\item
  Ito, S. and Dechant, A. (2020). Stochastic time evolution, information
  geometry, and the Cramér-Rao bound. \emph{Physical Review X}
  10:021056. https://doi.org/10.1103/PhysRevX.10.021056
\item
  Nakazato, M. and Ito, S. (2021). Geometrical aspects of entropy
  production in stochastic thermodynamics based on Wasserstein distance.
  \emph{Physical Review Research} 3:043093.
  https://doi.org/10.1103/PhysRevResearch.3.043093
\item
  Dechant, A., Sasa, S.-i. and Ito, S. (2022). Geometric decomposition
  of entropy production in out-of-equilibrium systems. \emph{Physical
  Review Research} 4:L012034.
  https://doi.org/10.1103/PhysRevResearch.4.L012034
\item
  Van Vu, T. and Saito, K. (2023). Thermodynamic unification of optimal
  transport: thermodynamic uncertainty relation, minimum dissipation,
  and thermodynamic speed limits. \emph{Physical Review X} 13:011013.
  https://doi.org/10.1103/PhysRevX.13.011013
\item
  Nagayama, R., Yoshimura, K. and Ito, S. (2025). Infinite variety of
  thermodynamic speed limits with general activities. \emph{Physical
  Review Research} 7:013307.
  https://doi.org/10.1103/PhysRevResearch.7.013307
\item
  Yoshimura, K. and Ito, S. (2021). Thermodynamic uncertainty relation
  and thermodynamic speed limit in deterministic chemical reaction
  networks. \emph{Physical Review Letters} 127:160601.
  https://doi.org/10.1103/PhysRevLett.127.160601
\item
  Aurell, E., Mejía-Monasterio, C. and Muratore-Ginanneschi, P. (2011).
  Optimal protocols and optimal transport in stochastic thermodynamics.
  \emph{Physical Review Letters} 106:250601.
  https://doi.org/10.1103/PhysRevLett.106.250601
\item
  Hasegawa, Y. (2023). Unifying speed limit, thermodynamic uncertainty
  relation and Heisenberg principle via bulk-boundary correspondence.
  \emph{Nature Communications} 14:2828.
  https://doi.org/10.1038/s41467-023-38074-8
\item
  Van Vu, T. and Hasegawa, Y. (2021). Geometrical bounds of the
  irreversibility in Markovian systems. \emph{Physical Review Letters}
  126:010601. https://doi.org/10.1103/PhysRevLett.126.010601
\item
  Delvenne, J.-C. and Falasco, G. (2024). Thermokinetic relations.
  \emph{Physical Review E} 109:014109.
  https://doi.org/10.1103/PhysRevE.109.014109
\item
  Seifert, U. (2017). Stochastic thermodynamics: from principles to the
  cost of precision. arXiv:1707.03759 (lecture notes).
\item
  Horowitz, J. M. and Gingrich, T. R. (2017). Proof of the finite-time
  thermodynamic uncertainty relation for steady-state currents.
  \emph{Physical Review E} 96:020103.
  https://doi.org/10.1103/PhysRevE.96.020103
\item
  Pietzonka, P., Barato, A. C. and Seifert, U. (2016). Universal bounds
  on current fluctuations. \emph{Physical Review E} 93:052145.
  https://doi.org/10.1103/PhysRevE.93.052145
\item
  Proesmans, K. and Van den Broeck, C. (2017). Discrete-time
  thermodynamic uncertainty relation. \emph{EPL} 119:20001.
  https://doi.org/10.1209/0295-5075/119/20001
\item
  Hasegawa, Y. and Van Vu, T. (2019). Fluctuation theorem uncertainty
  relation. \emph{Physical Review Letters} 123:110602.
  https://doi.org/10.1103/PhysRevLett.123.110602
\item
  Dechant, A. and Sasa, S.-i. (2018). Entropic bounds on currents in
  Langevin systems. \emph{Physical Review E} 97:062101.
  https://doi.org/10.1103/PhysRevE.97.062101
\item
  Lee, J. S., Park, J.-M. and Park, H. (2019). Thermodynamic uncertainty
  relation for underdamped Langevin systems driven by a
  velocity-dependent force. \emph{Physical Review E} 100:062132.
  https://doi.org/10.1103/PhysRevE.100.062132
\item
  Lee, J. S., Park, J.-M. and Park, H. (2021). Universal form of
  thermodynamic uncertainty relation for Langevin dynamics.
  \emph{Physical Review E} 104:L052102.
  https://doi.org/10.1103/PhysRevE.104.L052102
\item
  Dechant, A. and Sasa, S.-i. (2021). Improving thermodynamic bounds
  using correlations. \emph{Physical Review X} 11:041061.
  https://doi.org/10.1103/PhysRevX.11.041061
\item
  Horowitz, J. M. and Gingrich, T. R. (2020). Thermodynamic uncertainty
  relations constrain non-equilibrium fluctuations. \emph{Nature
  Physics} 16:15-20. https://doi.org/10.1038/s41567-019-0702-6
\item
  Barato, A. C. and Seifert, U. (2016). Cost and precision of Brownian
  clocks. \emph{Physical Review X} 6:041053.
  https://doi.org/10.1103/PhysRevX.6.041053
\item
  Salazar, D. S. P. (2021). Information bound for entropy production
  from the detailed fluctuation theorem. \emph{Physical Review E}
  103:022122. https://doi.org/10.1103/PhysRevE.103.022122
\item
  Landi, G. T. and Paternostro, M. (2021). Irreversible entropy
  production: from classical to quantum. \emph{Reviews of Modern
  Physics} 93:035008. https://doi.org/10.1103/RevModPhys.93.035008
\item
  Van Vu, T. and Hasegawa, Y. (2020). Thermodynamic uncertainty
  relations under arbitrary control protocols. \emph{Physical Review
  Research} 2:013060. https://doi.org/10.1103/PhysRevResearch.2.013060
\item
  Tanogami, T., Van Vu, T. and Saito, K. (2023). Universal bounds on the
  performance of information-thermodynamic engine. \emph{Physical Review
  Research} 5:043280. https://doi.org/10.1103/PhysRevResearch.5.043280
\item
  Kumasaki, K., Tojo, K., Sagawa, T. and Funo, K. (2026). Thermodynamic
  uncertainty relation for feedback cooling. \emph{Physical Review E}
  113:024134. https://doi.org/10.1103/4pq6-7djm
\item
  Honma, R. and Van Vu, T. (2026). Thermodynamic uncertainty relation
  with quantum feedback. \emph{Physical Review Letters}.
  https://doi.org/10.1103/6d8l-94pv. Preprint: arXiv:2602.22651.
\item
  Gingrich, T. R. and Horowitz, J. M. (2017). Fundamental bounds on
  first passage time fluctuations for currents. \emph{Physical Review
  Letters} 119:170601. https://doi.org/10.1103/PhysRevLett.119.170601
\item
  Schmiedl, T. and Seifert, U. (2007). Optimal finite-time processes in
  stochastic thermodynamics. \emph{Physical Review Letters} 98:108301.
  https://doi.org/10.1103/PhysRevLett.98.108301
\item
  Proesmans, K., Ehrich, J. and Bechhoefer, J. (2020). Optimal
  finite-time bit erasure under full control. \emph{Physical Review E}
  102:032105. https://doi.org/10.1103/PhysRevE.102.032105
\item
  Zhen, Y.-Z., Egloff, D., Modi, K. and Dahlsten, O. (2021). Universal
  bound on energy cost of bit reset in finite time. \emph{Physical
  Review Letters} 127:190602.
  https://doi.org/10.1103/PhysRevLett.127.190602
\item
  Esposito, M., Kawai, R., Lindenberg, K. and Van den Broeck, C. (2010).
  Efficiency at maximum power of low-dissipation Carnot engines.
  \emph{Physical Review Letters} 105:150603.
  https://doi.org/10.1103/PhysRevLett.105.150603
\item
  Tu, Z. C. (2008). Efficiency at maximum power of Feynman's ratchet as
  a heat engine. \emph{Journal of Physics A} 41:312003.
  https://doi.org/10.1088/1751-8113/41/31/312003
\item
  Shiraishi, N., Saito, K. and Tasaki, H. (2016). Universal trade-off
  relation between power and efficiency for heat engines. \emph{Physical
  Review Letters} 117:190601.
  https://doi.org/10.1103/PhysRevLett.117.190601
\item
  Pietzonka, P. and Seifert, U. (2018). Universal trade-off between
  power, efficiency, and constancy in steady-state heat engines.
  \emph{Physical Review Letters} 120:190602.
  https://doi.org/10.1103/PhysRevLett.120.190602
\item
  Ma, Y.-H., Xu, D., Dong, H. and Sun, C.-P. (2018). Universal
  constraint for efficiency and power of a low-dissipation heat engine.
  \emph{Physical Review E} 98:042112.
  https://doi.org/10.1103/PhysRevE.98.042112
\item
  Tsirlin, A. M., Mironova, V. A., Amelkin, S. A. and Kazakov, V.
  (1998). Finite-time thermodynamics: conditions of minimal dissipation
  for thermodynamic processes with given rate. \emph{Physical Review E}
  58:215-223. https://doi.org/10.1103/PhysRevE.58.215
\item
  Lee, J. S., Lee, S., Kwon, H. and Park, H. (2022). Speed limit for a
  highly irreversible process and tight finite-time Landauer's bound.
  \emph{Physical Review Letters} 129:120603.
  https://doi.org/10.1103/PhysRevLett.129.120603
\item
  Van Vu, T. and Saito, K. (2022). Finite-time quantum Landauer
  principle and quantum coherence. \emph{Physical Review Letters}
  128:010602. https://doi.org/10.1103/PhysRevLett.128.010602
\item
  Rolandi, A. and Perarnau-Llobet, M. (2023). Finite-time Landauer
  principle beyond weak coupling. \emph{Quantum} 7:1161.
  https://doi.org/10.22331/q-2023-11-03-1161
\item
  Kamijima, T., Funo, K. and Sagawa, T. (2024). Finite-time
  thermodynamic bounds and tradeoff relations for information
  processing. arXiv:2409.08606 (preprint).
\item
  Tkachenko, A. V. (2025). Thermodynamic cost of inference and learning
  in physical neural networks. arXiv:2503.09980 (preprint). Title of the
  current arXiv version.
\item
  Tatikonda, S. and Mitter, S. (2004). Control under communication
  constraints. \emph{IEEE Transactions on Automatic Control}
  49:1056-1068. https://doi.org/10.1109/TAC.2004.831187
\item
  Nair, G. N. and Evans, R. J. (2004). Stabilizability of stochastic
  linear systems with finite feedback data rates. \emph{SIAM Journal on
  Control and Optimization} 43:413-436.
  https://doi.org/10.1137/S0363012902402116
\item
  Nair, G. N., Fagnani, F., Zampieri, S. and Evans, R. J. (2007).
  Feedback control under data rate constraints: an overview.
  \emph{Proceedings of the IEEE} 95:108-137.
  https://doi.org/10.1109/JPROC.2006.887294
\item
  Sandberg, H., Delvenne, J.-C., Newton, N. J. and Mitter, S. K. (2014).
  Maximum work extraction and implementation costs for nonequilibrium
  Maxwell's demons. \emph{Physical Review E} 90:042119.
  https://doi.org/10.1103/PhysRevE.90.042119
\item
  Horowitz, J. M. and Sandberg, H. (2014). Second-law-like inequalities
  with information and their interpretations. \emph{New Journal of
  Physics} 16:125007. https://doi.org/10.1088/1367-2630/16/12/125007
\item
  Jacobs, K. (2009). Second law of thermodynamics and quantum feedback
  control: Maxwell's demon with weak measurements. \emph{Physical Review
  A} 80:012322. https://doi.org/10.1103/PhysRevA.80.012322
\item
  Funo, K., Watanabe, Y. and Ueda, M. (2013). Integral quantum
  fluctuation theorems under measurement and feedback control.
  \emph{Physical Review E} 88:052121.
  https://doi.org/10.1103/PhysRevE.88.052121
\item
  Guryanova, Y., Friis, N. and Huber, M. (2020). Ideal projective
  measurements have infinite resource costs. \emph{Quantum} 4:222.
  https://doi.org/10.22331/q-2020-01-13-222
\item
  Danageozian, A., Wilde, M. M. and Buscemi, F. (2022). Thermodynamic
  constraints on quantum information gain and error correction: a triple
  trade-off. \emph{PRX Quantum} 3:020318.
  https://doi.org/10.1103/PRXQuantum.3.020318
\item
  Wolpert, D. H. (2020). Minimal entropy production rate of interacting
  systems. \emph{New Journal of Physics} 22:113013.
  https://doi.org/10.1088/1367-2630/abc5c6
\item
  Tasnim, F. and Wolpert, D. H. (2023). Stochastic thermodynamics of
  multiple co-evolving systems: beyond multipartite processes.
  \emph{Entropy} 25:1078. https://doi.org/10.3390/e25071078
\item
  Wolpert, D. H. (2020). Uncertainty relations and fluctuation theorems
  for Bayes nets. \emph{Physical Review Letters} 125:200602.
  https://doi.org/10.1103/PhysRevLett.125.200602
\item
  Rolandi, A., Abiuso, P. and Perarnau-Llobet, M. (2023). Collective
  advantages in finite-time thermodynamics. \emph{Physical Review
  Letters} 131:210401. https://doi.org/10.1103/PhysRevLett.131.210401
\item
  Fujimoto, Y. and Ito, S. (2024). Game-theoretical approach to minimum
  entropy productions in information thermodynamics. \emph{Physical
  Review Research} 6:013023.
  https://doi.org/10.1103/PhysRevResearch.6.013023
\item
  Lloyd, S. (2000). Ultimate physical limits to computation.
  \emph{Nature} 406:1047-1054. https://doi.org/10.1038/35023282
\item
  Deffner, S. and Campbell, S. (2017). Quantum speed limits: from
  Heisenberg's uncertainty principle to optimal quantum control.
  \emph{Journal of Physics A} 50:453001.
  https://doi.org/10.1088/1751-8121/aa86c6
\item
  Lan, G., Sartori, P., Neumann, S., Sourjik, V. and Tu, Y. (2012). The
  energy-speed-accuracy trade-off in sensory adaptation. \emph{Nature
  Physics} 8:422-428. https://doi.org/10.1038/nphys2276
\item
  Mehta, P. and Schwab, D. J. (2012). Energetic costs of cellular
  computation. \emph{PNAS} 109:17978-17982.
  https://doi.org/10.1073/pnas.1207814109
\item
  Lang, A. H., Fisher, C. K., Mora, T. and Mehta, P. (2014).
  Thermodynamics of statistical inference by cells. \emph{Physical
  Review Letters} 113:148103.
  https://doi.org/10.1103/PhysRevLett.113.148103
\item
  Govern, C. C. and ten Wolde, P. R. (2014). Optimal resource allocation
  in cellular sensing systems. \emph{PNAS} 111:17486-17491.
  https://doi.org/10.1073/pnas.1411524111
\item
  Sartori, P., Granger, L., Lee, C. F. and Horowitz, J. M. (2014).
  Thermodynamic costs of information processing in sensory adaptation.
  \emph{PLoS Computational Biology} 10:e1003974.
  https://doi.org/10.1371/journal.pcbi.1003974
\item
  Barato, A. C., Hartich, D. and Seifert, U. (2013).
  Information-theoretic versus thermodynamic entropy production in
  autonomous sensory networks. \emph{Physical Review E} 87:042104.
  https://doi.org/10.1103/PhysRevE.87.042104
\item
  Barato, A. C., Hartich, D. and Seifert, U. (2014). Efficiency of
  cellular information processing. \emph{New Journal of Physics}
  16:103024. https://doi.org/10.1088/1367-2630/16/10/103024
\item
  Ouldridge, T. E., Govern, C. C. and ten Wolde, P. R. (2017).
  Thermodynamics of computational copying in biochemical systems.
  \emph{Physical Review X} 7:021004.
  https://doi.org/10.1103/PhysRevX.7.021004
\item
  Cao, Y., Wang, H., Ouyang, Q. and Tu, Y. (2015). The free-energy cost
  of accurate biochemical oscillations. \emph{Nature Physics}
  11:772-778. https://doi.org/10.1038/nphys3412
\item
  Zhang, D., Cao, Y., Ouyang, Q. and Tu, Y. (2019). The energy cost and
  optimal design for synchronization of coupled molecular oscillators.
  \emph{Nature Physics} 16:95-100.
  https://doi.org/10.1038/s41567-019-0701-7
\item
  Sartori, P. and Tu, Y. (2015). Free energy cost of reducing noise
  while maintaining a high sensitivity. \emph{Physical Review Letters}
  115:118102. https://doi.org/10.1103/PhysRevLett.115.118102
\item
  Horowitz, J. M., Zhou, K. and England, J. L. (2017). Minimum energetic
  cost to maintain a target nonequilibrium state. \emph{Physical Review
  E} 95:042102. https://doi.org/10.1103/PhysRevE.95.042102
\item
  Bryant, S. J. and Machta, B. B. (2023). Physical constraints in
  intracellular signaling: the cost of sending a bit. \emph{Physical
  Review Letters} 131:068401.
  https://doi.org/10.1103/PhysRevLett.131.068401
\item
  Tjalma, A. J., Galstyan, V., Goedhart, J., Slim, L. et al.~(2023).
  Trade-offs between cost and information in cellular prediction.
  \emph{Proceedings of the National Academy of Sciences}
  120:e2303078120. https://doi.org/10.1073/pnas.2303078120
\item
  Schreiber, T. (2000). Measuring information transfer. \emph{Physical
  Review Letters} 85:461-464. https://doi.org/10.1103/PhysRevLett.85.461
\item
  Tishby, N. and Polani, D. (2011). Information theory of decisions and
  actions. In \emph{Perception-Action Cycle}, Springer, pp.~601-636.
  https://doi.org/10.1007/978-1-4419-1452-1\_19
\item
  Stratonovich, R. L. (2020). \emph{Theory of Information and its Value}
  (English edition of the 1975 Russian monograph). Springer.
  https://doi.org/10.1007/978-3-030-22833-0
\item
  Pal, S., Saryal, S., Segal, D., Mahesh, T. S. and Agarwalla, B. K.
  (2020). Experimental study of the thermodynamic uncertainty relation.
  \emph{Physical Review Research} 2:022044.
  https://doi.org/10.1103/PhysRevResearch.2.022044
\item
  Barker, D., Lehmann, S., Dick, K. A., Samuelsson, P. et al.~(2025).
  Information thermodynamics in a quantum dot Szilard engine:
  experimentally investigating fluctuation theorems and thermodynamic
  uncertainty relations. arXiv:2511.08541 (preprint).
\item
  Friedman, H. M., Agarwalla, B. K., Shein-Lumbroso, O., Tal, O. and
  Segal, D. (2020). Thermodynamic uncertainty relation in atomic-scale
  quantum conductors. \emph{Physical Review B} 101:195423.
  https://doi.org/10.1103/PhysRevB.101.195423
\item
  Hwang, W. and Hyeon, C. (2018). Energetic costs, precision, and
  transport efficiency of molecular motors. \emph{Journal of Physical
  Chemistry Letters} 9:513-520.
  https://doi.org/10.1021/acs.jpclett.7b03197
\item
  Li, J., Horowitz, J. M., Gingrich, T. R. and Fakhri, N. (2019).
  Quantifying dissipation using fluctuating currents. \emph{Nature
  Communications} 10:1666. https://doi.org/10.1038/s41467-019-09631-x
\item
  Ness, G., Lam, M. R., Alt, W., Meschede, D., Sagi, Y. and Alberti, A.
  (2021). Observing crossover between quantum speed limits.
  \emph{Science Advances} 7:eabj9119.
  https://doi.org/10.1126/sciadv.abj9119
\item
  Pires, L. B., Goerlich, R., Luna da Fonseca, A., Debiossac, M.,
  Hervieux, P.-A. et al.~(2023). Optimal time-entropy bounds and speed
  limits for Brownian thermal shortcuts. \emph{Physical Review Letters}
  131:097101. https://doi.org/10.1103/PhysRevLett.131.097101
\item
  Oikawa, S., Nakayama, Y., Ito, S., Sagawa, T. and Toyabe, S. (2025).
  Experimentally achieving minimal dissipation via thermodynamically
  optimal transport. \emph{Nature Communications} 16:10424.
  https://doi.org/10.1038/s41467-025-66519-9
\item
  Proesmans, K., Dreher, Y., Gavrilov, M., Bechhoefer, J. and Van den
  Broeck, C. (2016). Brownian duet: a novel tale of thermodynamic
  efficiency. \emph{Physical Review X} 6:041010.
  https://doi.org/10.1103/PhysRevX.6.041010
\item
  Mancino, L., Cavina, V., De Pasquale, A., Sbroscia, M., Booth, R. I.
  et al.~(2018). Geometrical bounds on irreversibility in open quantum
  systems. \emph{Physical Review Letters} 121:160602.
  https://doi.org/10.1103/PhysRevLett.121.160602
\item
  Dago, S. and Bellon, L. (2022). Dynamics of information erasure and
  extension of Landauer's bound to fast processes. \emph{Physical Review
  Letters} 128:070604. https://doi.org/10.1103/PhysRevLett.128.070604
\item
  Mattingly, H. H., Kamino, K., Machta, B. B. and Emonet, T. (2021).
  \emph{Escherichia coli} chemotaxis is information limited.
  \emph{Nature Physics} 17:1426-1431.
  https://doi.org/10.1038/s41567-021-01380-3
\item
  Blickle, V. and Bechinger, C. (2012). Realization of a
  micrometre-sized stochastic heat engine. \emph{Nature Physics}
  8:143-146. https://doi.org/10.1038/nphys2163
\item
  Martínez, I. A., Roldán, É., Dinis, L., Petrov, D., Parrondo, J. M. R.
  and Rica, R. A. (2016). Brownian Carnot engine. \emph{Nature Physics}
  12:67-70. https://doi.org/10.1038/nphys3518
\item
  Roldán, É., Martínez, I. A., Parrondo, J. M. R. and Petrov, D. (2014).
  Universal features in the energetics of symmetry breaking.
  \emph{Nature Physics} 10:457-461. https://doi.org/10.1038/nphys2940
\item
  Krishnamurthy, S., Ghosh, S., Chatterji, D., Ganapathy, R. and Sood,
  A. K. (2016). A micrometre-sized heat engine operating between
  bacterial reservoirs. \emph{Nature Physics} 12:1134-1138.
  https://doi.org/10.1038/nphys3870
\item
  Krishnamurthy, S., Ganapathy, R. and Sood, A. K. (2023). Overcoming
  power-efficiency tradeoff in a micro heat engine by engineered
  system-bath interactions. \emph{Nature Communications} 14:6842.
  https://doi.org/10.1038/s41467-023-42350-y
\item
  Paneru, G., Dutta, S., Sagawa, T., Tlusty, T. and Pak, H. K. (2020).
  Efficiency fluctuations and noise induced refrigerator-to-heater
  transition in information engines. \emph{Nature Communications}
  11:1012. https://doi.org/10.1038/s41467-020-14823-x
\item
  Jung, T., Polani, D. and Stone, P. (2011). Empowerment for continuous
  agent-environment systems. \emph{Adaptive Behavior} 19:16-39.
  https://doi.org/10.1177/1059712310392389
\item
  Salge, C., Glackin, C. and Polani, D. (2013). Empowerment: an
  introduction. arXiv:1310.1863 (preprint).
\item
  Kolchinsky, A. (2024). Thermodynamic dissipation does not bound
  replicator growth and decay rates. arXiv:2404.01130 (preprint).
\item
  Adlam, E. C., McQueen, K. J. and Waegell, M. (2025). Agency cannot be
  a purely quantum phenomenon. arXiv:2510.13247 (preprint).
\item
  Goldt, S. and Seifert, U. (2017). Stochastic thermodynamics of
  learning. \emph{Physical Review Letters} 118:010601.
  https://doi.org/10.1103/PhysRevLett.118.010601
\item
  Still, S. (2020). Thermodynamic cost and benefit of memory.
  \emph{Physical Review Letters} 124:050601.
  https://doi.org/10.1103/PhysRevLett.124.050601
\item
  Boyd, A. B., Crutchfield, J. P. and Gu, M. (2022). Thermodynamic
  machine learning through maximum work production. \emph{New Journal of
  Physics} 24:083040. https://doi.org/10.1088/1367-2630/ac4309
\item
  Boyd, A. B., Crutchfield, J. P., Gu, M. and Binder, F. C. (2025).
  Thermodynamic overfitting and generalization: energetics of predictive
  intelligence. \emph{New Journal of Physics} 27:063901.
  https://doi.org/10.1088/1367-2630/addf71
\item
  Ehrich, J., Still, S. and Sivak, D. A. (2023). Energetic cost of
  feedback control. \emph{Physical Review Research} 5:023080.
  https://doi.org/10.1103/PhysRevResearch.5.023080
\item
  Ehrich, J. and Sivak, D. A. (2023). Energy and information flows in
  autonomous systems. \emph{Frontiers in Physics} 11:1108357.
  https://doi.org/10.3389/fphy.2023.1108357
\item
  Wolpert, D. H. (2019). The stochastic thermodynamics of computation.
  \emph{Journal of Physics A} 52:193001.
  https://doi.org/10.1088/1751-8121/ab0850
\item
  Levy, W. B. and Baxter, R. A. (1996). Energy efficient neural codes.
  \emph{Neural Computation} 8:531-543.
  https://doi.org/10.1162/neco.1996.8.3.531
\item
  Balasubramanian, V., Kimber, D. and Berry, M. J., II (2001).
  Metabolically efficient information processing. \emph{Neural
  Computation} 13:799-815. https://doi.org/10.1162/089976601300014358
\item
  Goldt, S. and Seifert, U. (2017). Thermodynamic efficiency of learning
  a rule in neural networks. \emph{New Journal of Physics} 19:113001.
  https://doi.org/10.1088/1367-2630/aa89ff
\item
  Fiderer, L. J., Barth, P. C., Smith, I. D. and Briegel, H. J. (2025).
  Information thermodynamics of agents: the work capacity of channels
  with memory. arXiv:2504.06209 (preprint).
\item
  Elliott, T. J., Gu, M., Garner, A. J. P. and Thompson, J. (2022).
  Quantum adaptive agents with efficient long-term memories.
  \emph{Physical Review X} 12:011007.
  https://doi.org/10.1103/PhysRevX.12.011007
\item
  Russell, S. J. and Subramanian, D. (1995). Provably bounded-optimal
  agents. \emph{Journal of Artificial Intelligence Research} 2:575-609.
  https://doi.org/10.1613/jair.133
\item
  Genewein, T., Leibfried, F., Grau-Moya, J. and Braun, D. A. (2015).
  Bounded rationality, abstraction, and hierarchical decision-making: an
  information-theoretic optimality principle. \emph{Frontiers in
  Robotics and AI} 2:27. https://doi.org/10.3389/frobt.2015.00027
\item
  Legg, S. and Hutter, M. (2007). Universal intelligence: a definition
  of machine intelligence. \emph{Minds and Machines} 17:391-444.
  https://doi.org/10.1007/s11023-007-9079-x
\item
  Takahashi, K. and Hayashi, Y. (2026). Thermodynamic limits of physical
  intelligence. In \emph{Artificial General Intelligence: 19th
  International Conference, AGI 2026, Proceedings, Part II}, Lecture
  Notes in Computer Science. Springer.
  https://doi.org/10.1007/978-3-032-33195-3\_24. arXiv:2602.05463.
\item
  Ortega, P. A., Braun, D. A., Dyer, J., Kim, K.-E. and Tishby, N.
  (2015). Information-theoretic bounded rationality. arXiv:1512.06789
  (preprint).
\item
  Zénon, A., Solopchuk, O. and Pezzulo, G. (2019). An
  information-theoretic perspective on the costs of cognition.
  \emph{Neuropsychologia} 123:5-18.
  https://doi.org/10.1016/j.neuropsychologia.2018.09.013
\item
  Sims, C. R. (2016). Rate-distortion theory and human perception.
  \emph{Cognition} 152:181-198.
  https://doi.org/10.1016/j.cognition.2016.03.020
\item
  Laughlin, S. B., de Ruyter van Steveninck, R. R. and Anderson, J. C.
  (1998). The metabolic cost of neural information. \emph{Nature
  Neuroscience} 1:36-41. https://doi.org/10.1038/236
\item
  Attwell, D. and Laughlin, S. B. (2001). An energy budget for signaling
  in the grey matter of the brain. \emph{Journal of Cerebral Blood Flow
  and Metabolism} 21:1133-1145.
  https://doi.org/10.1097/00004647-200110000-00001
\item
  Lennie, P. (2003). The cost of cortical computation. \emph{Current
  Biology} 13:493-497. https://doi.org/10.1016/S0960-9822(03)00135-0
\item
  Sengupta, B., Stemmler, M., Laughlin, S. B. and Niven, J. E. (2010).
  Action potential energy efficiency varies among neuron types in
  vertebrates and invertebrates. \emph{PLoS Computational Biology}
  6:e1000840. https://doi.org/10.1371/journal.pcbi.1000840
\item
  Sengupta, B., Stemmler, M. B. and Friston, K. J. (2013). Information
  and efficiency in the nervous system: a synthesis. \emph{PLoS
  Computational Biology} 9:e1003157.
  https://doi.org/10.1371/journal.pcbi.1003157
\item
  Levy, W. B. and Calvert, V. G. (2021). Communication consumes 35 times
  more energy than computation in the human cortex, but both costs are
  needed to predict synapse number. \emph{Proceedings of the National
  Academy of Sciences} 118:e2008173118.
  https://doi.org/10.1073/pnas.2008173118
\item
  Harris, J. J., Jolivet, R., Engl, E. and Attwell, D. (2015).
  Energy-efficient information transfer by visual pathway synapses.
  \emph{Current Biology} 25:3151-3160.
  https://doi.org/10.1016/j.cub.2015.10.063
\item
  Padamsey, Z., Katsanevaki, D., Dupuy, N. and Rochefort, N. L. (2022).
  Neocortex saves energy by reducing coding precision during food
  scarcity. \emph{Neuron} 110:280-296.e10.
  https://doi.org/10.1016/j.neuron.2021.10.024
\item
  Plaçais, P.-Y. and Preat, T. (2013). To favor survival under food
  shortage, the brain disables costly memory. \emph{Science}
  339:440-442. https://doi.org/10.1126/science.1226018
\item
  Plaçais, P.-Y., de Tredern, É., Scheunemann, L., Trannoy, S. et
  al.~(2017). Upregulated energy metabolism in the \emph{Drosophila}
  mushroom body is the trigger for long-term memory. \emph{Nature
  Communications} 8:15510. https://doi.org/10.1038/ncomms15510
\item
  Mery, F. and Kawecki, T. J. (2005). A cost of long-term memory in
  \emph{Drosophila}. \emph{Science} 308:1148.
  https://doi.org/10.1126/science.1111331
\item
  Hechler, A., de Lange, F. P. and Riedl, V. (2023). The energy
  metabolic footprint of predictive processing in the human brain.
  bioRxiv preprint, version 2 posted 2024.
  https://doi.org/10.1101/2023.12.08.570804
\item
  Malkin, J., O'Donnell, C., Houghton, C. J. and Aitchison, L. (2024).
  Signatures of Bayesian inference emerge from energy-efficient
  synapses. \emph{eLife} 12:RP92595. https://doi.org/10.7554/eLife.92595
\item
  Li, H. L. and van Rossum, M. C. W. (2020). Energy efficient synaptic
  plasticity. \emph{eLife} 9:e50804. https://doi.org/10.7554/eLife.50804
\item
  Lynn, C. W., Cornblath, E. J., Papadopoulos, L., Bertolero, M. A. and
  Bassett, D. S. (2021). Broken detailed balance and entropy production
  in the human brain. \emph{Proceedings of the National Academy of
  Sciences} 118:e2109889118. https://doi.org/10.1073/pnas.2109889118
\item
  Horowitz, M. (2014). Computing's energy problem (and what we can do
  about it). \emph{2014 IEEE International Solid-State Circuits
  Conference}, pp.~10-14. https://doi.org/10.1109/ISSCC.2014.6757323
\item
  Saad-Falcon, J., Narayan, A., Akengin, H. O., Griffin, J. W. et
  al.~(2025). Intelligence per watt: measuring intelligence efficiency
  of local AI. arXiv:2511.07885 (preprint).
\item
  Tschand, A., Rajan, A. T. R., Idgunji, S., Ghosh, A. et al.~(2025).
  MLPerf Power: benchmarking the energy efficiency of machine learning
  systems from μWatts to MWatts for sustainable AI. \emph{2025 IEEE
  International Symposium on High Performance Computer Architecture
  (HPCA)}, 1201-1216. https://doi.org/10.1109/HPCA61900.2025.00092.
  arXiv:2410.12032.
\item
  Chung, J.-W., Ma, J. J., Wu, R., Liu, J. et al.~(2025). The ML.ENERGY
  benchmark: toward automated inference energy measurement and
  optimization. arXiv:2505.06371 (preprint).
\item
  Niu, C., Zhang, W., Li, J., Zhao, Y. et al.~(2025). TokenPowerBench:
  benchmarking the power consumption of LLM inference. arXiv:2512.03024
  (preprint).
\item
  Jin, Y., Wei, G.-Y. and Brooks, D. (2025). The energy cost of
  reasoning: analyzing energy usage in LLMs with test-time compute.
  arXiv:2505.14733 (preprint).
\item
  Luccioni, S., Jernite, Y. and Strubell, E. (2024). Power hungry
  processing: watts driving the cost of AI deployment? \emph{ACM FAccT
  2024}, pp.~85-99. https://doi.org/10.1145/3630106.3658542
\item
  Samsi, S., Zhao, D., McDonald, J., Li, B. et al.~(2023). From words to
  watts: benchmarking the energy costs of large language model
  inference. \emph{2023 IEEE High Performance Extreme Computing
  Conference (HPEC)}, pp.~1-9.
  https://doi.org/10.1109/HPEC58863.2023.10363447
\item
  Dillavou, S., Stern, M., Liu, A. J. and Durian, D. J. (2022).
  Demonstration of decentralized physics-driven learning. \emph{Physical
  Review Applied} 18:014040.
  https://doi.org/10.1103/PhysRevApplied.18.014040
\item
  Stern, M., Dillavou, S., Jayaraman, D., Durian, D. J. and Liu, A. J.
  (2024). Training self-learning circuits for power-efficient solutions.
  \emph{APL Machine Learning} 2:016114.
  https://doi.org/10.1063/5.0181382
\item
  Saggio, V., Asenbeck, B. E., Hamann, A., Strömberg, T. et al.~(2021).
  Experimental quantum speed-up in reinforcement learning agents.
  \emph{Nature} 591:229-233. https://doi.org/10.1038/s41586-021-03242-7
\item
  Crosato, E., Spinney, R. E., Nigmatullin, R., Lizier, J. T. and
  Prokopenko, M. (2018). Thermodynamics and computation during
  collective motion near criticality. \emph{Physical Review E}
  97:012120. https://doi.org/10.1103/PhysRevE.97.012120
\item
  Chen, Q. and Prokopenko, M. (2025). Why collective behaviours
  self-organize to criticality: a primer on information-theoretic and
  thermodynamic utility measures. \emph{Royal Society Open Science}
  12:241655. https://doi.org/10.1098/rsos.241655
\item
  Khan, K. N., Hirki, M., Niemi, T., Nurminen, J. K. and Ou, Z. (2018).
  RAPL in action: experiences in using RAPL for power measurements.
  \emph{ACM Transactions on Modeling and Performance Evaluation of
  Computing Systems} 3:1-26. https://doi.org/10.1145/3177754
\item
  Desrochers, S., Paradis, C. and Weaver, V. M. (2016). A validation of
  DRAM RAPL power measurements. \emph{Proceedings of the Second
  International Symposium on Memory Systems}, pp.~455-470.
  https://doi.org/10.1145/2989081.2989088
\item
  Hackenberg, D., Schöne, R., Ilsche, T., Molka, D., Schuchart, J. and
  Geyer, R. (2015). An energy efficiency feature survey of the Intel
  Haswell processor. \emph{2015 IEEE IPDPS Workshops}, pp.~896-904.
  https://doi.org/10.1109/IPDPSW.2015.70
\item
  Apple. powermetrics(1) manual page. macOS developer tools
  documentation.
\item
  Hutter Prize for Lossless Compression of Human Knowledge.
  http://prize.hutter1.net/
\item
  Strubell, E., Ganesh, A. and McCallum, A. (2019). Energy and policy
  considerations for deep learning in NLP. arXiv:1906.02243 (preprint).
\item
  Patterson, D., Gonzalez, J., Le, Q., Liang, C. et al.~(2021). Carbon
  emissions and large neural network training. arXiv:2104.10350
  (preprint).
\item
  Lacoste, A., Luccioni, A., Schmidt, V. and Dandres, T. (2019).
  Quantifying the carbon emissions of machine learning. arXiv:1910.09700
  (preprint).
\item
  Luccioni, A. S., Viguier, S. and Ligozat, A.-L. (2022). Estimating the
  carbon footprint of BLOOM, a 176B parameter language model.
  arXiv:2211.02001 (preprint).
\item
  Anthony, L. F. W., Kanding, B. and Selvan, R. (2020). Carbontracker:
  tracking and predicting the carbon footprint of training deep learning
  models. arXiv:2007.03051 (preprint).
\item
  Elsworth, C., Huang, K., Patterson, D., Schneider, I. et al.~(2025).
  Measuring the environmental impact of delivering AI at Google scale.
  arXiv:2508.15734 (preprint).
\item
  Schmidhuber, J. (2010). Formal theory of creativity, fun, and
  intrinsic motivation (1990-2010). \emph{IEEE Transactions on
  Autonomous Mental Development} 2:230-247.
  https://doi.org/10.1109/TAMD.2010.2056368
\item
  Chollet, F. (2019). On the measure of intelligence. arXiv:1911.01547
  (preprint).
\item
  ARC Prize Foundation. ARC-AGI leaderboard.
  https://arcprize.org/leaderboard (accessed 8 October 2026).
\item
  Hernández-Orallo, J. (2017). \emph{The Measure of All Minds:
  Evaluating Natural and Artificial Intelligence}. Cambridge University
  Press. https://doi.org/10.1017/9781316594179
\item
  Hernández-Orallo, J. and Dowe, D. L. (2010). Measuring universal
  intelligence: towards an anytime intelligence test. \emph{Artificial
  Intelligence} 174:1508-1539.
  https://doi.org/10.1016/j.artint.2010.09.006
\item
  Gershman, S. J., Horvitz, E. J. and Tenenbaum, J. B. (2015).
  Computational rationality: a converging paradigm for intelligence in
  brains, minds, and machines. \emph{Science} 349:273-278.
  https://doi.org/10.1126/science.aac6076
\item
  Lieder, F. and Griffiths, T. L. (2020). Resource-rational analysis:
  understanding human cognition as the optimal use of limited
  computational resources. \emph{Behavioral and Brain Sciences} 43:e1.
  https://doi.org/10.1017/S0140525X1900061X
\item
  Balasubramanian, V. (2021). Brain power. \emph{PNAS} 118:e2107022118.
  https://doi.org/10.1073/pnas.2107022118
\item
  Silver, D., Singh, S., Precup, D. and Sutton, R. S. (2021). Reward is
  enough. \emph{Artificial Intelligence} 299:103535.
  https://doi.org/10.1016/j.artint.2021.103535
\item
  Vamplew, P., Smith, B. J., Källström, J., Ramos, G., Rădulescu, R. et
  al.~(2022). Scalar reward is not enough: a response to Silver, Singh,
  Precup and Sutton (2021). \emph{Autonomous Agents and Multi-Agent
  Systems} 36:41. https://doi.org/10.1007/s10458-022-09575-5
\item
  Hafez, W., Wei, C., Pena, R., Nazeri, A. et al.~(2026). A mathematical
  theory of agency and intelligence. arXiv:2602.22519 (preprint).
\item
  Karagoz, A. (2025). Energentic intelligence: from self-sustaining
  systems to enduring artificial life. arXiv:2506.04916 (preprint).
\item
  Chaisson, E. J. (2011). Energy rate density as a complexity metric and
  evolutionary driver. \emph{Complexity} 16:27-40.
  https://doi.org/10.1002/cplx.20323
\item
  Poplavskii, R. P. (1975). Thermodynamic models of information
  processes. \emph{Uspekhi Fizicheskikh Nauk} 115:465 (in Russian),
  https://doi.org/10.3367/UFNr.0115.197503d.0465. English translation:
  \emph{Soviet Physics Uspekhi} 18:222-241,
  https://doi.org/10.1070/PU1975v018n03ABEH001955
\item
  Ito, S. and Sagawa, T. (2017). Information thermodynamics on networks
  and its application to biological information processing.
  \emph{Butsuri} (日本物理学会誌) 72(9):658 (in Japanese).
  https://doi.org/10.11316/butsuri.72.9\_658
\item
  Sun, C.-P. and Quan, H.-T. (2013). Maxwell's demon and the physical
  limits on information processing. \emph{Wuli} (物理, Physics) 42(11)
  (in Chinese).
  \href{https://doi.org/10.7693/wl20131101}{doi:10.7693/wl20131101}
\item
  Quan, H.-T., Dong, H. and Sun, C.-P. (2023). Theory and experiments of
  mesoscopic statistical thermodynamics. \emph{Acta Physica Sinica}
  72:230501 (in Chinese, with English abstract).
  https://doi.org/10.7498/aps.72.20231608
\item
  Parrondo, J. M. R. (2023). Thermodynamics of information.
  arXiv:2306.12447 (encyclopedia chapter preprint).
\item
  Strasberg, P. Thermodynamics and information processing at the
  nanoscale. Doctoral thesis, Technische Universität Berlin.
\item
  Dago, S. (2022). Thermodynamique stochastique: pilotage de
  micro-oscillateurs et applications à l'étude et l'optimisation du
  traitement de l'information. Doctoral thesis, École normale supérieure
  de Lyon (in French). https://theses.fr/2022LYSEN019
\item
  Lagoin, M. (2023). Thermodynamique dans les systèmes stationnaires
  hors équilibre macroscopiques: démons de Maxwell et machines
  thermiques aléatoires. Doctoral thesis, École normale supérieure de
  Lyon (in French).
\item
  Ciampini, M. A., Mancino, L., Orieux, A., Vigliar, C. et al.~(2017).
  Experimental extractable work-based multipartite separability
  criteria. \emph{npj Quantum Information} 3:10.
  https://doi.org/10.1038/s41534-017-0011-9
\item
  Sudakov, K. V. (2004). Functional systems theory and the probabilistic
  prediction of behavior. \emph{Neuroscience and Behavioral Physiology}
  34:505-507. https://doi.org/10.1023/B:NEAB.0000022638.96382.dd
\item
  Sudakov, K. V. (2015). Theory of functional systems: a keystone of
  integrative biology. In M. Nadin (ed.), \emph{Anticipation: Learning
  from the Past}, Springer, pp.~153-173.
  https://doi.org/10.1007/978-3-319-19446-2\_9
\item
  Sudakov, K. V. (2022). Cognitive activity from the perspective of
  functional systems theory. In \emph{Russian Cognitive Neuroscience},
  Brill, pp.~87-117. https://doi.org/10.1163/9789004505667\_005
\item
  Shvyrkov, V. B. (1985). Toward a psychophysiological theory of
  behavior. \emph{Advances in Psychology}, pp.~47-71.
  https://doi.org/10.1016/S0166-4115(08)61596-4
\item
  Anokhin, K. V. (2021). The cognitome: seeking the fundamental
  neuroscience of a theory of consciousness. \emph{Neuroscience and
  Behavioral Physiology} 51:915-937.
  https://doi.org/10.1007/s11055-021-01149-4
\item
  Saltykov, A. and Grachev, S. (2015). Anticipation and the concept of
  system-forming factor in the theory of functional systems. In M. Nadin
  (ed.), \emph{Anticipation: Learning from the Past}, Springer,
  pp.~507-520. https://doi.org/10.1007/978-3-319-19446-2\_30
\item
  Vityaev, E. E. (2015). Purposefulness as a principle of brain
  activity. In M. Nadin (ed.), \emph{Anticipation: Learning from the
  Past}, Springer, pp.~231-254.
  https://doi.org/10.1007/978-3-319-19446-2\_13
\item
  Vityaev, E. E. and Demin, A. V. (2011). Recursive subgoals discovery
  based on the functional systems theory. \emph{Frontiers in Artificial
  Intelligence and Applications} (BICA 2011), IOS Press.
  https://doi.org/10.3233/978-1-60750-959-2-425
\item
  Vityaev, E. E. and Demin, A. V. (2018). Cognitive architecture based
  on the functional systems theory. \emph{Procedia Computer Science}
  145:623-628. https://doi.org/10.1016/j.procs.2018.11.072
\item
  Vityaev, E., Kolonin, A., Kurpatov, A., Molchanov, A. et al.~(2022).
  Brain principles programming. arXiv:2202.12710 (preprint).
\item
  Potapov, A., Belikov, A., Bogdanov, V., Scherbatiy, A. et al.~(2019).
  Differentiable probabilistic logic networks. arXiv:1907.04592
  (preprint).
\item
  Paninski, L. (2003). Estimation of entropy and mutual information.
  \emph{Neural Computation} 15:1191-1253.
  https://doi.org/10.1162/089976603321780272
\item
  Russell, S. and Wefald, E. (1991). Principles of metareasoning.
  \emph{Artificial Intelligence} 49:361-395.
  https://doi.org/10.1016/0004-3702(91)90015-C
\item
  Horvitz, E. J. (1987). Reasoning about beliefs and actions under
  computational resource constraints. \emph{Proceedings of the Third
  Conference on Uncertainty in Artificial Intelligence (UAI 1987)},
  429-447. arXiv:1304.2759.
\item
  Simon, H. A. and Kadane, J. B. (1975). Optimal problem-solving search:
  all-or-none solutions. \emph{Artificial Intelligence} 6:235-247.
  https://doi.org/10.1016/0004-3702(75)90002-8
\item
  Sekimoto, K. (1998). Langevin equation and thermodynamics.
  \emph{Progress of Theoretical Physics Supplement} 130:17-27.
  https://doi.org/10.1143/PTPS.130.17
\item
  Wolpert, D. H., Korbel, J., Lynn, C. W., Tasnim, F., Grochow, J. A.,
  Kardeş, G. et al.~(2024). Is stochastic thermodynamics the key to
  understanding the energy costs of computation? \emph{PNAS}
  121:e2321112121. arXiv:2311.17166.
  https://doi.org/10.1073/pnas.2321112121
\item
  Chen, L., Zaharia, M. and Zou, J. (2023). FrugalGPT: how to use large
  language models while reducing cost and improving performance.
  arXiv:2305.05176 (preprint).
\item
  Ong, I., Almahairi, A., Wu, V., Chiang, W.-L. et al.~(2024). RouteLLM:
  learning to route LLMs with preference data. arXiv:2406.18665
  (preprint).
\item
  Koyuk, T. and Seifert, U. (2020). Thermodynamic uncertainty relation
  for time-dependent driving. \emph{Physical Review Letters} 125:260604.
  arXiv:2005.02312. https://arxiv.org/abs/2005.02312
\item
  Jordan, S. P. (2017). Fast quantum computation at arbitrarily low
  energy. \emph{Physical Review A} 95:032305. arXiv:1701.01175.
  https://arxiv.org/abs/1701.01175
\item
  Nemenman, I., Shafee, F. and Bialek, W. (2002). Entropy and inference,
  revisited. \emph{Advances in Neural Information Processing Systems}
  14. arXiv:physics/0108025. https://arxiv.org/abs/physics/0108025
\item
  Norton, J. D. (2005). Eaters of the lotus: Landauer's principle and
  the return of Maxwell's demon. \emph{Studies in History and Philosophy
  of Modern Physics} 36:375-411.
  https://doi.org/10.1016/j.shpsb.2004.12.002
\item
  Maroney, O. J. E. (2005). The (absence of a) relationship between
  thermodynamic and logical reversibility. \emph{Studies in History and
  Philosophy of Modern Physics} 36:355-374. arXiv:physics/0406137.
  https://arxiv.org/abs/physics/0406137
\item
  Sagawa, T. (2014). Thermodynamic and logical reversibilities
  revisited. \emph{Journal of Statistical Mechanics} P03025.
  arXiv:1311.1886. https://arxiv.org/abs/1311.1886
\item
  López-Suárez, M., Neri, I. and Gammaitoni, L. (2016). Sub-kBT
  micro-electromechanical irreversible logic gate. \emph{Nature
  Communications} 7:12068. https://doi.org/10.1038/ncomms12068
\item
  Paneru, G., Lee, D. Y., Park, J.-M., Park, J. T., Noh, J. D. and Pak,
  H. K. (2018). Optimal tuning of a Brownian information engine
  operating in a nonequilibrium steady state. \emph{Physical Review E}
  98:052119. https://doi.org/10.1103/PhysRevE.98.052119
\item
  Panizon, E. (2026). Optimal control of a mesoscopic information
  engine. arXiv:2603.29804 (preprint). https://arxiv.org/abs/2603.29804
\item
  Rolandi, A., Abiuso, P., Lipka-Bartosik, P. et al.~(2026).
  Energy-time-accuracy tradeoffs in thermodynamic computing.
  arXiv:2601.04358 (preprint). https://arxiv.org/abs/2601.04358
\item
  Burnell, R., Schellaert, W., Burden, J. et al.~(2023). Rethink
  reporting of evaluation results in AI. \emph{Science} 380:136-138.
  https://doi.org/10.1126/science.adf6369
\item
  Zhou, L., Pacchiardi, L., Martínez-Plumed, F. et al.~(2025). General
  scales unlock AI evaluation with explanatory and predictive power.
  arXiv:2503.06378 (preprint). https://arxiv.org/abs/2503.06378
\item
  Laber, M., Shafi, Z., Savcisens, G. et al.~(2026). Scaling of
  capability and efficiency at inference time in large reasoning models.
  arXiv:2609.27166 (preprint). https://arxiv.org/abs/2609.27166
\item
  Erol, M. H., El, B., Suzgun, M., Yuksekgonul, M. and Zou, J. (2025).
  Cost-of-Pass: an economic framework for evaluating language models.
  arXiv:2504.13359 (preprint). https://arxiv.org/abs/2504.13359
\item
  Saad-Falcon, J., Narayan, A., Akengin, H. O. et al.~(2025).
  Intelligence per Watt: measuring intelligence efficiency of local AI.
  arXiv:2511.07885 (preprint). https://arxiv.org/abs/2511.07885
\item
  Standard Performance Evaluation Corporation. SPECpower\_ssj2008.
  https://www.spec.org/power\_ssj2008/
\item
  Krakauer, D. C., Mitchell, M. and Krakauer, J. W. (2026). Large
  language models and emergence: a complex systems perspective.
  \emph{Philosophical Transactions of the Royal Society A} 384:20250014.
  https://doi.org/10.1098/rsta.2025.0014. arXiv:2506.11135.
\item
  Krakauer, D. C. (2026). Competitive and complementary tools.
  arXiv:2607.18460 (preprint). https://arxiv.org/abs/2607.18460
\item
  Kimura, M. (1961). Natural selection as the process of accumulating
  genetic information in adaptive evolution. \emph{Genetical Research}
  2:127-140. https://doi.org/10.1017/S0016672300000616
\item
  Worden, R. P. (1995). A speed limit for evolution. \emph{Journal of
  Theoretical Biology} 176:137-152.
  https://doi.org/10.1006/jtbi.1995.0183
\item
  Hledík, M., Barton, N. and Tkačik, G. (2022). Accumulation and
  maintenance of information in evolution. \emph{PNAS} 119:e2123152119.
  https://doi.org/10.1073/pnas.2123152119
\item
  Perrier, E. (2025). Watts-per-intelligence: Part I (energy
  efficiency). arXiv:2504.05328 (preprint).
  https://arxiv.org/abs/2504.05328
\item
  Xu, J. and Li, Z. (2025). Information physics of intelligence:
  unifying derivation depth and entropy under thermodynamic constraints.
  arXiv:2511.19156v4 (preprint; the arXiv listing title reads ``unifying
  logical depth and entropy''). https://arxiv.org/abs/2511.19156
\item
  Fagan, P. D. (2025). Toward a physical theory of intelligence.
  arXiv:2601.00021v2 (preprint; first submitted 22 December 2025,
  revised 7 March 2026). https://arxiv.org/abs/2601.00021
\item
  Kumar, S. (2026). The Laws of Intelligence: Computation Efficiency and
  the Physical Scaling Limits of AI Systems. Zenodo (preprint), concept
  DOI 10.5281/zenodo.18993403, version DOI 10.5281/zenodo.18993404.
  https://doi.org/10.5281/zenodo.18993404
\item
  Zhang, J., Zhu, H., Zhang, J. J. and Hu, H. (2026). The semantic
  least-energy principle: a hypothesis for intelligence.
  arXiv:2607.24287 (preprint). https://arxiv.org/abs/2607.24287
\end{enumerate}

\end{document}
